% De Oratore (de-oratore) --- English translation
% Translated by Claude (Anthropic) from the Latin (Perseus).
% Style guide: STYLE.md.

\ciceroBook{1}

\ciceroSection{1} When I look back, my brother Quintus, in long thought and recollection of old times, the men who seem to me supremely happy are those who, when the commonwealth was in its best estate, while flourishing both in office and in the fame of their public deeds, could hold a course of life such that they were able to be either at work without danger or at leisure with dignity. There was a time when I supposed that for me, too, the beginning of rest, and a turning back of the mind to the noble pursuits both of us share, would be just and almost universally conceded — once the unending labour of public business and the demands of ambition had drawn to a stop with the running through of offices, indeed with the bend of life itself.

\ciceroSection{2} That hope of my plans and intentions has been disappointed both by the gravity of the common times and by the various changes of my own fortune. For in the very place which seemed likely to be most full of quiet and tranquillity, the greatest masses of trouble and the most turbulent storms have arisen. Nor in any case has the fruit of leisure been given me, who longed and yearned for it, that I might cultivate among us, and revisit, the arts to which from boyhood we were devoted.

\ciceroSection{3} For our earliest years coincided with the upheaval of the old discipline; in our consulship we came down into the very midst of the contest and crisis of all things; and the whole time after our consulship we have flung against those waves which, driven back through us from the common ruin, have come surging upon ourselves alone. But still, in these harsh straits and narrow times, I shall yield to our pursuits, and shall give over to writing whatever leisure the malice of enemies, the cases of friends, or the commonwealth allows me. To you, brother, I shall fail neither when you urge nor when you ask: for no one with me has greater weight either by authority or by good will.

\ciceroSection{4} I must call up the recollection of an old memory — not by any means fully worked out, but, as I think, suited to what you are looking for: that you may learn what those men, the most eloquent and most brilliant of all, have thought about the whole theory of speaking.

\ciceroSection{5} You wish — as you have often said to me — that since the things which slipped from us as boys or growing youths out of our little notebooks, half-finished and rough, are scarcely worthy now either of this age of ours or of the experience we have gained from those many great cases I have just mentioned, something on the same subjects, more polished and more finished, should be brought forth by me. And in our discussions you sometimes disagree with me on this point: that I hold eloquence is bound up with the arts of the most learned men, while you think it should be set apart from elegance of learning and placed in a certain kind of natural talent and practice. Indeed, when I look upon the foremost men endowed with the foremost natural gifts, it has seemed to me again and again that I must ask why more men have proved admirable in everything than in speaking; for whithersoever you turn your mind and thought, you will see very many of the first rank in every kind of art — not modest arts, but almost the greatest.

\ciceroSection{6} For who, if he wished to measure the knowledge of distinguished men by the usefulness or the magnitude of their deeds, would not set the general above the orator?

\ciceroSection{7} And who would doubt that we can produce, from this one state, leaders in war of the very first rank almost without number, but barely a few who are excellent in speaking?

\ciceroSection{8} What is more, in counsel and wisdom, men capable of guiding and steering the commonwealth have arisen many in our own day, more in our fathers', more still in our ancestors' memory; while for a very long stretch there were no good orators at all, and barely one tolerable orator could be found in any one age. And lest anyone should think this rationale of speaking should rather be compared with other studies, which deal with hidden arts and a wide variety of letters, than with the praise of a general or the wisdom of a good senator, let him turn his mind to those very kinds of arts and look about: who has flourished in them, and how many. So he will most easily judge how few orators there are and have ever been.

\ciceroSection{9} Nor does it escape you that the procreatrix and as it were the parent of all praised arts is judged by the most learned men to be that thing the Greeks call \textit{philosophia}; in which it is hard to count up how many men have been outstanding in knowledge and in the variety and abundance of their pursuits — men who have not laboured separately in some one subject, but have grasped everything they could either by the investigation of knowledge or by the method of reasoning.

\ciceroSection{10} Who is unaware in what darkness of subject and in what hidden, manifold and subtle art those who are called mathematicians move? And yet in that field so many men have proved themselves perfect that hardly anyone seems to have applied himself with any vigour to that science without attaining what he set out to. Who has given himself wholly to musicians, who to that study of letters which the men called grammatici profess, without grasping by knowledge and understanding the well-nigh limitless extent and content of those arts?

\ciceroSection{11} I think I shall be saying this truly: that of all those who have given themselves to the most liberal pursuits and learning of these arts, the smallest abundance has been of outstanding poets and orators. And in this very class — in which one excellent man arises only very rarely — if you should compare diligently both from our own and from the Greek supply, you will still find far fewer good orators than poets.

\ciceroSection{12} Which ought to seem the more remarkable since the studies of the other arts are drawn for the most part from hidden, recondite springs; whereas the whole rationale of speaking is set out in the open, in a sort of common usage, in the mouth and conversation of men. So in the others, what most excels is what is furthest removed from the comprehension and feeling of the unlearned; but in speaking, the greatest fault of all is to recoil from the everyday kind of language and from the habits of common feeling.

\ciceroSection{13} Nor in fact can it truly be said that more men devote themselves to other arts, or are stirred by greater pleasure or richer hope or larger rewards to mastery of them. To pass over Greece — which has always wished to be first in eloquence, and to leave Athens, that mother-inventor of all branches of learning, in which the highest force of speaking was both invented and brought to perfection — even in this state of ours, surely no studies have flourished more vehemently than the studies of eloquence.

\ciceroSection{14} For after the empire of all peoples was established and the long stretch of peace had confirmed leisure, hardly any youth eager for fame did not think it incumbent on himself to bend every effort of zeal towards speaking. And at first, indeed, ignorant of the whole science, supposing that there was no force of practice nor any precept of art, our men got as far as they could by their own native ability and thought. Afterwards, having heard the Greek orators and become acquainted with their writings and called in their teachers, our men burned with an incredible passion for learning.

\ciceroSection{15} They were stirred by the magnitude, the variety, and the multitude of cases of every kind, so that to the learning each had got from his own study was added constant practice, which surpasses the precepts of all teachers. And the rewards set out for this study were, as they still are, the very greatest: for popular favour, or for power, or for high station. The native talents of our men too, as we may judge from many things, have far surpassed those of all other peoples.

\ciceroSection{16} For these reasons, who would not justly wonder that out of the whole record of ages, of times, of states, so small a number of orators is to be found? But the matter is something greater than men suppose, and is gathered out of more arts and pursuits. For what other cause should anyone suppose, in the great multitude of learners, the abundant supply of teachers, the most distinguished native talents, the limitless variety of cases, the most ample prizes set before eloquence — what other cause, but a certain incredible greatness and difficulty in the thing itself?

\ciceroSection{17} For there must be acquired a knowledge of very many subjects (without which a fluency of words is empty and laughable); the speech itself must be shaped, not only by the choice but by the arrangement of words; and all the movements of feeling which the nature of things has assigned to humankind must be thoroughly known, since the whole force and method of speaking has to be brought out in calming or in stirring the minds of those who listen. To these must be added a certain charm and wit, a learning worthy of a free-born man, and a quickness and brevity, both in answering and in attacking, joined with subtle elegance and urbanity. All antiquity must besides be held fast, with the force of historical examples; and the knowledge of laws and the civil law must not be neglected.

\ciceroSection{18} For why should I say more about delivery itself, which has to be regulated by the motion of the body, by gesture, by countenance, by the shaping and variety of the voice? How much it amounts to all by itself, the light art of actors and the stage make plain: in this art all labour to control face, voice and motion, but who is unaware how few there are, and have been, whom we can watch with composure? What shall I say of memory, the storehouse of all things? If memory is not brought in as the keeper of the matter and the words discovered and thought out, we know that everything, however brilliant in the orator, is doomed to perish.

\ciceroSection{19} Therefore let us cease to wonder what the cause is of the scarcity of eloquent men, since eloquence is composed of all those things, in any single one of which to do good work is a very great matter. Rather let us urge our children and the others whose glory and dignity is dear to us, that they grasp the magnitude of the thing in their minds, and that they trust to attain what they aspire to — not by those precepts or teachers or exercises which all use, but by certain other means.

\ciceroSection{20} And in my judgment no one will be able to be the orator complete in every grace, unless he has gained the knowledge of all the great matters and arts. For from the knowledge of things speech ought to flower and overflow; and unless the matter is grasped and known by the speaker, the speech itself has a kind of empty and well-nigh childish utterance.

\ciceroSection{21} Yet I shall not lay so great a burden upon our orators, especially in this great preoccupation of city and life, as to think nothing may be lawfully unknown to them — although the very force of the orator and the profession of speaking well does seem to undertake and promise that on whatever subject is set before him, he should speak with ornament and abundance.

\ciceroSection{22} But because I do not doubt that this seems to most an immense and infinite undertaking, and because I see that the Greeks, abounding not only in talent and learning but also in leisure and zeal, have already made a kind of partition of the arts, and that no one of them has laboured in the whole sphere, but they have set apart from the other branches of speaking that part which deals in the forensic disputes of trials or deliberations, and have left this one kind to the orator: I shall therefore include in these books no more than what nearly by the consensus of the greatest men, after the matter has been debated and much discussed, has been allotted to this kind of speaking.

\ciceroSection{23} And I shall trace, not from the cradle of our old, schoolboy learning, a certain sequence of precepts, but those things which I once received as having been treated in the discussion of the most eloquent men among us and the foremost in every dignity. Not because I despise what the Greek artists and teachers of speaking have left behind — but because, since those things lie open and ready for everyone, and cannot be either more handsomely set out or more plainly expressed by my interpretation, you will, I think, grant me this, my brother: that I prefer the authority of the men to whom the highest praise of speaking has been conceded by our own people, over the Greeks.

\ciceroSection{24} Now when the consul Philippus was inveighing more and more vehemently against the cause of the leading men, and the tribunate of Drusus, undertaken in defence of the senate's authority, seemed already to be cracking and weakening, I remember being told that, in the days of the Roman Games, Lucius Crassus had withdrawn to his Tusculan villa as if to gather himself; and that there had also come, his former father-in-law Quintus Mucius (Scaevola), and Marcus Antonius, the man who was both a partner in the political plans and joined to Crassus in the closest intimacy.

\ciceroSection{25} They had also gone out with Crassus himself, the young men, especially friends of Drusus, in whom the older men were then placing the great hope of their own dignity: Gaius Cotta, who was at that time canvassing for the tribunate of the people, and Publius Sulpicius, who was thought likely to seek the same office in turn.

\ciceroSection{26} On the first day, until the very end of it, these men talked among themselves about the times and about the whole commonwealth — the very matter for which they had come. In that conversation Cotta used to say there were many things lamented and recalled with prophetic foresight by those three men of consular rank: nothing of evil afterwards befell the state which they had not seen impending so long before.

\ciceroSection{27} When that whole conversation was done, there was so much courteous spirit in Crassus that, when they had bathed and reclined, all the gloom of the earlier talk was dispelled, and there was such cheerfulness in the man and so much charm in his speaking that the day among them seemed to have been one of the senate-house, the dinner one of a Tusculan villa.

\ciceroSection{28} On the next day, when those older men had rested enough and they had gone out for a walk, Cotta would relate that Scaevola, after two or three turns, said: ``Why don't we imitate, Crassus, that Socrates of yours in Plato's \textit{Phaedrus}? For this plane-tree of yours has put me in mind of him, spread out with its broad branches no less to shade this place than that one was, whose shade Socrates followed — a tree which seems to me to have grown not so much from that little stream which is described, as from Plato's prose. And what he, with the toughest feet, did when he flung himself onto the grass and so spoke those words which the philosophers report as having been said with divine inspiration — that, surely, is fairer to be granted to my own feet.''

\ciceroSection{29} Then Crassus: ``On the contrary, even more comfortably than that''; and he called for cushions and they all sat on those couches that were under the plane-tree. There, Cotta used to say, that all their minds might unbend from the earlier talk, Crassus brought up a discussion about the study of speaking.

\ciceroSection{30} He began by saying that Sulpicius and Cotta had no need of his exhortation; rather, both should be praised, since they had already attained such ability that they were not only set ahead of their contemporaries but compared with their elders. ``And nothing,'' he said, ``seems to me more excellent than to be able by speaking to hold the minds of human assemblies, to draw their wills along, to drive them where one wishes, and from where one wishes to draw them away. This one thing, in every free people, and most of all in peaceful and quiet states, has always pre-eminently flourished and always held sway.

\ciceroSection{31} For what is so wonderful as that, out of the limitless multitude of men, there should arise one — alone, or with very few — who is able to do that very thing which nature has given to all? Or so pleasant to know and to hear, as a speech adorned and polished with wise opinions and weighty words? Or so powerful and so majestic, that the movements of a people, the consciences of judges, the gravity of the senate are turned about by the speech of one man? What, further, is so kingly, so generous, so munificent, as to bring help to the suppliant, to raise up the broken, to give safety, to free from danger, to keep men within the state?

\ciceroSection{32} What, again, is so necessary as always to have weapons in hand, with which one may either keep oneself sheltered, or challenge the safe, or, if provoked, take revenge? Come now — that you may not always think of the Forum, the benches, the rostra and the curia — what can there be in leisure either more pleasant or more characteristic of human civility, than wit and conversation rough in nothing? For in this one thing especially we surpass even the beasts: that we converse among ourselves, and that we can express by speaking what we feel.

\ciceroSection{33} Therefore who would not justly wonder at this, and think it should be most laboured at, that in the very thing wherein men most excel beasts, a man should excel his fellow-men? But to come now to the highest things: what other power could either gather scattered men into one place, or lead them away from a wild and rustic life into this human and civil cultivation, or, when states were already established, set out laws, courts, and rights?

\ciceroSection{34} And not to chase too many things, which are well-nigh innumerable, I shall sum up briefly. For I lay it down: in the moderation and wisdom of the perfect orator is contained, in the highest measure, not only his own dignity, but the safety also of very many private men and of the whole commonwealth. Therefore go on as you are doing, young men, and lay yourselves to that pursuit in which you are engaged, that you may be at once an honour to yourselves, a usefulness to your friends, and a benefit to the commonwealth.''

\ciceroSection{35} Then Scaevola, kindly as was his way: ``In other things,'' he said, ``I agree with Crassus — let me not detract from the art or the glory of either Gaius Laelius, my father-in-law, or of his son-in-law here. But on those two points, Crassus, I am afraid I cannot grant your case: first, that you said states have been founded in the beginning, and often preserved, by orators; second, that, leaving aside the Forum, the assembly, the courts, the senate, you have decreed that the orator is perfect in every kind of speech and of culture.

\ciceroSection{36} For who would grant you this — that in the beginning the human race, scattered in mountains and forests, hedged itself within towns and walls because it was driven thither by the counsels of the wise, rather than because it was charmed thither by the speech of the eloquent? Or that the rest of the benefits, either in founding or in preserving states, were established not by wise and brave men but by men speaking eloquently and ornately?

\ciceroSection{37} Or do you think that famous Romulus gathered the shepherds and the migrants, or joined the marriage-bonds with the Sabines, or held back the violence of his neighbours by eloquence — and not by extraordinary counsel and wisdom? In Numa Pompilius, in Servius Tullius, in the other kings — many of whose deeds were excellent for the founding of the commonwealth — does any trace of eloquence appear? When the kings were driven out — though the very driving-out we see was carried through by the mind, not the tongue, of L. Brutus — do we not see all that follows full of counsels, empty of words?

\ciceroSection{38} If I wished to use both the examples of our own state and of others, I could bring forward more harms than helps to public affairs imported by the most eloquent men. To pass over the rest: of all the men I have heard — except, Crassus, you two — I judge the most eloquent to have been Tiberius and Gaius Sempronius (Gracchus). Their father, a prudent and weighty man but by no means eloquent, was again and again, and most of all as censor, the salvation of the commonwealth: not by some careful abundance of speech but by a nod and a word he transferred the freedmen into the urban tribes — and if he had not done so, the commonwealth which we now scarcely keep, we should long ago have ceased to have at all. But his sons, eloquent men and equipped with all the resources both of nature and of learning for speaking, when they had received a state most flourishing both by their father's counsel and by ancestral arms, scattered the commonwealth — by that splendid governess of states, as you call it: eloquence.

\ciceroSection{39} What of the old laws and the customs of our ancestors? What of the auspices, over which you and I, Crassus, preside with great safety to the commonwealth? What of the religions and ceremonies? What of these civil laws, which have for so long been studied in our family without any praise of eloquence — were they discovered, or known, or in any way handled by the order of orators?

\ciceroSection{40} For my part I remember Servius Galba, a divine man in speaking, and Marcus Aemilius Porcina, and Gaius Carbo himself, whom you, when only a young man, struck down: men ignorant of the laws, hesitant in the institutions of our ancestors, untaught in the civil law. And this generation of yours — except for you, Crassus, who learned the civil law from us, more by your own zeal than by any function proper to the eloquent — is so ignorant of the law that one is sometimes ashamed of it.

\ciceroSection{41} As for what you took up at the end of your speech, as if of right — that the orator can move with great fluency in the discussion of every kind of subject — that, if we were not here in your own kingdom, I should not have endured. I should have set many at the head of the cause who would either contend with you by the interdict or summon you, by law, to lay hand on the property: for having broken so rashly into other men's possessions.

\ciceroSection{42} For first the Pythagoreans and Democriteans would prosecute you under the law, and the rest of the natural philosophers would claim what is theirs by right — men adorned and weighty in speaking, with whom you would not be allowed to enter the lawful wager of contestation. Then the herds of the philosophers would press in: from that fount and head, Socrates himself, they would convict you of having learned nothing about the goods of life, nothing about the evils, nothing about the movements of the soul, nothing about human character, nothing about the science of life — of having looked into nothing, of knowing nothing. And when the whole army had charged you in mass, then the individual schools would file their suits one by one. The Academy would press, who would force you to deny your own every word.

\ciceroSection{43} Our Stoics would hold you tangled in the snares of their disputations and questionings. The Peripatetics would prove that even those very things which you take to be the proper helps and ornaments of orators must be sought from them; they would show that Aristotle and Theophrastus have written about those subjects not only better but far more abundantly than all teachers of speaking.

\ciceroSection{44} I leave aside the mathematicians, the grammarians, the musicians, with whose arts that speaking-power of yours is not joined by even the slightest connection. Therefore, Crassus, I would not have you profess these so many and so great matters. It is enough — and it is much — that you can do this: that in the courts whatever cause you plead seems the better and the more probable, that in assemblies and in the giving of opinions your speech avails the most for persuasion, and finally that you seem to speak eloquently in the eyes of the prudent, and even truly in the eyes of the foolish. If you can do anything beyond this, that, it will seem to me, will be no orator, but Crassus, by some power proper to himself, not common to the orators.''

\ciceroSection{45} Then he: ``I am not unaware, Scaevola,'' said Crassus, ``that those things are debated and disputed among the Greeks. For when, as quaestor, I came to Athens from Macedonia, I heard the foremost men there, the Academy in flower (so it was held in those times), at a time when Charmadas, Clitomachus and Aeschines held it. There was Metrodorus too, who together with them had heard with most diligence that Carneades himself, a man, as report had it, of all in speaking the keenest and most copious. And Mnesarchus the hearer of your Panaetius was flourishing then, and Diodorus the Peripatetic, hearer of Critolaus.

\ciceroSection{46} There were besides many illustrious and noble men in philosophy, by all of whom in well-nigh one voice I saw the orator being driven away from the helms of states, shut out from all learning and the knowledge of greater matters, and thrust together into courts and small assemblies as into a kind of mill.

\ciceroSection{47} But I neither agreed with them, nor with the inventor and prince of these debates, by far the most weighty and eloquent of all in speaking — Plato, whose \textit{Gorgias} I read with most diligence at Athens then with Charmadas. In that book what I most admired in Plato was that, in deriding orators, he himself seemed to me the highest orator. The dispute is over a word, and has long tortured Greeklings more eager for contention than for truth.

\ciceroSection{48} For if anyone defines the orator as one who can only speak fluently before the law, in the courts, before the people, or in the senate, even to that man he must concede and grant much. For without much going-over of all public matters, and without knowledge of laws, customs, and right, and without an understanding of the nature and characters of men, he cannot move with sufficient skill and craft even in those very matters. And the man who has learned these — without which no one can defend rightly even the smallest things in cases — what can be lacking to him from the knowledge of the greatest things? But if you wish nothing of the orator except to speak in well-arranged, ornate, and copious style, I ask: how can he attain that very thing without that knowledge which you do not concede to him? For the virtue of speaking cannot exist unless the things he is to say have been grasped by the speaker.

\ciceroSection{49} Therefore, if that natural philosopher Democritus spoke ornately — as both is reported and seems to me — that subject-matter was the natural philosopher's, on which he spoke, but the ornament of the words must be reckoned the orator's. And if Plato spoke divinely on subjects most remote from civil disputes — which I grant — and if Aristotle, Theophrastus, Carneades likewise were eloquent in the things on which they disputed and pleasing and ornate in their speaking, even though those subjects on which they argue lie in certain other studies, the speech itself is proper to this one art of which we speak and inquire.

\ciceroSection{50} For we see that on the same subjects some have argued meagrely and dryly — that Chrysippus, the man they call the keenest, for instance — and have not on that account fallen short of philosophy because they did not have this faculty of speaking from another, alien art. What then is the difference, or how will you tell apart the abundance and copiousness of those I have named in speaking, from the dryness of those who do not use this variety and elegance of speech? There will surely be one thing which those who speak well bring as their own: an arranged and ornate speech, marked by a kind of art and polish. But this speech, if there is no matter beneath it grasped and known by the orator, must be either nothing, or laughed at by everyone.

\ciceroSection{51} For what is so mad as the empty resounding of words, even of the best and most ornate, with no thought or knowledge under them? Whatever it shall be, then, from whatever art, of whatever kind, the orator — if he has learned it, as it were the case of a client — will speak it better and more handsomely than the very inventor and craftsman of the matter himself.

\ciceroSection{52} For if there shall be anyone to say that there are certain thoughts and cases proper to orators, and a knowledge of certain things bounded within the lattices of the courts, I shall confess indeed that this our way of speaking moves more constantly in those subjects. But even within these subjects there are very many things which the masters themselves, who are called rhetoricians, neither pass on nor possess.

\ciceroSection{53} For who does not know that the greatest force of the orator stands forth in the minds of men, either in stirring them to anger or hatred or grief, or in calling them back from those very emotions to mildness and pity? But unless one has thoroughly seen into the natures of men and the whole force of human feeling and the causes by which men's minds are stirred or turned aside, he will not be able to accomplish by speaking what he wishes. And this whole region seems proper to the philosophers, and the orator on my counsel will never dispute it;

\ciceroSection{54} but when he has yielded to them the knowledge of subjects, since they wished to labour in that alone, the handling of speech (which without that knowledge is nothing) he will take to himself. For this is the orator's own — as I have said often by now — speech weighty and ornate and adapted to the feelings and minds of men.

\ciceroSection{55} I confess that Aristotle and Theophrastus have written on these subjects; but, Scaevola, take care that this is not all on my side. For I do not borrow from them what is shared between the orator and them; rather, those men concede that what they argue about these subjects belongs to orators, and so they entitle and call their other books by the names of their own art, but these they entitle and call \textit{rhetorical}.

\ciceroSection{56} For when those topics fall into a man's speech, as very often happens — that he must speak about the immortal gods, about piety, about concord, about friendship, about the common right of citizens, of men, of nations, about equity, about temperance, about greatness of soul, about every kind of virtue — all the gymnasiums and all the schools of the philosophers will cry out, I suppose, that all these things are properly theirs and have nothing whatever to do with the orator.

\ciceroSection{57} Even when I have conceded to them that they may argue these matters in their corners, to spend their leisure as they choose, I shall still grant and assign to the orator this: that he should set out the very same things on which they argue with a thin and bloodless speech, with all pleasantness and weight. These things I argued at the time at Athens with the very philosophers, for our Marcus Marcellus pressed me to it — the man who is now curule aedile, and who, certainly, were he not now putting on the games, would be present at this conversation of ours; even then, when only a young man, he was wonderfully devoted to these studies.

\ciceroSection{58} And as for the founding of laws, war, peace, allies, taxes, the right of citizens distributed in classes by orders and ages — let the Greeks, if they wish, declare that Lycurgus or Solon (though we ourselves do reckon those men among the eloquent) knew these things better than Hyperides or Demosthenes, men by then perfect and wholly polished in speaking; or let our own people put the Decemvirs who set down the Twelve Tables — men who must have been wise — ahead, in this kind of matter, of Servius Galba and your father-in-law Gaius Laelius, men who confessedly excelled in the glory of speaking.

\ciceroSection{59} For I shall never deny that there are certain departments proper to those who have laid all their study on the discovery and treatment of those subjects. But I say the full and perfect orator is the one who can speak on every subject copiously and variously. For often even in those cases which all confess to be proper to orators, there is something which must be drawn out and taken not from the forensic practice — which alone you concede to the orator — but from some more obscure science.

\ciceroSection{60} For I ask whether one can speak either against or for a general without practical experience of warfare, or even, in many cases, without a knowledge of regions of land and sea? Whether one can speak before the people about laws to be passed or vetoed, or in the senate about every kind of public matter, without the highest knowledge and prudence of civic affairs? Whether speech can be brought to bear on the senses and movements of men's minds — to set them on fire or to extinguish them, the one thing which dominates in the orator — without the most diligent investigation of all those reasonings concerning the natures and morals of the human race which the philosophers expound?

\ciceroSection{61} I do not know but that I shall convince you of this less easily; but I shall not hesitate to say what I feel. Those very subjects you set down a little while ago as proper to the other arts — natural philosophy, mathematics — are the knowledge of those who profess them; but if anyone wishes to set those very arts off in light by speech, he must take refuge in the orator's faculty.

\ciceroSection{62} For if it is agreed that Philo, the architect who built the arsenal for the Athenians, gave the people a most eloquent account of his work, that should not be reckoned the eloquence of an architect rather than of an orator; nor, had it fallen to our Marcus Antonius here to speak on Hermodorus's behalf about a navy-yard's construction, would he have failed, after he had learned the case from him, to speak ornately and copiously about another man's craft. Nor again did Asclepiades, the doctor and friend whose services we used, surpass the rest of the doctors in eloquence in those very matters in which he spoke ornately by using the faculty of medicine, but of eloquence.

\ciceroSection{63} And what Socrates was wont to say is more probable, though not in fact true: that all are sufficiently eloquent in what they know. The truer thing is this: no one can be eloquent on what he does not know; and if he knows it best but is ignorant of how to make and polish a speech, he cannot speak even that very subject he knows about with eloquence.

\ciceroSection{64} Therefore, if anyone wishes to define and embrace the universal and proper force of the orator, he, in my judgment, will be the orator worthy of so weighty a name who, on whatever subject shall fall to be set out by speech, will speak with prudence, in good order, with ornament and with memory — together with a certain dignity also of delivery.

\ciceroSection{65} If anyone thinks that I have set the field too unbounded with my ``on whatever subject,'' let him pare and trim it however much he likes; yet I shall maintain this — that even if the orator is ignorant of what is set in the other arts and pursuits and grasps only what is in disputes and forensic practice, even on those very subjects, when he must speak, after he has learned from those who hold them what each thing's matter is, the orator will speak much better than those very men whose arts they are.

\ciceroSection{66} So if our Sulpicius here must speak on military matters, he will inquire of our kinsman Gaius Marius; and when he has been told, he will pronounce in such a way that to Gaius Marius himself the man here will seem almost to know those things better than Marius does. If on the civil law, he will consult you, and you, the most prudent and experienced man in those very things which he will have learned from you, he will surpass in the art of speaking.

\ciceroSection{67} If some matter shall fall in which he must speak about nature, or about the vices of men, about desires, about moderation, about self-control, about pain, about death — perhaps, if it has seemed best to him (although these things the orator ought to know) — he will consult Sextus Pompey, a learned man in philosophy. He will surely accomplish this: that whatever subject he has learned from anyone, he will speak on it far more handsomely than the very man from whom he has learned.

\ciceroSection{68} But if he will hear me, since philosophy is divided into three parts — into the obscurity of nature, the subtlety of disputation, life and morals — let us leave the first two and yield them to our laziness. The third, which has always been the orator's, unless we hold it, we shall leave nothing to the orator in which he can be great.

\ciceroSection{69} Therefore this whole region — life and morals — must be thoroughly learned by the orator. The rest, even if he has not learned them, he will still be able to set out by speech, when need shall arise, provided they have been brought and handed over to him. For if it is agreed among the learned that Aratus, a man ignorant of astronomy, spoke about the heavens and the stars in the most ornate and finished verses; and if Nicander of Colophon, a man as far from the country as could be, wrote brilliantly on rural matters by a poetic, not a rustic, faculty — what reason is there that the orator should not speak with the highest eloquence on the subjects he has learned for a particular cause and occasion?

\ciceroSection{70} For the poet is the orator's neighbour: a little more bound by metre, but freer in the licence of his words; a partner and almost equal in many kinds of ornament; and in this surely much the same — that he confines and defines his right by no boundaries, but is free to range with the same faculty and abundance wherever he wishes.

\ciceroSection{71} Now as to what you said, Scaevola, you would not have endured if you were not in my own kingdom — that I had said the perfect orator must be perfect in every kind of speech, in every department of culture: by Hercules, I should never say this, if I supposed that the man I am picturing were myself.

\ciceroSection{72} But, as Gaius Lucilius used often to say — a man somewhat angry at you, less close to me than he wished for that very reason, but learned and most urbane — so I feel: that no one should be reckoned in the number of orators who is not polished in all those arts which are worthy of a free man. And even if we do not use these arts in our speech, it still appears and stands forth whether we are unschooled in them or have learned them. So those who play ball do not use in the play itself the artistry proper to the wrestling-school, but the very motion shows whether they have learned the wrestling-school's training or not; and those who model figures, even though they then make no use of painting, it is still not obscure whether they know how to paint or not.

\ciceroSection{73} Just so, in these very speeches before the courts, the assemblies, the senate, even when the rest of the arts are not properly brought in, it is still easily made clear whether the speaker has only been knocked about in this declaiming work, or has come to speaking equipped with all the liberal arts.''

\ciceroSection{74} Then, laughing, Scaevola: ``I shall not contend with you any longer, Crassus. For that very thing you spoke against me you have got by a kind of artistry: that you both conceded to me what I myself wished should not belong to the orator, and yet, somehow or other, you twisted those very things back, and handed them over as proper to the orator.

\ciceroSection{75} For when I came to Rhodes as praetor, and brought to that great teacher of that discipline, Apollonius, the things I had received from Panaetius, he, as was his way, laughed at philosophy and despised it, and said many things not so much weightily as wittily. But your speech was of such a kind, not as to despise any art or learning, but to say that all are companions and handmaidens of the orator.

\ciceroSection{76} If anyone embraced all of these in himself, and added besides to them this faculty of the most ornate speech, I cannot say that he would not be some splendid and admirable man. But that man, if he existed, or if he ever had existed, or, indeed, if he could exist, you alone would surely be — you who, in my judgment and in everyone's, have left scarcely any praise to the rest of the orators (with no offence to those here).

\ciceroSection{77} But if you yourself lack nothing of what concerns forensic and civil matters, that you should know it, and yet you have not embraced that knowledge which you are joining to the orator, let us see that we do not assign him more than the matter and truth itself concedes.''

\ciceroSection{78} Here Crassus: ``Remember,'' he said, ``that I spoke not of my own faculty but of the orator's. For what have we either learned or been able to know — we who have come to action before we have come to learning, whom the matter itself in the Forum, in canvassing, in public affairs, in our friends' business has finished off before we could have any inkling of so great matters?

\ciceroSection{79} If, then, you think there is so much in us — in whom, even if natural ability, as you suppose, was not utterly lacking, learning surely was lacking, and leisure, and, by Hercules, that keenest zeal of learning — what do you suppose, if those things which I have not touched should be added to the talent of someone perhaps even greater than ours? What sort of orator and how great will he be?''

\ciceroSection{80} Then Antonius: ``You convince me, Crassus, of those things you say; nor do I doubt that anyone will be far more richly equipped in speaking who has

\ciceroSection{81} grasped the principle and nature of all subjects and arts. But first, this is a hard thing to do, especially in this life of ours and in our occupations. Then there is also this to be feared, lest we be drawn off from this practice and habit of speaking among the people and in the Forum. For another kind of speech seems to me to belong to those men of whom you spoke a little while ago, however ornately and weightily they speak either of nature or of human matters: a certain glossy and rich kind of words, of the wrestling-school and the oil rather than of this civic press and the Forum.

\ciceroSection{82} For I myself, who had laid hand on Greek learning late and lightly, when, on my way as proconsul to Cilicia, I came to Athens, lingered there several days because of the difficulty of sailing. And as every day I had with me men of the highest learning — almost the same men just named by you — when somehow it had got around among them that I, like you, was wont to take part in the greater cases, each one of them, as he was able, argued about the duty and the principle of the orator.

\ciceroSection{83} Some of these — like that very Mnesarchus — said that those whom we call orators were nothing but a kind of workmen with a swift and exercised tongue; that no one was an orator unless he was wise, and that eloquence itself, since it consisted of the science of speaking well, was a single virtue; and that he who had one virtue had all, and these were equal and on a level among themselves; so that he who was eloquent had all the virtues and was wise. But this was a thorny and meagre kind of speech, and far removed from our common sense.

\ciceroSection{84} Charmadas, however, used to speak much more copiously on the same subjects — not that he opened his own opinion (this was the ancestral custom of the Academy: to argue against everyone in disputation); but most clearly he used to signify this — that those who were called rhetoricians and who handed down precepts of speaking, plainly held nothing, and that no one could attain the faculty of speaking who had not learned the discoveries of the philosophers.

\ciceroSection{85} Eloquent Athenian men, versed in public life and in cases, argued against him; among whom was the man who recently was at Rome — Menedemus, my own guest. When he said there was a certain prudence which was concerned with seeing through the principles of founding and governing states, the other man would be roused — a man ready, abounding in every kind of learning, with a kind of incredible variety and abundance of subjects — and he would teach that all the parts of that very prudence had to be sought from philosophy; and that those things which were laid down in a state about the immortal gods, about the discipline of the youth, about justice, patience, temperance, the limit of all things, and the rest, without which states either cannot exist or be well ordered, were nowhere to be found in those rhetoricians' books.

\ciceroSection{86} If those teachers of rhetoric embraced so great a force of the greatest matters in their art, he would ask, why were their books filled with prefaces and epilogues and trifles of that kind (so he called them), while not a letter on the founding of states, on the writing of laws, on equity, justice, fidelity, the breaking of desires, the moulding of human characters, was to be found in their books?

\ciceroSection{87} The very precepts he was wont to mock so as to show that those men were not only unskilled in the prudence which they claimed for themselves, but knew not even this very rationale and method of speaking. For he held the orator's first concern was that he himself should appear to those before whom he was pleading such as he wished himself to seem; and this was done by dignity of life, of which those rhetoricians had left nothing in their precepts. And his second concern was that those who heard should be so affected in mind as the orator wished them to be affected; which likewise could in no way be done unless the speaker had learned in how many ways, and by what means, and by what kind of speech the minds of men are moved in each direction. But these things were buried and hidden deep in the middle of philosophy, which those rhetoricians had not touched even with the tip of their lips.

\ciceroSection{88} These things Menedemus tried to refute more by examples than by arguments. For from memory he would deliver many splendidly written passages from Demosthenes's speeches, and would teach that, in stirring the minds of judges or of the people in any direction by speaking, that man had not been ignorant of the means by which he was attaining what Charmadas denied anyone could know without philosophy.

\ciceroSection{89} To this Charmadas would reply that he did not deny Demosthenes had had the highest prudence and the highest force of speech, but, whether he had been able to do it by his own genius, or — what was agreed — had been an eager hearer of Plato, the question must be not what he had been able to do, but what those rhetoricians taught.

\ciceroSection{90} Often too he was carried in his speech to such a point that he argued that there was no art of speaking at all. He would teach this with arguments — that we were so born that we could both flatter subtly those from whom we had to seek something, and frighten our adversaries with menacing tones, and set out the matter at hand, and confirm what we intended, and refute what was said on the other side, and at the last beg something off and complain — those very things in which all the orator's faculty moves; and that habit and exercise sharpened the prudence of understanding and quickened the swiftness of expression.

\ciceroSection{91} He also relied on a wealth of examples. For first, as if of set purpose, he would say that no writer of an art on speaking had been even moderately eloquent, going back to one Corax and Tisias, who, it was agreed, had been the inventors and chiefs of that art; while of men most eloquent who had neither learned those things nor cared at all to know them, he would name countless ones. Among whom too — whether deriding me, or because he so thought and had so heard — he would put me forward in that number, who had not learned those things and yet, as he himself said, could do something in speaking. Of which I readily granted him the one — that I had learned nothing — but in the other I judged either that he was making fun of me, or that he himself was mistaken.

\ciceroSection{92} He denied indeed that there was any art at all save what was contained in things known and seen through to the bottom, looking to one outcome, and never deceiving. But all those things which were treated by orators were doubtful and uncertain, since they were said by men who did not plainly hold them all, and were heard by those to whom no science was being handed over but the false or at any rate obscure opinion of a brief moment.

\ciceroSection{93} What more is to be said? So he then seemed to me to argue that there was no art at all of speaking, and that no one could speak either skilfully or copiously without knowing those things which were said by the most learned men in philosophy. In which Charmadas was wont vehemently to admire your talent, Crassus: for, as he said, I had appeared to him most easygoing in listening, you most pugnacious in dispute.

\ciceroSection{94} And so I, drawn on by this same opinion, even wrote in a certain little book — which slipped from me unawares and against my will and came into the hands of men — that I had known some skilful speakers, but as yet no eloquent one; for I held the skilful speaker to be one who could speak before middling men out of a sort of common opinion of mankind clearly enough and acutely enough, but the eloquent one to be he who could amplify and adorn what he wished more wonderfully and magnificently, and who held in his mind and memory all the springs of all subjects pertaining to speaking. If this is hard for us, who, before we have begun to learn, are buried in canvassing and the Forum, let it still be set in the matter and the nature of the thing.

\ciceroSection{95} For I — as far as I divine by conjecture, and seeing how great talents there are in our men — do not despair that there will be someone someday who, with keener zeal than we are or have been, with greater leisure and faculty for learning and at a riper age, with greater labour and industry, when he has given himself over to hearing and reading and writing, will stand forth such an orator as we are seeking — one who can rightly be called not only skilful but eloquent. Yet in my judgment that is either Crassus here already, or, if there shall be any of equal talent who has heard and read and written more than he, the man will be able to add a little something to him.''

\ciceroSection{96} At this point Sulpicius: ``Beyond hope of mine and of Cotta, but greatly wished by both of us, has it fallen out, Crassus, that you have slipped into this kind of conversation. For as we were coming up here, it seemed pleasant enough if, while you were speaking of other things, we might still pick up something out of your conversation worthy of memory. But that you should come right into the very heart, all but the inmost part, of the discussion of this whole pursuit, art, faculty — that hardly seemed to us a thing to be hoped for.

\ciceroSection{97} For I, who from earliest youth had been fired with zeal for both of you, and with love also for Crassus — when I was constantly with him I could never draw a word from him about the force and method of speaking, though I had pressed him myself and often tried through Drusus. In that line you, Antonius — I shall speak truly — never once failed me when I was inquiring or asking, and very often you taught me what you used to observe in speaking.

\ciceroSection{98} Now, since both of you have opened the way of access to those very things we are seeking, and since Crassus was the first to begin this conversation, grant us this favour, that you set out subtly all that you feel about every kind of speaking. If we obtain this from you, I shall give great thanks, Crassus, to this wrestling-school of yours and your Tusculan villa, and I shall set this suburban gymnasium of yours far above that famous Academy and Lyceum.''

\ciceroSection{99} Then he: ``Nay, Sulpicius,'' he said, ``let us rather ask Antonius — who is both able to do what you ask, and accustomed to (so I hear you say). For I confess that I have always shrunk from this whole kind of speech, and have very often, as you yourself said a little earlier, refused you when you were eager and pressing. I did so not from pride or rudeness, nor because I did not wish to obey your most upright and excellent zeal — especially since I had recognised you as the one out of all most born and apt for speaking — but, by Hercules, from my unfamiliarity with that kind of disputation and from my ignorance of those things which are handed down, as it were, in an art.''

\ciceroSection{100} Then Cotta: ``Since the thing which seemed to us the hardest of all — that you should speak about these subjects at all, Crassus — we have attained, henceforth it will be our own fault if we let you go without your having unfolded everything that we shall ask.''

\ciceroSection{101} ``About those things, I take it,'' said Crassus, ``in the form usual in declarations of inheritance: \textit{which I shall know and shall be able}.'' Then he: ``As for what you cannot or do not know — which of us is so impudent as to demand to know it himself, or be able?'' ``On that condition then,'' said Crassus, ``provided I am free to deny that I am able when I shall not be able, and to confess that I do not know what I do not know — you may ask at your own pleasure.''

\ciceroSection{102} ``Well then,'' said Sulpicius, ``the thing of which Antonius just expounded what you think — we ask of you: do you suppose there is any art of speaking?'' ``What — do you,'' said Crassus, ``set me, as you might some idle and chattering little Greekling, perhaps learned and well-read, a little question on which I am to speak at my own pleasure? When have you ever supposed I cared about or thought of those things, and did not always rather laugh at the impudence of those men, who, when they had taken their seat in the school, would order any one of a great audience to ask whatever he wished?

\ciceroSection{103} The man they say first did this was Gorgias of Leontini, who seemed to be undertaking and professing something tremendous, when he proclaimed that he was prepared on every subject any single hearer wished to ask about. Afterwards they began to do this commonly, and do it to this day: that there is no matter so great or so unforeseen or so new that they do not profess they will say all that can be said about it.

\ciceroSection{104} If I supposed, Cotta, that you, or you, Sulpicius, wished to hear about such things, I should have brought up here some Greek to delight us with disputations of that kind; nor is it even now hard to do. For with M. Piso, the young man already given to this study, a man of the highest natural ability and most fond of us, there is the Peripatetic Staseas, a man indeed familiar to us and, as I see is agreed among experts, in his own kind the leader of all.''

\ciceroSection{105} ``What Staseas, what Peripatetic, do you tell me of?'' said Mucius. ``You must humour the young men, Crassus, who do not seek the daily babble of some Greek without practice, nor the chant of the schools, but, from the wisest and most eloquent of all men — and from one who is the leader, in counsel and in tongue, not in little books but in the greatest cases and in this home of empire and glory — they are inquiring after the opinion of the man whose footsteps they desire to follow.

\ciceroSection{106} For my part, I have always thought you a god in speaking, but never have I assigned you greater praise of eloquence than of human kindness — which it most befits you now to use, and not to flee that disputation to which two young men of excellent talent desire that you should come.''

\ciceroSection{107} ``I do then desire,'' he said, ``to comply with these men, and I shall not grudge to say briefly, in my way, what I feel about each subject. And first this — since I think it not lawful for me to neglect your authority, Scaevola — I answer: there seems to me to be either no art of speaking at all, or only a very thin one, and the whole contention among learned men rests on a dispute about the word.

\ciceroSection{108} For if art is so defined as Antonius set it out a little while ago — out of subjects seen through to the bottom, plainly known and separated from the dominion of opinion, and grasped by knowledge — there does not seem to me to be any art of the orator. For all the kinds of this our forensic speech are various and adapted to common, popular feeling.

\ciceroSection{109} But if those things which have been observed in the practice and treatment of speaking have been noted and marked by skilful and experienced men, defined in words, illustrated by kinds, distributed in parts (which I see could happen), I do not see why, if not by that subtle definition then by this common opinion, it should not seem to be an art. But whether it is an art, or some likeness of an art, it is in any case not to be neglected. Yet we must understand that there are certain other things, greater than this, for the attaining of eloquence.''

\ciceroSection{110} Then Antonius said he heartily agreed with Crassus, in that he neither embraced art so as did those who placed all the force of speaking in art, nor again rejected it whole, as most philosophers do. ``But I judge,'' he said, ``that you will gratify our friends here, Crassus, if you set out those things which you suppose can avail more than the very art for the speaking.''

\ciceroSection{111} ``I shall speak indeed, since I have begun, and I shall ask of you,'' he said, ``that you do not carry these absurdities of mine outside. I shall be moderate myself, lest I seem to be some master and craftsman, but rather one out of the number of men in the toga, a middling man with forensic experience and not wholly unschooled, who has not promised something on his own account but has stumbled by chance into your conversation.

\ciceroSection{112} For my part, when I was canvassing for office, I used to dismiss Scaevola from me as I went the rounds with my hand outstretched, telling him I wished to play the fool — that is, to canvass more flatteringly, which could not be done well unless it were done foolishly; and that he was the one man of all in whose presence I should least wish to be a fool. And now fortune has set him down as the witness and spectator of my fooleries: for what is more foolish than to speak about speaking, since the very speaking is never not foolish, save when it is necessary?''

\ciceroSection{113} ``Go on then, Crassus,'' said Mucius; ``the fault you fear, I shall stand surety for it.'' ``I think then,'' said Crassus, ``that nature first and natural ability bring the greatest force to speaking; and indeed it was not the principle and method of speaking that was lacking to those writers of art Antonius spoke of just now, but nature itself. For there must be certain swift movements of mind and ability — swift and sharp for invention, abundant for setting out and adorning, firm and lasting for memory.

\ciceroSection{114} And if anyone should think these things can be received by art (which is false: the matter would do excellently if these could be kindled or stirred by art; but they cannot be inserted and given by art — for they are all gifts of nature) — what shall I say of those things which assuredly are born with the man himself: a free tongue, a sound of voice, sides, strength, a certain shape and figure of the whole face and body?

\ciceroSection{115} For I do not say this in such a way as that art cannot polish some — for I am not unaware that good things can become better by teaching, and what is not best can in some way still be sharpened and corrected — but there are some who either stammer so in tongue, or are so out of tune in voice, or so coarse and rustic in face and bodily motion, that even if they have force in talent and art, still they cannot come into the number of orators. And there are some so apt in those very things, so adorned with the gifts of nature, that they seem not to have been born but to have been moulded by some god.

\ciceroSection{116} It is a great burden and great task to undertake and to profess oneself the one man, while all are silent, to be heard about the greatest matters in a great gathering of people. For there is hardly anyone present who does not see the faults in the speaker more sharply and keenly than the right things; so that whatever it is in which the speaker offends overwhelms even the things that should be praised.

\ciceroSection{117} Nor do I argue this so as wholly to deter young men, if they happen not to have something natural, from the pursuit of speaking. For who does not see that to Gaius Coelius, my own contemporary, a new man, that mediocrity in speaking which alone he had been able to attain was a great honour? Who does not understand that your contemporary Quintus Varius, a coarse and ugly man, by that very faculty, whatever it was, attained great popular favour in the state?

\ciceroSection{118} But because we are inquiring about the orator, we must shape him in our discourse with all faults stripped away and crowned with every praise. For if the multitude of suits, the variety of cases, this crowd and forensic barbarity gives a place even to the most faulty orators, we shall not for that reason abandon what we are seeking. Therefore in those arts where the necessary thing sought is not utility but a free delight of the mind, how diligently and how nearly fastidiously we judge! For there are no quarrels or disputes which compel men, as in the Forum to endure bad orators, so in the theatre to endure bad actors.

\ciceroSection{119} The orator must therefore see to it diligently, not that he shall satisfy those for whom it is necessary, but that he shall seem admirable to those who are free to judge. And, if you ask, I shall declare plainly what I feel, before men of closest intimacy, what hitherto I have always kept silent and thought ought to be kept silent: that those who speak best, and who can do it most easily and ornately, still seem to me almost shameless if they do not approach speaking with fear and are not flustered in beginning the speech — though that cannot happen.

\ciceroSection{120} For the more excellent any one is in speaking, the more he dreads the difficulty of speaking and the various outcomes of a speech and the expectations of men. But the man who can produce nothing worthy of the matter, of the name of orator, of the ears of men, even if he is moved when he speaks, still seems to me shameless. For we should escape the name of shamelessness not by being ashamed but by not doing what is unbecoming.

\ciceroSection{121} The man who is not ashamed — which I see in most — him I judge worthy not only of reproach but also of punishment. For my part, I am wont to notice in you, and very often experience in myself, both that I grow pale at the beginnings of speaking, and tremble in my whole mind and in all my limbs. As a young man, indeed, I was so unmanned at the beginning of a prosecution that I owe it as a great kindness to Quintus Maximus that he immediately dismissed the council, as soon as he had seen me broken and weakened by fear.''

\ciceroSection{122} Here all began to give signs of agreement among themselves and to converse. For there was a certain wonderful modesty in Crassus, which not only did not weigh against his speech but added to it through the recommendation of his uprightness. Then Antonius: ``Often, as you say,'' said he, ``I have observed, Crassus, both in you and in the rest of the foremost orators (although in my view no one has ever been your equal) that they are deeply moved at the beginning of speaking.

\ciceroSection{123} Looking for the cause of this — why it should be that the orator with most ability in him is the most afraid — I found these two reasons. One is that those whom practice and nature have taught understand that even to the foremost orators the outcome of speaking sometimes does not turn out as they wish; and so, not without justice, whenever they speak, they fear lest at that very moment the thing happen that on occasion can

\ciceroSection{124} happen. The other is this, of which I am wont often to complain: that men of other arts, when of established and approved standing, if they ever have done something less well than they are accustomed, are thought either to have been unwilling, or to have been kept from attaining what they knew through illness — `Roscius did not wish to act today,' they say, or `he had a bad stomach.' But in the orator, if any fault is observed, it seems a fault of stupidity.

\ciceroSection{125} And stupidity has no excuse, for no one is thought to have been stupid because he had a bad stomach or wished to be so. Whereby we undergo a yet heavier judgement in speaking; for as often as we speak, so often are we judged. And he who has gone wrong in gesture once is not at once thought ignorant of gesture; but as for him who is reproached for anything in his speaking, an opinion of his slowness either eternal or at least long-lasting takes hold.

\ciceroSection{126} As for what you said — that there are very many things which the orator, if nature had not given them, would not be much helped by a teacher in — I strongly agree with you, and on that point I most approved that great teacher, Apollonius of Alabanda. He, although he taught for pay, would not allow those whom he judged unable to come out as orators to lose their effort with him. He used to send them away, and was wont to drive each man towards the art for which he thought him fit, and to encourage him.

\ciceroSection{127} For it is enough, in mastering the other crafts, only to have the likeness of a man, and to be able to grasp in mind and store in memory whatever is handed down or even hammered in, even if one happens to be slower. There is no demand for nimbleness of tongue, swiftness of words, nor finally those things which we cannot fashion ourselves: face, look, voice.

\ciceroSection{128} But in the orator must be required the keenness of the dialecticians, the thoughts of the philosophers, the words almost of the poets, the memory of the jurists, the voice of the tragic actors, the gesture almost of the foremost actors. Therefore nothing in the human race can be found rarer than the perfect orator. For those things which the masters of single arts are approved for, if they have got them moderately, in the orator are not approved unless they are all of the highest order.''

\ciceroSection{129} Then Crassus: ``But see,'' he said, ``in the very thinnest and lightest craft, how much more diligence is brought to bear than in this matter, which is agreed to be the greatest. For I am wont often to hear Roscius say that he has thus far been able to find no pupil whom he could approve — not because there were not some commendable ones, but because, if there was any fault even small, he himself could not bear it. For nothing is so striking, nothing so durably fixed in memory, as something on which you have stumbled.

\ciceroSection{130} And so, that we may direct the praise of the orator by this likeness of the actor, do you not see that nothing is done by him save perfectly, nothing save with the highest grace, save in such a way as is becoming, and so as to move and please all? Therefore he has long since attained this — that whoever excels in any craft is called the Roscius of his kind. And I, in demanding this completion and perfection in the orator, from which I myself am far removed, act shamelessly. For I wish pardon for myself, while I do not pardon the rest. The man who cannot, who acts faultily, who, in short, is not becoming — he, as Apollonius bade, must, I judge, be thrust off to what he can do.''

\ciceroSection{131} ``You bid me, then, or this Cotta here,'' said Sulpicius, ``learn the civil law or military matters? For who can attain to those highest things, perfect in every kind?'' Then he: ``I myself,'' said he, ``because I have recognised in you a certain extraordinary and brilliant disposition for speaking, for that reason set out all these things; and I directed my speech not so much to deter those who could not, as to sharpen you, who can. And although I have seen the highest talent and zeal in each of you, still those things which lie in outward appearance, of which I perhaps said more than the Greeks are wont to say, are divine in you, Sulpicius;

\ciceroSection{132} for I do not seem to have heard anyone more apt either in motion of body or in very bearing and form, nor anyone with a fuller or sweeter voice. To those who have less of these from nature, it is still possible to attain that they use what they have moderately and skilfully, and so as not to be unbecoming. For unbecomingness is most of all to be avoided, and on this one matter it is least easy to give precept — not for me only, who speak as a single head of a household about these things, but for that very Roscius, whom I often hear say that the head of his art is becomingness; and yet that this is the one thing that cannot be handed over by art.

\ciceroSection{133} But if you please, let us turn the conversation elsewhere and speak in our own way for once, not in the rhetorician's.'' ``By no means,'' said Cotta. ``For now we must beseech you, since you keep us in this pursuit and do not dismiss us to another art, to set out to us, whatever it is, what you can do in speaking. We are not too greedy: we are content with that middling eloquence of yours; and we ask of you (so that we may not attain more than the little you have attained in speaking) — since you say that what is to be sought from nature is not too lacking to us — what besides you suppose we should take to ourselves?''

\ciceroSection{134} Then Crassus, smiling: ``What do you suppose, Cotta,'' said he, ``but zeal and a certain fire of love? Without which, just as in life no one ever does anything outstanding, so most certainly this thing you are after no one will ever attain. And I do not see that you need to be exhorted to it, since, while you are also troublesome to me, I see that you both burn too much already with the desire.

\ciceroSection{135} But indeed pursuits avail nothing for arriving anywhere unless that which leads and conducts you to the place you aim at has been recognised. So since you are laying on me a lighter burden, and you ask of me not about the orator's art, but about this faculty of mine, however small it is, I shall set out to you not some abstruse or very difficult or grand or weighty principle, but the customary practice of mine, which I once was wont to use when, as a young man, I was at liberty to spend my time in that pursuit.''

\ciceroSection{136} Then Sulpicius: ``O day, Cotta, much-wished for to us! For what I could never attain by entreaties or by laying snares or by spying — to be able not only to see what Crassus was doing for the sake of meditating or speaking, but even to suspect it from his secretary and reader Diphilus — that I hope we have got, and that now from the man himself we shall learn all those things we have long desired.''

\ciceroSection{137} Then Crassus: ``But, Sulpicius, I think that when you have heard, you will not so much wonder at these things which I shall say as judge that, when you wished to hear them, there was no reason why you wished it. For I shall say nothing recondite, nothing worthy of your expectation, nothing either unheard by you or new to anyone. For at the start — what is worthy of a free-born man liberally educated — I shall not deny that I have learned all those common, well-worn precepts: that the first duty of the orator is to speak with adaptation to persuasion;

\ciceroSection{138} that all speech is either about the question of an unbounded matter without designation of persons or times, or about a matter set in fixed persons and times.

\ciceroSection{139} In each kind, whatever comes into dispute is wont to be questioned, either as to whether it has been done; or, if done, of what kind it is; or further, by what name it should be called; or, what some add, whether it seems to have been done rightly.

\ciceroSection{140} Disputes arise also from the interpretation of writing — where what has been written is either ambiguous or contradictory, or in such a way that the writing differs from the intention. Under all these parts, certain proper arguments are set.

\ciceroSection{141} Of cases not separated from common questioning, some are concerned with judgements, some with deliberations; there is also a third kind, which is set in the praising or blaming of men. There are fixed topics which we use in judgements where equity is sought; others in deliberations, all directed to the advantage of those for whom we are giving counsel; others again in laudations, in which all things are referred to the dignity of persons.

\ciceroSection{142} The whole force and faculty of the orator is divided into five parts: that he ought first to find what to say; then to dispense and arrange the things found, not only in order, but also with a certain weight and judgement; then to clothe and adorn them with speech; afterwards to fence them about with memory; at last to deliver them with dignity and grace.

\ciceroSection{143} I had also learned and received that, before we speak about the matter, we must at the start win over the minds of those who hear; then the matter must be set out; afterwards the dispute established; then what we maintain must be confirmed; afterwards what is said on the other side must be refuted. At the end of the speech the points in our favour must be amplified and increased, and those for our adversaries weakened and broken.

\ciceroSection{144} I had also heard the precepts handed down concerning the ornaments of the speech itself: that we should speak first purely and in good Latin; then plainly and clearly; then with ornament; afterwards aptly to the dignity of the matter and as it were becomingly. And I had learned the precepts for each thing.

\ciceroSection{145} Indeed I had seen that art is brought to bear even on the very things most proper to nature; for I had tasted certain brief but well-practised precepts on delivery and on memory. For in these matters in the main consists all the doctrine of those craftsmen — which, if I should say in nothing helps, I should lie. For it has, as it were, certain things to remind the orator whither he should refer each matter, and looking on which he will less stray from whatever he has set before himself.

\ciceroSection{146} But I understand the force of all the precepts to be this: not that orators by following them have attained praise of eloquence, but rather that what eloquent men used to do of their own accord, certain men have observed and gathered. So eloquence is not born of art, but art of eloquence — which I do not, however, as I said, throw out: for it is, even if less necessary for speaking well, yet not unworthy of a free man's knowing.

\ciceroSection{147} And a certain practice must be undertaken by you — although you, indeed, have long been on the course. But for those who are entering the stadium and who can even now, by play-like exercise, learn beforehand and meditate upon those things which are to be done in the Forum as in a battle-line.''

\ciceroSection{148} ``This very thing,'' said Sulpicius, ``we wish to know. And yet we are eager to hear those things on art that have been run through briefly by you — although they too are not unheard by us. But that later: now we ask what you feel about practice itself.''

\ciceroSection{149} ``For my part,'' said Crassus, ``I approve those things you are wont to do — that, given a case set down similar to those brought into the Forum, you should speak with the closest accommodation to truth. But most men in this exercise only train their voice (and that not skilfully) and their lungs, and quicken the speed of the tongue, and delight in the abundance of words. And what deceives them is that they have heard that men, by speaking, are wont to be made fit to speak.

\ciceroSection{150} For there is also a true saying, that men attain to speaking badly most easily by speaking badly. Therefore in those very exercises, although it is useful even often to speak on the spur, that is yet more useful — having taken time to think about it — to speak more prepared and with more care. The chief thing, however — which, to speak truly, we very seldom do (for it is of great labour, which most of us shrink from) — is to write as much as possible. The pen is the best and most outstanding maker and master of speaking; and not without reason. For if a sudden, off-hand speech is easily surpassed by reflection and thought, the very reflection itself will surely be surpassed by constant and diligent writing.

\ciceroSection{151} For all the topics, whether of art or of some natural ability and prudence, that are at all in the matter on which we write — when we are inquiring and contemplating with the whole edge of our wit — show themselves and come up to meet us; and all the thoughts and all the words of every kind which are most striking must rise and yield themselves to the point of the pen. Then the very arrangement and shaping of words is brought to perfection in writing — not in poetic, but in a certain oratorical rhythm and measure.

\ciceroSection{152} These are the things that produce the cries and the wonderings at good orators; and no one will attain them, unless he has long and copiously written, even if he has exercised himself most vehemently in these sudden speakings. And he who comes from the habit of writing to speaking, brings this faculty: that even if he speaks suddenly, what he says still seems like things written. And even if he sometimes brings something written into his speech, when he has departed from it, the rest of the speech follows in a similar manner.

\ciceroSection{153} As in a swiftly moving ship, when the rowers have ceased to row, the ship still keeps her motion and her course, although the impulse and stroke of the oars has been intermitted; so in continuous speech, when written things fail, the rest of the speech still keeps an equal course, driven on by the likeness and force of the things written.

\ciceroSection{154} In daily exercises, indeed, as a young man I was wont to set myself most of all that exercise which I knew our enemy Gaius Carbo had been wont to use: that, with the most weighty verses set forth, or some speech read up to the point I could grasp by memory, I should declare that very matter which I had read in other words, the most choice ones I could pick. But afterwards I noticed this fault in this practice — that those words which were the most proper to each thing and most ornate and best, either Ennius (if I exercised myself on his verses) or Gracchus (if I happened to set his speech before me) had got hold of. So if I used the same words, it gained me nothing; if I used others, it actually hindered me, because I was getting accustomed to use less suitable ones.

\ciceroSection{155} Afterwards I came to like, and used as a young man, this practice: to render Greek speeches of the foremost orators. By reading these I attained this — that, when I rendered into Latin what I had read in Greek, I not only used the best words, and yet familiar ones, but even, by imitating, expressed certain words new to our language, provided they were suitable.

\ciceroSection{156} The motions and exercises of voice and breath, of the whole body, and of the tongue itself, need labour rather than art. In which, careful regard must be had to whom we imitate, whom we wish to be like. We must look not only at orators but also at actors, lest by bad habit we come to some misshapenness and crookedness.

\ciceroSection{157} The memory must also be exercised, by learning by heart as much as possible both of our own writings and of others'; and in that exercise it does not displease me, if you are accustomed, to bring in even that method of places and images which is handed down in the art. The diction must then be brought out of this private and shaded exercise into the middle of the column, into the dust, into the shouting, into the camp and the battle-line of the Forum. The looks of men must be undergone and the strength of the talent put to the test, and that meditation hidden away must be brought into the light of truth.

\ciceroSection{158} The poets too must be read; the histories made known; the teachers and writers of all the good arts chosen and turned over and, for exercise, praised, interpreted, corrected, blamed, refuted. We must argue every matter on opposing sides, and whatever there is in any subject which can seem probable must be drawn out and spoken.

\ciceroSection{159} The civil law must be thoroughly learned, the laws known, all antiquity grasped, the senatorial usage, the discipline of the commonwealth, the rights of the allies, treaties, agreements, the cause of the empire learned. There must also be sipped, from every kind of urbanity, a certain wit's charm, with which, as with salt, all speech may be sprinkled. I have poured out to you everything I felt, which perhaps any head of a household whom you might have grabbed from a circle of bystanders would have answered to your same questions.''

\ciceroSection{160} When Crassus had said this, silence followed. But although those present thought what had been said sufficient for what was proposed, they nevertheless felt that he had concluded much sooner than they themselves would have wished. Then Scaevola: ``What is it, Cotta?'' he said. ``Why are you silent? Does nothing come into your mind which besides you would ask of Crassus?''

\ciceroSection{161} ``By Hercules,'' he said, ``that very thing I am attending to. For so great was the rush of words and so the speech flew off, that I have looked upon its force and impulse but scarcely seen its footsteps and its tread; and as if I had come into some rich and crammed house, with the clothing not unfolded, the silver not set forth, the paintings and statues not displayed in the open, but with all these many magnificent things heaped up and shut away — so just now in Crassus's speech I have seen the riches and ornaments of his talent through certain wrappings and coverings, but when I wished to gaze on them, hardly was the chance of looking given. So I cannot say that I am wholly ignorant of what he possesses, nor that I plainly know it and have seen it.''

\ciceroSection{162} ``Why then,'' said Scaevola, ``do you not do the same thing you would do if you had come into some house full of ornaments, or villa? If they were laid aside, as you say, and you were eager to look — would you not, free of doubt, ask the master to have them brought forth, especially if he were a friend? Likewise now you will ask Crassus to bring out into the light that abundance of his ornaments, which built up in one place we have, as it were through a lattice, glanced at as we passed, and to set each in its own place.''

\ciceroSection{163} ``I beg you to do this, Scaevola,'' said Cotta. ``For shame holds me and Sulpicius back from inquiring of the most weighty man of all, who has always despised this kind of disputation, the things which to him perhaps may seem the elements taught to boys. But you grant us this favour, Scaevola, and bring it about that Crassus enlarge for us and unfold the things he has hemmed in and packed most narrowly in his speech.''

\ciceroSection{164} ``By Hercules,'' said Mucius, ``before this I wished it more for your sake than for my own; for I did not so much desire this discussion from Crassus as I am pleased by his speech in cases. But now, Crassus, I ask you also for my own sake — that, since we have so much leisure as has not befallen us for a long time — that you not be reluctant to build up the work which you have begun. For I see the shape of the whole business better and greater than I supposed, which I strongly approve.''

\ciceroSection{165} ``In truth,'' said Crassus, ``I cannot wonder enough that even you, Scaevola, desire these things — which I do not hold as those who teach do, nor are of such a kind that, if I held them best, they would be worthy of your wisdom and ears.'' ``Do you say so?'' said the other. ``If you think hardly even those common and well-worn things ought to be heard at this age — can we still neglect those things which you said the orator ought to know: about the natures of men, about morals, about those reasonings by which the minds of men are stirred up and held back; about history, about antiquity, about the administration of the commonwealth, lastly about our own civil law? For I knew that all this knowledge and abundance of subjects was set in your wisdom; but in the orator's equipment I had never seen so handsome a furniture.''

\ciceroSection{166} ``Can you, then, Crassus,'' he said, ``— to leave aside countless other immense things and come to your own civil law — think those orators who, when Hypsaeus in the loudest voice and with the most words contended with the praetor M. Crassus that the man he was defending should be allowed to lose his case, and Cn. Octavius, a consular, refused with no shorter speech that the adversary should lose his case, and that the man for whom he himself spoke should be freed from a shameful judgment of guardianship and from all annoyance through his opponent's stupidity — those whom Publius Scaevola, smiling and choking with anger, watched for many hours when he was hurrying off to the Campus?''

\ciceroSection{167} ``I, indeed, those men,'' he said — ``for I remember Mucius telling me — should hardly judge worthy not only of the name of orator but even of the Forum.'' ``And yet those advocates,'' said Crassus, ``did not lack eloquence or any rationale or abundance of speech — but the science of civil law: for one was demanding by his action more than what the law in the Twelve Tables had permitted, so that, if he attained it, his client would lose the case; the other thought it unfair that more was being claimed against his client than what was in the action; nor did he understand that, if it had been so done, the adversary would have lost the suit.

\ciceroSection{168} What of this? In these last few days, while we were sitting on the bench of our friend Quintus Pompeius, the urban praetor, did not a man of the number of the eloquent demand that the old and customary exception be granted to him from whom payment was being claimed: \textit{the day of which money has been due}? — without understanding that this was provided in the plaintiff's interest, so that, if the denier had proved to the judge that the money had been claimed before it had begun to be owed, the plaintiff might, when he claimed again, not be excluded by the exception \textit{because that matter has come to judgment before}.

\ciceroSection{169} What, then, can be done or said more disgraceful than this — that he who has taken on this character: to defend the disputes and cases of his friends, to come to the help of the labouring, to heal the sick, to raise up the broken — that this man should so trip in the smallest and slightest things that he seems to some pitiable, to others ridiculous?

\ciceroSection{170} I judge that our kinsman, that Publius Crassus the Rich, an elegant and accomplished man in many other things, must especially be brought forward and praised for this: that, although he was Publius Scaevola's brother, he was wont very often to say to him both that the latter could not satisfy his art in the civil law unless he had taken to himself an abundance of speaking — which his son, the man who was consul with me, attained — and that he himself had not begun to handle and plead his friends' cases before he had learned the civil law.

\ciceroSection{171} What of that great Marcus Cato? Was he not at once both as eloquent as the times and the age could bear in this state, and most expert of all in civil law? I speak more modestly of this matter for some time now, because there is present a man supreme in speaking, whom I most admire as a unique orator; but yet he too has always despised the civil law.

\ciceroSection{172} But, since you have wished to be partners in my opinion and feeling, I shall hide nothing, and as far as I shall be able, I shall set out to you what I think of each thing. The incredible, almost singular and divine force of Antonius's talent seems able, even when stripped of this knowledge of the law, easily to defend itself with the other arms of prudence; therefore let him be excepted by us. But the rest I shall not hesitate, in my opinion, first to condemn for slothfulness, then also for shamelessness.

\ciceroSection{173} For to flit about the Forum, to be at home in the law and at the praetors' tribunals, to take part in private suits over great matters in which the contest is often not about fact but about equity and right, to throw oneself about in the suits of the Hundred Men — in which there are tossed about the rights of usucapion, guardianship, kinship, agnation, alluvial gain, eddying loss, bond, mancipation, walls, lights, eaves-drips, broken or valid wills, and countless other things — when one altogether does not know what is one's own, what another's, by what difference, in short, anyone is a citizen or a foreigner, a slave or a free man — this is mark of a peculiar shamelessness.

\ciceroSection{174} That, indeed, is laughable arrogance: to confess oneself unskilled in smaller boats, but to have learned to steer galleys of five oars or even larger. When you are tripped up in a circle of bystanders by your opponent's little stipulation, when you set your seal to your client's tablets in which is written the very thing whereby he is taken — am I to think any larger case should be entrusted to you? More quickly, by Hercules, would the man who has overturned a two-oared little boat in harbour be the steersman of the Argonauts' ship in the Euxine Sea.

\ciceroSection{175} What of this? If even small cases are not — but often great ones — in which the contest is about civil law, what kind of face is on that advocate who dares come to such cases without any knowledge of the law? What case could there be greater than that famous one of the soldier? When false news of his death had come home from the army, and his father had believed it and changed his will, and made the man who pleased him heir, and had himself died — when the soldier returned home and brought action by law for the paternal inheritance, the matter was referred to the Hundred Men: and the son had been, by the will, struck out as heir. In that case it was asked of the civil law whether a son could be excluded from his father's goods, whom the father in his will had neither named heir nor expressly disinherited.

\ciceroSection{176} What of this? In the case the Hundred Men decided between the patrician Marcelli and the patrician Claudii — when the Marcelli claimed that the inheritance of one and the same man had returned to themselves by descent through a freedman's son, and the patrician Claudii by clan — was it not in that case necessary for the orators to speak about the whole law of descent and clan-right?

\ciceroSection{177} What of this — likewise contested in a Hundred-Men trial, as we hear: when one had come to Rome in exile, who had at Rome the right of exile if he attached himself to someone as patron, and had died intestate — was not in that case the law of attachment, sufficiently obscure and unknown, opened in the trial and brought to light by the patron?

\ciceroSection{178} Indeed, just lately, when I was defending the case of Gaius Sergius Orata against this our Antonius in a private trial, was not all our defence rooted in law? For when Marcus Marius Gratidianus had sold Orata's house and had not stated in the conveyance that a certain part of the house was burdened with a servitude, we maintained that whatever inconvenience was attached to the conveyance, the seller, if he had known of it and had not declared it, was bound to make good.

\ciceroSection{179} In which kind of matter our friend Marcus Buculeius, a man not foolish in my judgement and very wise in his own, and not averse to the study of law, recently in some way erred in a similar matter. For when he was selling a house to Lucius Fufius, he kept by the conveyance the lights as they then were. But Fufius, as soon as building was begun in some part of the city which could only just be seen from those houses, brought action against Buculeius — for he supposed that, however small a particle of sky was blocked, however far away, the lights were altered.

\ciceroSection{180} What of this? The most famous case of M.' Curius and Marcus Coponius lately before the Hundred Men — with what concourse of men, with what expectation was it pleaded? When Quintus Scaevola, my contemporary and colleague, a man of all the most learned in the discipline of the civil law, the keenest in talent and prudence, the most polished and subtle in speech, and, as I am wont to say, the most eloquent of the law-skilled and the most law-skilled of the eloquent, defended the rights of wills out of the written word and denied that, unless a posthumous child were both born and had died before he came of age, the man could be heir who had been instituted heir after the posthumous child should be born and dead — and I, on the other side, defended that the maker of the will had this then in mind: that, if there were no son to come of age, M.' Curius should be heir — did either of us in that case keep from working, in authorities, in examples, in the formulas of wills, that is, in the very middle of the civil law?

\ciceroSection{181} I leave aside more examples of the greatest cases, which are countless. It can often happen that cases concerning our very life turn on the law. For if Gaius Mancinus, a most noble and excellent man and consular — when on account of the hatred for the Numantine treaty the \textit{pater patratus} by senatorial decree had handed him over to the Numantines and they had not received him, and Mancinus had afterwards returned home and had not hesitated to enter the senate — Publius Rutilius, son of Marcus, tribune of the plebs, ordered him to be led out, on the ground that he was no citizen, because it was so handed down by tradition that whoever had been sold by his father or by the people, or had been handed over by the \textit{pater patratus}, had no \textit{postliminium}: what greater contest can we find in all civil affairs, than that about a man's status, citizenship, freedom, the life of a consular man — especially when this rested not on some criminal charge, which he could deny, but on civil law?

\ciceroSection{182} On a similar kind of question, of a lower order, if anyone of an allied people had served as a slave among us and had freed himself and afterwards returned home, it was asked among our ancestors whether he had returned to his own people by \textit{postliminium} and lost our citizenship.

\ciceroSection{183} What of this? About freedom, than which no judgment can be weightier, can there not be a contest from civil law, when it is asked whether the man who has been registered by the wish of his master is free at once, or only when the lustrum has been closed? What of the case which came in usage in our fathers' memory: when a head of a house had come to Rome from Spain leaving his pregnant wife in the province, and had taken another wife at Rome, and had not sent notice of divorce to the first, and had died intestate, and a son had been born to each — was a slight matter brought into contest, when there was the question concerning two citizens' status: about the boy who had been born of the second, and about the mother — for if it were judged that, by certain words, divorce was made with the first not by new marriage, she would be reckoned in the place of a concubine?

\ciceroSection{184} A man ignorant of these and like rights of his state, walking erect and tall, with cheerful and ready face and look, looking now this way now that, ranging through the whole Forum with a great crowd, holding out and offering himself as a defence to clients, a help to friends, a light of his talent and counsel almost to all citizens — should he not be reckoned among the chief disgraces?

\ciceroSection{185} And since I have spoken of shamelessness, let us also chastise men's slothfulness and laziness. For if that knowledge of the law were great and difficult, the magnitude of the usefulness still ought to drive men to take up the labour of learning. But, immortal gods, I should not be saying this in Scaevola's hearing if he himself were not wont to say that the knowledge of no art seemed to him easier.

\ciceroSection{186} It is for certain reasons judged otherwise by most: first, because those ancients who presided over this knowledge, in order to keep and increase their power, did not wish their art to be made common; then, after it was published, the actions having first been set out by Cn. Flavius, there were none who would arrange those laws skilfully and put them together in classes. For there is nothing which can be reduced to art unless he who would set up the art of those things he holds, has first that knowledge whereby he can make an art out of those subjects whose art does not yet exist.

\ciceroSection{187} I see that, in wishing to speak briefly, I have spoken a little obscurely; but I shall try and say more plainly, if I can. Almost all that is now confined within arts was once scattered and dispersed: as in music, the rhythms and tones and modes; in geometry, lines, shapes, intervals, magnitudes; in astronomy, the conversion of the heaven, the risings and settings and motions of the stars; in grammar, the going-over of the poets, the knowledge of the histories, the interpretation of words, a certain manner of pronunciation; lastly, in this very rationale of speaking, the inventing, the adorning, the arranging, the remembering, the delivering — once unknown to all, and seemed widely diffuse.

\ciceroSection{188} An art was therefore brought in from outside, from another kind of subject, which the philosophers wholly take to themselves, which would glue together a thing dissolved and torn apart and bind it together with a certain method. Let then this be the goal in civil law: the preservation of the lawful and customary equity in the affairs and cases of the citizens.

\ciceroSection{189} Then the kinds must be marked and reduced to a fixed and small number. The genus is what, by a certain community with itself but with the species differing, embraces two or more parts. Parts are what are placed under those genera from which they flow. And of all things which are either names of genera or of parts, the force they have must be expressed by definitions; for definition is a brief and bounded explanation of those things which belong properly to that thing which we wish to define.

\ciceroSection{190} I should add examples to these subjects, were I not aware among whom this speech is being held. Now I shall sum up briefly what I have proposed: if either I shall be allowed to do what I have long been thinking of, or some other, while I am hindered, shall take it up, or, when I am dead, shall accomplish — that he should first arrange all the civil law in classes, which are very few; then divide as it were certain limbs of those classes; then declare the proper force of each by definition — you will have the perfect art of civil law, more great and rich than difficult and obscure.

\ciceroSection{191} And meanwhile, while these scattered things are being collected, you may, by gathering and gleaning here and there, be filled with a sufficient knowledge of the civil law. Do you not see that the Roman knight, a man of the keenest talent of all, but least learned in the other arts, Gaius Aculeo, who lives with me and always has lived, holds the civil law in such a way that, when you have set Crassus aside, no one of the most expert is preferred to him?

\ciceroSection{192} For all things are set before our eyes, established in daily use, in the meeting of men, in the Forum; and they are not contained in so many letters or great volumes. For the same things are first set out by many, then with a few words changed, written by the same writers more often.

\ciceroSection{193} Indeed there is added to make the civil law more easily perceived and known — what most men least suppose — a certain wonderful pleasure and delight in the learning. For if anyone is delighted by these Aelian studies, there is the most ample image of antiquity in all civil law, in the books of the pontiffs, and in the Twelve Tables; for the ancient old age of words is recognised, and certain kinds of actions declare the custom and life of our ancestors. Or if anyone is delighted by civic knowledge — which Scaevola does not think proper to the orator, but to a kind of prudence of another sort — he will see this whole, with all the uses and parts of the state described, contained in the Twelve Tables. Or if that overpowering and glorious philosophy delights anyone — I shall speak more boldly — it has these springs of all its disputations: those things contained in the civil law and the laws. From them we see that dignity is most to be desired, since true virtue and honest labour is honoured with offices, rewards, splendour, while men's vices and frauds are punished with losses, disgrace, chains, beatings, exiles, death.

\ciceroSection{194} And we are taught, not by endless and contention-filled disputations but by the authority and nod of the laws, to keep our lusts tamed, to restrain all our desires, to preserve our own, to keep our minds, eyes, and hands from what is another's.

\ciceroSection{195} Let all rage if they will, I shall say what I think: the little book of the Twelve Tables alone, by Hercules, seems to me to surpass the libraries of all the philosophers — to anyone who looks upon the springs and heads of laws — both in the weight of authority and in the abundance of utility.

\ciceroSection{196} And if our country delights us, as it most ought (for the force and nature of which is so great that the wisest of men preferred his Ithaca, fixed like a little nest among the roughest crags, to immortality) — with what love at last must we be inflamed for such a country as ours, which is alone in all the lands the home of virtue, of empire, of dignity? Of which the mind, the custom, the discipline ought first to be known to us, both because she is the parent of us all, and because so much wisdom must be reckoned to have been in establishing her law as there was in winning the great resources of empire.

\ciceroSection{197} You will derive also that joy and pleasure from the knowledge of the law: that you will most easily understand how much our ancestors excelled the rest of the nations in prudence, if you should wish to compare with their Lycurgus, Draco, and Solon, our laws. For it is incredible how unformed and well-nigh ridiculous all civil law is besides this of ours; about which I am wont to say much in everyday conversations, when I prefer the prudence of our men to all others and most of all to the Greeks. For these reasons, Scaevola, I had said that the knowledge of the civil law was necessary for those who wished to be perfect orators.

\ciceroSection{198} Now of itself, who is unaware how much honour, favour, and dignity it brings to those who preside over it? And so, just as among the Greeks the lowest men, drawn by a small fee, offer themselves as helpers to orators in trials — those who among them are called \textit{pragmatikoi} — so in our state, on the contrary, the most eminent and most distinguished man — like the one who, on account of this knowledge of civil law, was called by the supreme poet, ``a man of remarkable heart, the canny Aelius Sextus,'' and many besides who, when by their natural ability they had won themselves dignity, brought it about that, in giving legal opinions, they had more weight by authority even than by talent itself.

\ciceroSection{199} For the celebrating and adorning of old age, what more honourable refuge can there be than the interpretation of law? For my part, I have prepared this support for myself even from young manhood, not only for the use of forensic cases but also for the grace and ornament of old age, so that, when my powers should begin to fail me — which time is now nearly at hand — I might rescue my house from solitude. For what is more splendid than that an old man, having gone through the offices and burdens of public life, should be able by his own right to say what in Ennius the Pythian Apollo says: that he is the one from whom — if not peoples and kings — at least all his fellow-citizens seek counsel, when in doubt about the highest matters: ``whom I, by my help, send away certain instead of uncertain, and equipped with counsel, that they may not deal rashly with troubled affairs:''

\ciceroSection{200} for without doubt the house of the jurisconsult is the oracle of the whole state. Witness this Quintus Mucius's door and entrance-hall, which, in his most infirm health and now broken old age, is daily thronged by the greatest crowd of citizens and the splendour of the highest men.

\ciceroSection{201} The reasons why I think also that the public laws, which are proper to the state and to the empire, and the records of past affairs and the examples of antiquity, ought to be known to the orator, do not now require a long discourse. For just as in cases and trials of private affairs the speech must often be drawn from civil law, and so, as we said before, the science of civil law is necessary to the orator, in like manner in public causes — of trials, assemblies, the senate — all this memory of antiquity, this authority of public law, this rationale and science of governing the commonwealth, ought, as a sort of subject-matter, to lie under those orators who move in public life.

\ciceroSection{202} For we are not seeking, in this conversation of ours, some pleader of cases or shouter or wrangler, but the man who first is a high-priest of that art (which, since nature itself gives so great a faculty to man, is yet thought to have a god as its author, so that the very thing which was proper to man should seem given to us not by our own getting but by divine bestowal); then, the man who can, equipped not so much with the herald's wand as with the orator's name, move unhurt even amid the weapons of enemies; then, the man who can by his speech subject the crime and fraud of the guilty to the hatred of the citizens and bind it down with punishment; the same man, by the protection of his talent, set free innocence from the courts' punishment; the same man rouse a flagging and faltering people either to nobility, or lead it back from error, or kindle it against the wicked, or, when stirred against the good, soothe it; lastly, the man who can, by speaking, either rouse or calm whatever movement in the minds of men the matter and case may require.

\ciceroSection{203} If anyone supposes this force has been set out either by those who have written about the rationale of speaking, or can be set out by me in so brief a space, he is greatly mistaken, and sees not only my own want of knowledge but the very greatness of the matter. For my part, since you so wished, I have judged it necessary to point out to you the springs from which you should draw, and the very paths — not as if I myself were your guide, which is both endless and unnecessary, but only that I might point out the way and, as is wont to be done, point my finger at the springs.''

\ciceroSection{204} ``In my judgment,'' said Mucius, ``you have done enough and more for the studies of these young men, if they are studious. For as Socrates, they say, was wont to say, his work was complete if anyone was sufficiently roused by his exhortation to the pursuit of recognising and grasping virtue: for those whom this had persuaded — that they would prefer nothing to being good men — to them the rest of the teaching was easy. So I understand: if you wish to enter into those things which Crassus has opened up by his speech, you will most easily come to what you desire, by this entrance and gate now thrown open.''

\ciceroSection{205} ``Most welcome and most pleasing to us are these things,'' said Sulpicius. ``But we still ask a few things, and especially those which were run through very briefly by you, Crassus, about the very art, when you confessed that you neither despised it nor failed to learn it. If you should speak about these a little more broadly, you would fulfil our every expectation of long-standing desire. For now, what we should be devoted to we have learned, which is itself a great thing; but we wish to learn the paths and rationale of these things.''

\ciceroSection{206} ``What if,'' said Crassus, ``— since I, that I may keep you with me the more easily, have rather followed your wish than either my custom or my nature — we ask Antonius to set out for us those things which he holds in himself but has not yet brought forth, and from which he has long since complained one little book has slipped from him? — and to declare to us those mysteries of speaking?'' ``As you please,'' said Sulpicius; ``for as Antonius speaks, we shall sense your own opinion as well.''

\ciceroSection{207} ``I ask you then,'' said Crassus, ``since this burden, Antonius, is laid on us men of this age by the studies of these young men, set out what you feel about those things which you see asked of you.'' ``I see and feel,'' said Antonius, ``that I am clearly caught: not only because those things are asked of me of which I am ignorant and unfamiliar, but because what in cases I am wont strongly to flee — to follow you, Crassus — those friends will not let me now avoid.

\ciceroSection{208} But I shall enter on what you wish more boldly because I hope the same thing will fall out for me in this disputation as is wont to happen in my speaking: that no ornate speech is to be expected. For I shall not speak about an art which I have never learned, but about my own custom; and those very things which I have written down in my notebook are of such a kind: not handed down to me by some teaching, but treated in the practice of affairs and cases. And if these things, before such learned men as you, do not win approval, lay the blame on your own unfairness, who have asked of me what I should not know; praise my own ease, when, drawn on not by my own judgment but by your own zeal, I have answered without complaint.''

\ciceroSection{209} Then Crassus: ``Go on, Antonius. There is no danger that you should say anything otherwise than so prudently that none of us shall regret to have driven you to this conversation.'' ``I shall go on then,'' he said, ``and I shall do what I think ought to be done at the start in all disputations: that what the matter under dispute is should be made plain, lest the speech be forced to wander and stray, if those who differ among themselves do not understand that the very thing they are dealing with is the same.

\ciceroSection{210} For if it should chance to be asked what the art of the general is, I should think we should establish at the start who the general is. When he had been established as a kind of administrator of waging war, then we should add about the army, the camp, the line of march, the joinings of standards, the storming of towns, the supply, the laying and avoiding of ambushes, the rest which were proper to the management of war. Those who were endowed in mind and knowledge with these things, I should say were generals; and I should use the examples of the Africani and the Maximi, and should name Epaminondas and Hannibal and men of that kind.

\ciceroSection{211} If it were asked who that man was, who had brought to bear his experience and knowledge and zeal upon the directing of the commonwealth, I should define him thus: he who held those things by which the usefulness of the commonwealth was secured and increased, and used them — him I should hold to be the director of the commonwealth and the author of public counsel. And I should set forth Publius Lentulus the famous \textit{princeps}, Tiberius Gracchus the father, Quintus Metellus, Publius Africanus, Gaius Laelius, and countless others both from our own state and from the rest.

\ciceroSection{212} If it were asked who could truly be called a jurisconsult, I should say: he who is expert in the laws and the custom by which private men in the state act, both for giving opinions, and for litigating, and for warning. And from that kind I should name Sextus Aelius, M.' Manilius, Publius Mucius. And, to come now to the lighter studies of the arts, if a musician, a grammarian, a poet were sought, I could similarly set out what each of them professes, and beyond what should be required of him. Even of the philosopher himself — who alone of his own force and wisdom professes well-nigh everything — there is yet a certain definition: that he who studies the force, nature, and causes of all things divine and human, and holds and pursues all the rationale of living well, should be called by this name.

\ciceroSection{213} As for the orator, since concerning him we are inquiring, I do not make him the same as Crassus does, who seemed to me to embrace the knowledge of all things and arts in the one office and name of the orator. I think the orator is rather he who can use words pleasing to hear and thoughts adapted for proving, in forensic and common cases. Him I call orator, and I wish him to be furnished besides with voice, with delivery, and with a certain charm.

\ciceroSection{214} But our friend Crassus seemed to me to draw the orator's faculty by no boundaries of that art, but by the well-nigh limitless boundaries of his own talent. For he handed over to the orator, by his own opinion, the helm of governing states. In which it seemed to me marvellous, Scaevola, that you conceded this to him, when very often the senate has agreed with you, who speak briefly and unpolishedly, on the greatest matters. M. Scaurus indeed — whom I hear is at his country house not far off, a man most knowing in the directing of the commonwealth — if he hears that this authority of his weight and counsel is being claimed by you, Crassus, on the ground that you say it is the orator's own, even now, I suppose, would come here and dismay this loquacity of ours by his very face and look. Although he is not at all to be despised in speaking, his strength rests on prudence in great matters more than on the art of speaking.

\ciceroSection{215} Nor, even if anyone can do both, is the man on that account an orator, who is the author of public counsel and a good senator; nor, on the other hand, has this skilful and eloquent man, if he is also outstanding in directing the state, attained that knowledge by the abundance of speaking. Those faculties stand far apart from one another and are far different and separate, nor by the same reason and method did M. Cato, P. Africanus, Q. Metellus, C. Laelius (all of whom were eloquent) adorn at once their own speech and the dignity of the commonwealth. For it is not forbidden either by the nature of things or by any law and custom that single men should know more than single arts.

\ciceroSection{216} Therefore, even though Pericles was the most eloquent man at Athens, and was for many years the leader of public counsel in that state, both faculties are not on that account to be reckoned of the same man and art. Nor, even if Publius Crassus was both eloquent and law-skilled, is the science of civil law on that account included in the faculty of speaking.

\ciceroSection{217} For if it is so — that whoever excels in some art and faculty, if he should take up another art also, brings it about that what he knows besides should seem to be a part of that in which he excels — by that line of reasoning we may say that to play ball well and at the twelve-line game is proper to the civil law, since Publius Mucius did each excellently. By the same reasoning let those who are called by the Greeks \textit{phusikoi} also be called poets, because the natural philosopher Empedocles wrote an outstanding poem. But this not even the philosophers themselves dare to say — who would have all things as their own and possessed by themselves — that geometry or music is the philosopher's, because everyone confesses that Plato was supreme in those arts.

\ciceroSection{218} And if now you would prefer to subjoin all the arts to the orator, it is more tolerable to put it like this: that since the faculty of speaking ought not to be meagre and bare but sprinkled and adorned by a kind of pleasant variety of many subjects, it is the business of a good orator to have heard much, seen much, run over much in mind and thought, much also in reading; and not to have possessed these things as his own but to have tasted them as another's. For I confess that he must be skilled and a beginner in nothing and not a stranger and a guest in pleading.

\ciceroSection{219} Nor am I disturbed, Crassus, by those tragedies of yours which the philosophers most use, when you said that no one of them can either kindle the minds of his hearers by speaking or extinguish them when kindled — wherein the force and greatness of the orator is most discerned — unless one has thoroughly seen the nature of all things and the morals and reasonings of men. In which study we see that the whole lifetimes even of the most ingenious and leisure-rich men are worn out. The abundance and magnitude of their knowledge and art I not only do not despise, but vehemently admire. To us, however, who move in this people and Forum, it is enough to know and to say about the movements of minds those things which do not jar with the morals of men.

\ciceroSection{220} For what great and weighty orator, when he wished to make a judge angry at his adversary, ever hesitated for the reason that he did not know what anger was — whether a fervour of mind or a desire to punish pain? Who, when he wished to mix and stir other movements of the soul, by speaking, in judges or in the people, said the things that the philosophers are wont to say — who in part deny that there should be any movements at all in souls, and that those who stir them up in the minds of judges commit a wicked crime; in part, who wish to be more tolerable and to come nearer to the truth of life, say that the movements ought to be very moderate and rather slight?

\ciceroSection{221} The orator, however, makes by words much greater and more bitter all those things which are reckoned bad and troublesome and to be avoided in the common practice of life; and likewise he amplifies and adorns by speaking those things which are commonly thought to be sought after and desirable. Nor does he wish to seem so wise among fools that those who hear think him either inept and a Greekling, or, even if they highly approve his talent, marvel at the wisdom of the orator and take it ill that they themselves are fools.

\ciceroSection{222} Rather, he so passes through their souls, so handles the senses and minds of men, that he does not need the philosophers' descriptions, nor inquire in his speech whether the highest good is in the soul or in the body, defined by virtue or by pleasure, or whether these things can be joined and coupled with one another; or whether — as has seemed to some — nothing certain can be known, nothing plainly recognised and grasped. The science of these things I confess is great and manifold and there are many copious and various reasonings.

\ciceroSection{223} But we are seeking something else, far else, Crassus: we have need of a man keen and skilled by nature and use, who shall sagaciously trace out what his fellow-citizens, and the men whom he wishes to persuade something by speaking, think, feel, suppose, expect. He must hold the veins of every kind, age, and order, and taste the minds and senses of those before whom he is to plead or shall plead.

\ciceroSection{224} The books of the philosophers, however, let him reserve for the rest and leisure of such a Tusculan villa as this — lest, if at any time he must speak about justice and good faith, he borrow from Plato. Plato, when he thought these things must be set in words, fashioned in his books a certain new state — to such a degree did the things which he judged should be said about justice diverge from the use of life and the morals of states.

\ciceroSection{225} If those things were approved among peoples and states, who would have conceded to you, Crassus — most distinguished, most ample man and chief of the state — to say in the largest assembly of your fellow-citizens what you said? ``Snatch us out of miseries, snatch us out of the jaws of those whose cruelty cannot be filled save by our blood. Do not allow us to be slaves to anyone, save to all of you, to whom we both can and ought to be slaves.'' I leave aside ``miseries,'' in which, as those men say, the strong man cannot be; I leave aside the ``jaws'' from which you wish to be snatched, lest by an unfair judgment your blood be sucked out — what they say cannot happen to the wise man. But to be slaves — not only you, but the whole senate, whose cause you were then pleading — you dared to say it?''

\ciceroSection{226} ``Can virtue, Crassus, be a slave — by those authorities of yours, whose precepts you embrace in the orator's faculty? Virtue is always and alone free, and even if bodies are taken in arms or bound in chains, virtue ought still to keep its own right and the unpunished freedom of all things. As to what you added — that the senate not only could but even ought to be a slave to the people — what philosopher could approve this, however soft, however slack, however nerveless, however referring all things to bodily pleasure and pain — that the senate should be a slave to the people, to whom the people itself had handed over, as it were certain reins, the power of moderating and ruling itself?

\ciceroSection{227} Therefore, although I judged this to be said divinely by you, P. Rutilius Rufus, a learned man devoted to philosophy, was wont to say that it had been said not only too imprudently but even disgracefully and shamefully. The same man used to reprehend Servius Galba most weightily — a man whom he claimed to remember well — because Galba, when L. Scribonius was bringing a question before him, had stirred the people's pity, although M. Cato, Galba's heavy and keen enemy, had spoken harshly and vehemently before the Roman people — which speech he himself set out in his \textit{Origines}.

\ciceroSection{228} Rutilius therefore reprehended Galba, because he had himself almost lifted onto his own shoulders the orphan son of his kinsman C. Sulpicius Gallus, who, by the recollection and memory of his most distinguished father, was to move tears in the people; and he had commended his own two small sons to the people's guardianship, and had said, as if making his will in battle-line without scales and tablets, that he was making the Roman people their guardian for their bereavement. And so, when Galba was hard pressed by both unpopularity and the people's hatred, he said he had been freed by these tragic shows. And I likewise see this written in Cato — that, if he had not used the boys and tears, Galba would have paid the penalty. These things Rutilius vehemently censured, and said that to that lowliness either exile or death was to be preferred.

\ciceroSection{229} Nor did he say this only — but he himself both felt and did it. For when he was, as you know, that very example of innocence, when no one was either of more upright character or more holy in the state, he was not only unwilling to be a suppliant before the judges, but even refused that his case be pleaded more ornately or freely than the simple rule of truth bore. He gave a small share to this Cotta here, a most eloquent young man, his sister's son. The case was also pleaded in some part by Q. Mucius, in his customary way, with no display, purely and clearly.

\ciceroSection{230} If you, Crassus, had then spoken — you, who said a little before that the orator's reinforcement for the abundance of speaking should be sought from those disputations the philosophers use; and if for Rutilius you had been allowed to speak, not after the philosophers' fashion but your own — however wicked those men had been (as they were pestilent citizens worthy of punishment), still the force of your speech would have torn out from their inmost minds all their unforgiveness. Now, such a man was lost, while the case was being pleaded as if the matter were being conducted in that fictitious state of Plato's. No one of the patrons groaned, no one cried out, nothing pained anyone, no one complained, no one called on the commonwealth, no one supplicated; what more? No one stamped his foot in that trial — to keep, I suppose, his own report from the Stoics.

\ciceroSection{231} The man — a Roman and a consular — imitated that ancient Socrates, who, although he was the wisest of all men and had lived most holily, so spoke for himself in the trial for his life that he seemed not a suppliant or a defendant but the master and lord of his judges. Indeed, when the most eloquent orator Lysias had brought him a written speech which, if it pleased him, he might learn by heart and use for himself in the trial — he read it, not unwillingly, and said it had been suitably written. ``But,'' he said, ``as if you had brought me Sicyonian shoes, I should not use them, however well-fitted to the foot, because they were not manly,'' so that speech seemed to him eloquent and oratorical, but not strong and manly. So he too was condemned — not only by the first verdicts, in which the judges only decided whether to condemn or acquit, but also by those which they had to give again according to the laws.

\ciceroSection{232} For at Athens, when the defendant had been condemned, if the offence was not capital, there was, as it were, an estimation of the penalty; and when the verdict was being given by the judges, the defendant was asked what kind of estimation he most confessed himself to have deserved. When Socrates was asked this, he answered that he had deserved to be honoured with the most ample distinctions and rewards, and that his daily nourishment should be furnished him publicly in the Prytaneum — which honour is held the highest among the Greeks.

\ciceroSection{233} At his answer the judges so flared up that they condemned a most innocent man to death. If he had been acquitted — what, by Hercules, even though it does not concern us, I should still wish for, on account of the magnitude of his talent — how could we bear those philosophers, who now say, when he has been condemned for no other fault save by reason of ignorance of speaking, that nevertheless precepts of speaking ought to be sought from them? With whom I do not contend whether their thing be better or truer: I only say that what they teach is one thing and what we are talking of is another, and the latter can without the former be supreme.

\ciceroSection{234} As for the civil law, Crassus, which you embraced so vehemently, I see what you have done. Even when you were speaking, I saw it. First you gave yourself over to Scaevola, whom we all are most justly bound to love for his outstanding sweetness; whose art, when you saw it was without dowry and unadorned, you enriched and adorned with the dowry of words. Then because you had spent more pains and labour on it, since the encourager and master of that pursuit was at home with you, you feared that, unless you had heaped up that art by your speech, you should have lost your trouble.

\ciceroSection{235} But I do not contend even with that art. Let it be as great as you wish — for, without dispute, both is it great and lies wide open and reaches to many things, and has always been in the highest honour, and the most distinguished citizens preside over its study even today — but see, Crassus, lest while you wish to adorn the science of civil law with new and foreign ornament, you also strip and bare it of its own conceded and handed-down adornment.

\ciceroSection{236} For, if you said that he who was a jurisconsult was also an orator, and likewise he who was an orator was also a jurisconsult, you would establish two distinguished arts and equal to one another and partners of the same dignity. But now you confess that the jurisconsult can exist without this eloquence we are seeking, and that very many such men have existed; but you deny that the orator can exist unless he has taken on that science. So in your view the jurisconsult by himself is nothing except a kind of cautious and shrewd pettifogger, a herald of actions, a chanter of formulas, a snarer of syllables; but because the orator often uses the help of law in cases, you have therefore appended this science of law as a little handmaid and follower to eloquence.

\ciceroSection{237} As to your wonder at the shamelessness of those advocates who either, when they did not know small things, professed great ones, or who dared to handle in cases the greatest things in civil law when they did not know them and had never learned them — to either matter the defence is easy and ready. For nor is it surprising that one who does not know in what words a \textit{coemptio} is made could still defend the case of the woman who has made the \textit{coemptio}; and if the science of steering is the same in a small ship as in a great, it is not therefore so that one who does not know the words by which property must be partitioned could not plead the case of the dividing of an estate.

\ciceroSection{238} As to your bringing forward those greatest cases of the Hundred Men as set in law — what one of them at last was there which by an eloquent man unskilled in law could not have been most splendidly pleaded? In all of which cases — as in M.' Curius's very case which you lately pleaded, and in the dispute of Gaius Hostilius Mancinus, and in the case of that boy who had been born of the second wife when no notice of divorce had been sent to the first — there was the highest disagreement on law among the most expert men. I ask, then, what aided the orator in those cases the science of law, when the jurisconsult was about to depart from court the better — but supported, not by his own art, but by another's, that is, not by the science of law but by eloquence.

\ciceroSection{239} I have often heard this: when Publius Crassus was canvassing for the aedileship and the older and now consular Ser. Galba was attending him on his rounds, because he had betrothed Crassus's daughter to his own son Gaius, a certain rustic came up to Crassus to consult him. He drew Crassus aside, and laid the matter before him, and got from him an answer truer than convenient to his own affair. When Galba saw the man dejected, he called him by name and asked on what matter he had laid the case before Crassus. When he had heard the answer, and saw the man

\ciceroSection{240} disturbed, ``Crassus,'' he said, ``has answered you, I see, with mind preoccupied and busy.'' Then he took Crassus himself by the hand and said: ``Hey there — what came into your mind to answer so?'' Then Crassus, a most expert man, confidently maintained that the matter so stood as he had answered, and that there could be no doubt. But Galba, playing on the matter variously and copiously, brought forward many likenesses, and said many things for equity against the strict law. And the other (since he could not be his match in argument — although Crassus was in the number of the eloquent, by no means a match for Galba), took refuge in the authorities, and brought forth what he himself was saying as written both in his brother Publius Mucius's books and in Sextus Aelius's commentaries — and yet conceded that Galba's argument seemed to him probable and almost true.

\ciceroSection{241} And yet, the cases which are of such a kind that no doubt about their law can exist, are not generally brought to judgment at all. Does anyone seek an inheritance by that will which the head of a household has made before the son was born to him? No one. For it is agreed that the will is broken by his being born. Therefore in this kind of law there are no judgments. The orator may then with impunity be ignorant of all this part of the law not in dispute, which without doubt is by far the greatest;

\ciceroSection{242} and on the law which is uncertain among the most expert, it is not difficult for the orator on whichever side he shall defend, to find some authority. From him, when he has received as it were thonged spears, he will himself hurl them by the orator's arms and powers. Unless indeed — let me speak with the leave of this excellent man — you defended the case of M.' Curius from Scaevola's little books or from your father-in-law's precepts; you did not, but you snatched up the patronage of equity and the defence of wills and of the wishes of the dead.

\ciceroSection{243} And in my judgment — I have often heard you and been at hand — by far the greater part of the votes you charmed by your salt and your charm and your most polished wit; when you derided that excessive subtlety, and marvelled at Scaevola's talent — that he had thought out that one ought to be born before one died; and when you collected many things from laws and from senatorial decrees and from the life and conversation of every day, not only sharply but also wittily and humourously, where, if we should follow the words and not the matter, nothing could be brought to a head. So that trial was full of merriment and joy. In which what your practice of civil law profited you, I do not understand; an outstanding force of speaking, joined with the highest gaiety and grace, profited you.

\ciceroSection{244} That very Mucius, the defender of his ancestral law and as it were the champion of his own patrimony, what did he bring in that case, when he was speaking against you, that seemed drawn out of civil law? What law did he recite? What did he reveal by speaking that had been more obscure to the unlearned? Surely his whole speech turned on this — that the written word ought to be of the highest weight. But on this kind of thing all boys are exercised by their masters, when in such cases they are taught now to defend the written letter, now equity.

\ciceroSection{245} And, I suppose, in that case of the soldier, if you had defended either the heir or the soldier, you would have had recourse to Hostilius's actions, not to your own force and oratorical faculty. Nay, even if you defended the will, you would so plead that all the law of all wills should seem set in that judgment; or if you pleaded the soldier's case, you would, as is your way, by speaking have raised his father from the dead. You would have set him before our eyes; he would have embraced his son and, weeping, would have commended him to the Hundred Men. By Hercules, you would have forced all the stones to weep and lament — so that all that ``\textit{as the tongue shall have pronounced}'' should seem written not in the Twelve Tables, which you set above all libraries, but in some schoolboy's exercise.

\ciceroSection{246} As to your accusing the laziness of young men, who do not first learn that easiest of all arts — how easy it may be, those will see who, on the strength of that art's arrogance as if it were the hardest, walk about so propped up; and finally you yourself will see, who say that the art is easy which you concede is not yet wholly an art at all but, sometime, if anyone shall have learned another art so that he may make this an art, then it will be an art. Then, that it is full of delight: in which all men remit you that delight to yourself and bear to be without it. Nor is there any one of them who, if anything must be learned by him, would not rather memorise the \textit{Teucer} of Pacuvius than Manilius's laws on the sale of slaves.

\ciceroSection{247} Then in that, by love of country, you think we ought to know our ancestors' inventions — do you not see that the old laws either have grown old by their very age or have been removed by new laws? And as to your supposing that good men are made by civil law, because the laws set out rewards for the virtues and punishments for the vices: I should have thought virtue was handed down to men, if it can be handed down by reason at all, by teaching and persuading, not by threats and force and fear. For that very thing — how good it is to take care against evil — we can know without any knowledge of law.

\ciceroSection{248} As for me — to whom alone you concede that I can satisfy in cases without any knowledge of law — I answer you, Crassus, that I never have learned the civil law, and yet in those cases which I could defend in the law-court, I have never wished for that knowledge. For it is one thing to be the master of a particular kind and art, and another in common life and in the ordinary practice of men to be neither dull nor unschooled.

\ciceroSection{249} Which of us is at liberty to go round our own farms or to visit our country property whether for profit or pleasure? And yet no one lives so without eyes, so without sense, that he is wholly ignorant of what sowing and harvest are, what the pruning of trees and vines is, at what time of year and in what way these things are done. Must, then, if any farm has to be inspected or anything ordered to a steward about the cultivation of land or commanded to a bailiff, the books of Mago of Carthage be thoroughly mastered? Or can we be content with this common understanding? Why, therefore, in the civil law also — especially when we are worn down in cases and business and the Forum — can we not be sufficiently equipped at least so that we should not seem strangers and aliens in our own country?

\ciceroSection{250} If now some more obscure case has been brought to us, it would be hard, I suppose, to share it with this Scaevola — although the people whose business it is bring all questions to us already considered and worked out. But if the dispute is about the matter itself, about boundaries when we have not gone to the spot, about tablets and registers, we necessarily learn intricate and often difficult matters; if laws or the answers of expert men have to be known to us, are we afraid that, if from youth we have less devoted ourselves to civil law, we shall not be able to learn them? Therefore is the science of civil law in nothing useful to the orator? I cannot deny that any kind of knowledge is useful, especially to him whose eloquence ought to be adorned with abundance of subjects. But the things necessary to the orator are many and great and difficult — so much so that I should not wish to draw off his industry into too many studies.

\ciceroSection{251} Who would deny that for the orator there is need, in this oratorical motion and bearing, of Roscius's gesture and grace? Yet no one would advise young men eager for speaking to labour in learning gesture in the way of actors. What is so necessary to the orator as voice? Yet no one eager for speaking, on my counsel, will be a slave to the voice in the way of the Greek tragic actors, who both for many years sit declaiming and daily, before they speak, lying down they gradually call up the voice and the same, when they have spoken, sitting they bring back from the highest tone down to the lowest tone, and as it were in a certain way collect it. If we would do this, those whose cases we have undertaken would be condemned before we had recited the Paean or hymn as many times as is prescribed.

\ciceroSection{252} If, then, in gesture, which much aids the orator, and in voice, which alone most either commends or supports eloquence, we are not at liberty to labour, and we can attain in each only what the time of our daily duty allows in this battle-line — how much less must we descend into the occupation of mastering civil law? — which both can be acquired summarily without instruction, and has this unlikeness with those other things: that voice and gesture cannot be suddenly taken on or snatched from elsewhere, while the usefulness of law for any case can on the spot be drawn either from the experts or from the books.

\ciceroSection{253} Therefore those most eloquent men have, as helpers in their cases, men skilled in the law, although they themselves are most unskilled — those who, as you said a little before, are called \textit{pragmatici}. In which our people altogether do far better, in that by the authority of the most distinguished men they have wished the laws and rights to be guarded. But still, the Greeks would not have failed to grasp this, had they thought it so necessary, that they should have trained the orator himself in the civil law, not given him a \textit{pragmaticus} for a helper.

\ciceroSection{254} As to your saying that old age is freed from solitude by the science of civil law — perhaps even by a great quantity of money. But we are inquiring not what is useful to us, but what is necessary to the orator. Although, since we take many things for the orator's likeness from one craftsman, the same Roscius is wont to say that, the more years are added to him, the slower will be his pipe-player's measures and the easier his songs. If he, bound by certain measure of rhythm and feet, contrives something for the rest of old age, how much more easily can we, not loosen our measures, but change them altogether?

\ciceroSection{255} Nor does this escape you, Crassus, how many and various are the kinds of speaking — a thing which I scarcely know whether you yourself first showed, who for a long time have spoken much more remissly and gently than you used to. And no less is this gentleness of your most weighty speech approved than that highest force and contention. Many were the orators — as we hear of that famous Scipio and Laelius — who carried through everything in conversation, only a little more intent, and never strove (as Servius Galba did) with their lungs and shouting. If you can no longer do this, or do not wish to, do you fear that your house — that of such a man and such a citizen — if not frequented by litigious men, will be deserted by all the rest? For my part I am so far from that opinion, that I not only do not think the support of old age must be set on the multitude of those who come to consult, but I look forward to that solitude you fear, as to some haven. For I judge that the best support of old age is leisure.

\ciceroSection{256} As for the rest, even if they help — I mean history, and prudence in public law, and the memory of antiquity, and an abundance of examples — when at any time I have need, I shall borrow them from a most excellent man and most equipped in those very things, my friend Congus, and I shall not refuse — what you have just exhorted — that they should read everything, hear everything, busy themselves in every honourable study and in culture. But, by Hercules, they do not seem to me to have so much spare time, if they shall wish to do and pursue what has been laid down by you, Crassus. You seemed to me to lay laws on this age almost too hard, but yet for the attaining of what they desire, almost necessary.

\ciceroSection{257} For both the sudden exercises in proposed cases, and the careful and meditated reflections, and that pen of yours which you have truly said is the maker and master of speaking — that takes much sweat. And that comparison of one's own speech with others' writings, and the sudden disputation, on another's writing, for the sake of praising or blaming or approving or refuting, is no slight contention either for the memory or for imitation.

\ciceroSection{258} But there was one thing horrible — which, by Hercules, I fear has had more force to deter than to exhort: for you wished each one of us in his kind to be, as it were, a Roscius. And you said that what is right is not so much approved as what is wrong sticks to the fastidious; which I think is observed not so fastidiously in us as in actors.

\ciceroSection{259} For so I see we are often listened to most attentively even when we are hoarse — for the very matter and case holds them; while Aesopus, if he croaks even a little, is hissed off the stage. From those of whom nothing is sought except pleasure of the ears, offence is taken as soon as anything of pleasure is diminished. But in the eloquent man there are many things which hold; and if all of these are not of the highest quality but the most are still great, those very things which are present must necessarily seem wonderful.

\ciceroSection{260} Therefore, to come back to my first point: let the orator with us be, as Crassus has described him, the man who can speak with adaptation to persuading. But let him be confined to those things which lie in the common practice of states and the Forum, and, with the rest of his pursuits laid aside, however ample and brilliant they may be, let him in this one work, so to say, be pressed night and day. And let him imitate that Athenian Demosthenes, to whom without doubt the highest force of speech is conceded — in whom there is said to have been so great a zeal and so great a labour that, first of all, he overcame by his diligence and industry the impediments of nature; and although he was so stammering that he could not say the first letter of the very art he was pursuing, he so brought it about by reflection that no one was thought to have spoken more plainly. Then, when his breath was rather short, he so achieved by holding his breath in speaking, that — as his writings show — in one continuous stretch of words he kept up two risings and fallings of the voice;

\ciceroSection{261} who also, as memory has handed down, having put pebbles into his mouth, used to declaim many verses at the highest pitch in one breath; nor did he stand on one spot, but walked about and went forward up a steep ascent.

\ciceroSection{262} Crassus, with these exhortations to rouse young men to study and labour I heartily agree. The rest, which you have collected from various and diverse studies and arts — although you yourself have attained them all — I think to be set apart from the orator's proper office and duty.'' When Antonius had said this, Sulpicius and Cotta seemed indeed to be in doubt as to which speech seemed to come closer to the truth.

\ciceroSection{263} Then Crassus: ``You are making the orator into a kind of journeyman for us, Antonius — and I am not sure but that you feel otherwise, and are using that wonderful habit of yours of refutation, in which no one has ever surpassed you. The exercise of which faculty itself is proper to orators, but is also now in vogue with the philosophers, especially with those who on every subject set forth are wont to speak most copiously on either side.

\ciceroSection{264} For my part, I thought — especially with these as our hearers — that I ought to picture what kind of man he could be, who lived on the benches and brought nothing more than the necessity of cases demanded. But I saw something greater, when I judged that the orator, especially in our state, ought to be without none of the ornaments. You, however, since you have closed the orator's whole duty within certain narrow boundaries, will the more easily set out for us those things which have been asked of you about the duties and precepts of the orator. But, I think, after this day; for enough has been said by us today.

\ciceroSection{265} Now both Scaevola, since he has decided to go to his Tusculan villa, will rest a little until the heat breaks; and we ourselves, since the time is what it is, will give attention to our health.'' This pleased everyone. Then Scaevola: ``I should much wish,'' he said, ``that I had not arranged with L. Aelius that I should come to his Tusculan villa today; gladly would I hear Antonius.'' And as he was rising up, smiling at the same time: ``Nor was he so troublesome to me,'' he said, ``in that he plucked at our civil law, as he was pleasing in confessing that he himself did not know it.''

\ciceroBook{2}

\ciceroSection{1} Quintus, my brother, if you remember, when we were boys we held the firm opinion that Lucius Crassus had touched no more of learning than the first schoolboy training had been able to give him; and that Marcus Antonius had been wholly without and ignorant of all education. Many there were who, although they did not so judge, were yet glad — that they might more easily deter us from learning when we were burning with zeal — to say what I have said about those orators. So that, if uneducated men had attained the highest prudence and incredible eloquence, all our labour might seem empty, and the studious zeal of our father — that excellent and most prudent man — for our education might seem foolish.

\ciceroSection{2} Whom we, as boys, were wont to refute by witnesses from our own household: by my father and by C. Aculeo, our kinsman, and by my uncle L. Cicero. For my father and Aculeo (whose wife was our maternal aunt, and whom Crassus loved best of all) and my uncle (who had set out with Antonius to Cilicia and had returned with him) often told us many things about Crassus's zeal and learning. And when we were learning, with our cousins, Aculeo's sons, those things which pleased Crassus, and were being trained by those teachers whom he employed, we often understood — when we were at his house, what even boys could feel — that he both spoke Greek so as to seem to know no other tongue, and set our teachers in inquiry such questions, and himself in every conversation handled such matters, that nothing seemed new to him, nothing unheard.

\ciceroSection{3} About Antonius, although we had often heard from our most cultivated uncle how the man at Athens or at Rhodes had given himself to the conversations of the most learned men, still, even as a young lad, as far as the modesty of that early age of mine permitted, I often asked him many things. To you what I am writing here will not be new; for already even then you would hear from me that he, from his many and various conversations, did not seem to me to be raw or ignorant in any matter at all which lay in those arts about which I could form a judgment.

\ciceroSection{4} But there was this in each of them: that Crassus wished to be thought, not so much not to have learned, as rather to despise those things and to set Roman prudence in every kind ahead of Greek; while Antonius judged his own speech would be more credited by the people if he were thought never to have learned at all. And so each of them thought he would be the weightier, if the one seemed to despise the Greeks, the other not even to know them.

\ciceroSection{5} What their counsel was is nothing to the matter at present. But this belongs to the subject and the time of this composition: that no one ever could either have flourished in eloquence or stood out in it, not only without the doctrine of speaking, but not even without all wisdom. For the rest of the arts almost defend themselves each by themselves; but to speak well — which is to speak knowingly, expertly, ornately — has no defined territory whose boundaries hedge it in. All things, whatever can fall into the discussion of men, must be well spoken by him who professes himself able to do this — or the name of eloquence must be given up.

\ciceroSection{6} Therefore I confess that, both in our own state and in Greece itself (which has always considered these things the highest), there have been many men of excellent talent and of great praise of speaking, without the highest knowledge of all things. But that such an eloquence stood forth as was in Crassus and Antonius, without all the things known which pertained to so great a prudence, and so great an abundance of speaking as was in them — this I affirm could not have been.

\ciceroSection{7} I have done this all the more willingly: that I have committed to writing the conversation which they once held between themselves on these matters — both that the opinion which had always existed might be removed, that the one was not most learned, the other plainly unlearned; both that I might keep in writing those things which I judged to have been said divinely by the foremost orators about eloquence, if in any way I could attain and embrace them; or by Hercules even that I might rescue their praise, now well-nigh growing old, as far as I could, from men's forgetfulness and silence.

\ciceroSection{8} For if they could have been known from their own writings, perhaps I should have judged this to be a less laborious task for me. But since the one had left not much that survived, and that itself written when he was a young man, and the other had left almost nothing of writing at all, I thought it owed by me to such great talents that, while we were still holding their living memory, I should make it immortal if I could.

\ciceroSection{9} Which work I approach with the greater hope of carrying it through to approval — because I am not writing about the eloquence of Servius Galba or of Gaius Carbo, in which I might be allowed to invent if I wished, with no man's memory now refuting me. Rather, I am setting these things forth to be known by those who often heard the very men of whom I speak; so that, by the witness of the memory of those living and present in whom both of those orators are known, I may commend two supreme men to those who have seen neither.

\ciceroSection{10} Nor in fact am I now pursuing you, my dearest and best brother, to teach you with certain rhetorical books which you call rustic; for what could be more subtle or more ornate than your own speech? But whether by judgment, as you are wont to say, or as that father of eloquence, Isocrates, wrote of himself, you have shrunk back from speaking by modesty and a kind of native timidity; or whether (as I myself am wont to joke) you have thought one rhetorician sufficient not only in one family but in nearly the whole state — I do not, however, suppose that these books will be of you of that kind which can deservedly be mocked because of the meagreness of the good arts in those who have argued about the rationale of speaking. For nothing in Crassus's and Antonius's discussion seems to me to be passed over which anyone with the highest talent, the keenest zeal, the best learning, the greatest experience, could think capable of being known and grasped — which you can most easily judge, who have wished to grasp the prudence and rationale of speaking by yourself, and the practice through us.

\ciceroSection{11} But that we may carry through this no slight task we have undertaken the more quickly, with our own exhortation laid aside, let us come to the conversation and discussion of those whom we have set before us.

\ciceroSection{12} On the day after these things had been done, then, about the second hour, when Crassus was still in bed and Sulpicius sat by him, and Antonius was walking with Cotta in the colonnade, suddenly Q. Catulus the elder came up with his brother C. Julius. When Crassus heard this, he rose up startled, and all wondered, suspecting that there must be some weighty cause for their coming.

\ciceroSection{13} When they had greeted one another in the friendliest way, as their accustomed intimacy bore, ``What of it now?'' said Crassus. ``Surely there is no news?'' ``Nothing at all,'' said Catulus, ``for you see there are the games. But — even if you think us either inept or troublesome — when Caesar came to me yesterday evening at my Tusculan villa from his own Tusculan villa, he told me Scaevola had been met by him on his way from here, and from him he had heard wonderful things: that you, whom I had never been able by every kind of trying to draw out for disputation, had argued with Antonius about eloquence at length, and as it were in a school, almost after the Greek fashion.

\ciceroSection{14} So my brother prevailed on me — myself not very averse to listening, but, by Hercules, fearing lest we should intrude troublesomely upon you — to come here with him. For Scaevola said a good part of the conversation had been put off till today. If you think this too eager, you will assign it to Caesar; if too familiar, to both of us. We — unless we have intruded troublesomely — are glad to have come.''

\ciceroSection{15} Then Crassus: ``I should be glad, whatever cause had brought you here, when I see at my house men dearest and most friendly to me; but yet, to speak truly, I should rather it had been any cause than this you mention. For — to speak as I feel — never have I been less pleased with myself than yesterday. This befell me more by my own ease of nature than by any other fault: for, while I humoured the young men, I forgot that I was an old man, and did what not even as a young man I had done — to dispute about subjects which lay in some learning. But yet this has fallen out very opportunely for me: that, my own part being now done, you have come to hear Antonius.''

\ciceroSection{16} Then Caesar: ``For my part, Crassus,'' he said, ``I am so eager to hear you in that longer and continuous disputation that, if I cannot have it, I shall be content with this everyday conversation of yours. I shall try then so that my friend Sulpicius and Cotta should not seem to have more weight with you than I do; and I shall surely beg you to share some part of your sweetness with me also, and with Catulus. But if it pleases you the less, I shall not press you, nor bring it about that, while you fear to be foolish, you make me think I have been so.''

\ciceroSection{17} Then he: ``By Hercules, Caesar, of all the Latin words I have always thought the force of this word \textit{ineptus} the very greatest. He whom we call \textit{ineptus} seems to me to have got his name from this — that he is not \textit{aptus} (fitting); and that has very wide currency in our common conversation. For one who either does not see what the time demands, or talks too much, or shows himself off, or has no regard for the dignity or the convenience of those with whom he is, or finally is, in some respect, ungainly or excessive — that man is called \textit{ineptus}.

\ciceroSection{18} With this fault that most learned Greek nation is heaped up; and so, because the Greeks do not see the force of this evil, they have not even imposed a name on the fault. For though you should hunt all through, you will not find any way the Greeks call \textit{ineptus}. Of all imperitnences, which are countless, I do not know whether any is greater than (as those Greeks are wont) in any place, among any men one likes, to dispute most subtly about the most difficult or unnecessary subjects. This we, against our will and protesting, were forced to do yesterday by these young men.''

\ciceroSection{19} Then Catulus: ``Not even the Greeks, Crassus, who have been distinguished and great in their cities — as you are, and as we all wish to be in our own commonwealth — were like these Greeks who force themselves on our ears; nor did they shun this kind of conversation and disputation in their leisure.

\ciceroSection{20} And if those who have no regard for the time, the place, the persons, seem to you (as they ought) inept — does at last either this place seem unsuitable, in which this very colonnade where we are now walking, and the wrestling-ground, and so many seats of gymnasia stir somehow the memory of the Greeks' disputations? Or is the time inopportune, in so much leisure, which is rarely given and now has been given to us most opportunely? Or are the persons foreign to this kind of disputation, who all of us are such that we judge there is no life without these pursuits?''

\ciceroSection{21} ``All of those things,'' said Crassus, ``I interpret in another way. First, I think the wrestling-ground and seats and colonnades the Greeks themselves invented for exercise and delight, not for disputation. For the gymnasia were invented many ages before the philosophers began to chatter in them; and at this very time, when the philosophers occupy all the gymnasia, still those who hear them prefer to hear the discus rather than the philosopher. As soon as it has clanged, in the very middle of his speech they all leave the philosopher disputing about the greatest and weightiest things, to anoint themselves; so they prefer the lightest pleasure to the weightiest usefulness, as they themselves call it.

\ciceroSection{22} As to the leisure you speak of, I agree. But the fruit of leisure is not the straining of the mind, but its relaxation. I have often heard from my father-in-law that he said his own father-in-law Laelius used always to retire into the country with Scipio, and that they two were wont to grow young again incredibly when they had flown out from the city to the country as out of chains. I do not dare to say it of such men, but Scaevola is wont to say that they used to gather shells and pebbles at Caieta and Laurentum, and to come down to every relaxation of mind and play.

\ciceroSection{23} For so the matter is: as we see birds for the sake of breeding and their own usefulness fashion and build nests, and the same, when they have completed something, free themselves from their work and fly about here and there in the open to lighten their toil — so our minds, wearied by forensic business and city work, leap up and long to fly free from care and labour.

\ciceroSection{24} And so what I said in Curius's case to Scaevola I did not say otherwise than I felt. For ``If, Scaevola,'' said I, ``no will shall be rightly made except what you have written, all the citizens will come to you with our tablets, and you alone will write all the wills. What then?'' I said. ``When will you do public business? When the friends'? When your own? When, in short, will you do nothing?'' Then I added: ``A man does not seem to me free who does not, sometimes, do nothing.'' In which opinion I stand fast, Catulus; and when I have come here, this very thing — to do nothing and to be plainly idle — delights me.

\ciceroSection{25} As to your third point — that you are men who would think life unsweet without these pursuits — that not only does not exhort me to dispute but even deters me. For as Gaius Lucilius, a learned and most urbane man, was wont to say, what he wrote he wished read neither by the most unlearned nor by the most learned, because the one would understand nothing, the other perhaps more than he himself; about whom he even wrote: ``I do not care to be read by Persius'' (this man, as we knew, was nearly the most learned of all our men) ``but I want Laelius Decumus,'' whom we knew, a good man and not unlettered, but nothing compared with Persius. So I, if I now had to dispute about these our pursuits, should not wish to do so among rustics, but much less among you. For I should rather not be understood than be reproved.''

\ciceroSection{26} Then Caesar: ``For my part, Catulus,'' he said, ``I now seem to myself to have wasted no labour in coming here. For this very refusal of disputation has been a kind of disputation, very pleasant to me. But why do we hinder Antonius, whose part I hear it is to argue about the whole of eloquence, and whom Cotta and Sulpicius have been awaiting for some time?''

\ciceroSection{27} ``For my part,'' said Crassus, ``I shall neither permit Antonius to say a word, and shall myself fall silent, unless I first obtain something from you'' — ``What?'' said Catulus. ``That you stay here today.'' Then, when the other was hesitating because he had promised his brother, ``For both of us,'' said Iulius, ``I answer: so we shall do; and on this very condition you would hold me, even though you should not say a word.''

\ciceroSection{28} Here Catulus smiled and at the same time, ``My doubt,'' he said, ``has been cut off, since I had given no orders at home and the man whose guest I was to be has so easily promised without my opinion.'' Then they all turned their eyes upon Antonius. And he: ``Listen, then, listen,'' he said. ``You will be hearing a man trained in the school and from a master and in Greek letters, and I shall speak the more confidently because Catulus has come up as a hearer, to whom not only do we Latin speakers but the Greeks themselves are wont to concede the subtlety and elegance of his own language.

\ciceroSection{29} But because nevertheless this whole matter, whatever it is — whether art or pursuit of speaking — without the addition of cheek, can be nothing, I shall teach you, my disciples, that which I myself have not learned: what I feel about every kind of speaking.''

\ciceroSection{30} Here, after they had smiled: ``The matter,'' he said, ``seems to me distinguished by faculty, mediocre in art. For art is of those things which are known; but the orator's whole action is contained in opinions, not in knowledge. For we both speak before those who know nothing, and we say what we ourselves do not know. Therefore both they at different times feel and judge differently about the same subjects, and we often plead opposite cases — not only that Crassus speaks against me sometimes or I against Crassus (when one of us must necessarily say what is false), but even that each of us at one time defends one thing on the same subject, at another the opposite, when there cannot be more than one truth. So in such a matter — which rests on falsehood, which does not often arrive at knowledge, which hunts the opinions of men and often their errors — I shall speak, if you think there is reason for you to listen.''

\ciceroSection{31} ``We do think so, and very strongly,'' said Catulus, ``and the more so because you do not seem to me about to use any ostentation. For you have begun, not as you say grandly, but rather from truth than from any kind of dignity.''

\ciceroSection{32} ``As, therefore,'' said Antonius, ``I have confessed about the kind itself that the art is not very great, so I affirm this: that certain very keen precepts can be given, both for handling the minds of men and for catching their wills. The science of this thing, if anyone wishes to call a great art, I shall not contend; for since most men plead cases in the Forum rashly and on no plan, while some — by exercise or by some habit — do this more skilfully, there is no doubt that, if anyone has noticed what it is by reason of which some speak better than others, he can mark it. Therefore he who has done this in the whole sphere will have invented, if not plainly an art, yet a kind of art.

\ciceroSection{33} And would that, as I seem to myself to see those things in the Forum and in cases, so I could now hunt out for you how they are discovered! But I shall see to my own part. Now I lay this down — what I am persuaded of — that even if it be no art, still nothing is more brilliant than the perfect orator. For — to omit the practice of speaking, which dominates in every peaceful and free state — there is so great a delight in the very faculty of speaking that nothing can be perceived more pleasant either to men's ears or to their minds.

\ciceroSection{34} For what song can be found sweeter than measured speech? What poem more apt than the artful close of words? What actor more pleasant in imitating truth than the orator in undertaking it? What is more subtle than thoughts close-set and keen? What more wonderful than a matter set off by the splendour of words? What more full than a speech crammed with every kind of subject? Nor is there any matter not proper to the orator, that ought indeed to be spoken with ornament and weight.

\ciceroSection{35} It is his to give counsel on the greatest matters with an opinion set out with dignity; his is both the rousing of a flagging people and the moderating of an unbridled one. By the same faculty are the frauds of men called to ruin, and integrity to safety. Who can urge to virtue more ardently? Who recall from vices more keenly? Who rebuke the wicked more harshly? Who praise the good more handsomely? Who break desire more vehemently by accusing? Who lighten grief more gently by consoling?

\ciceroSection{36} History indeed — the witness of times, the light of truth, the life of memory, the teacher of life, the messenger of antiquity — by what other voice than the orator's is it commended to immortality? For if there is some other art which professes the science of making or selecting words; or if anyone but the orator is said to shape the speech and to vary it and to mark it with, as it were, certain badges of words and thoughts; or if any way is handed down save from this one art, of arguments or of thoughts or, in short, of distribution and order, let us confess either that what this art professes is alien to it, or that it is shared with some other art. But if in this one art alone is that rationale and instruction, even if some men of other arts have spoken well, that is not the less proper to this one.

\ciceroSection{37} But just as the orator about those subjects which belong to other arts, if only he has come to know them, can speak excellently (as Crassus said yesterday), so the men of the other arts speak those things of theirs more handsomely if they have learned anything from this art.

\ciceroSection{38} For if any farmer about country matters or, what is common, a doctor about diseases, or a painter about painting, has spoken or written eloquently, eloquence is not on that account to be reckoned of his art. In which, since there is great force in the talents of men, many even without instruction attain something in every kind and art. But what is each thing's own — even though it can be judged from this when you have seen what each art teaches — yet nothing can be more certain than this: that all the other arts can perform their function without eloquence, but the orator without it cannot keep his name. So that the others, if they shall be eloquent, have something from this art, while this man, unless he has equipped himself with the resources of his own household, cannot seek the abundance of speaking from elsewhere.''

\ciceroSection{39} Then Catulus: ``Although, Antonius, this rush of your speech is by no means to be hindered by interruption, you will yet bear with me, and pardon me. `For I cannot but cry out,' as the man says in the \textit{Trinummus}: so subtly have you seemed to express the force of the orator and so copiously to praise it. The eloquent man indeed ought to do this best of all — to praise eloquence; for he ought to bring to bear in praising it that very thing which he praises. But go on; I agree with you that this whole — to speak eloquently — is yours; and that if anyone in another art does it, he is using a good thing taken from elsewhere, not proper or his own.''

\ciceroSection{40} And Crassus: ``The night, Antonius,'' he said, ``has polished you and given you back to us as a man. For yesterday, in your conversation, you had described the orator to us as some kind of man of one trade only — as Caecilius says, a kind of rower or porter, a man poor of culture and uncivil.'' Then Antonius: ``Yesterday I had set this before me — that, if I had refuted you, I might draw away these pupils from you. Now, with Catulus and Caesar listening, I think I ought not so much to fight with you as to say what I myself feel.

\ciceroSection{41} It follows, therefore, since this man we are speaking of must be set up by us in the Forum and before the eyes of the citizens, that we see what business we shall give him, and over what office we shall wish him set. Crassus yesterday — when you, Catulus and Caesar, were not present — laid down briefly in the division of the art the same that most Greeks have laid down (and showed not what he himself felt, but what they say): there are two first kinds of questions in which eloquence moves, the one unbounded, the other fixed.

\ciceroSection{42} The unbounded I took him to mean: in which something is asked in general, like this — `Should eloquence be sought after? Should offices?' The fixed: in which something is asked in particular persons and in a settled and defined matter — of the kind which arise in the Forum and in the cases and disputes of citizens.

\ciceroSection{43} These seem to me set either in pleading a suit or in giving counsel. For that third kind, which both Crassus touched on and (as I hear) the famous Aristotle himself, who most of all set these matters in light, joined to the others — even if there is room for it, is yet less necessary.'' ``What kind?'' said Catulus. ``Is it the laudations? For I see that as a third kind set down.''

\ciceroSection{44} ``Yes,'' said Antonius, ``and in that kind I know both I myself and all who were present were vehemently delighted, when Popilia your mother was praised by you — for she was the first woman, I think, to whom this honour was paid in our state. But all things which we say should not, in my view, be referred to art and precepts;

\ciceroSection{45} for from those springs whence the precepts for every ornament of speaking are drawn, it will be possible also to adorn a laudation, and there will be no need of those elementary precepts. For if no one handed them down, who is there who would not know what should be praised in a man? With these things laid down — those which Crassus said in the opening of his speech which he made as censor against his colleague: that whatever nature or fortune gave to men, in those things he could allow himself to be surpassed with equanimity; but those things which men could prepare for themselves, in those things he could not allow himself to be surpassed — he who shall praise anyone will understand that he must set out fortune's goods.

\ciceroSection{46} These are: family, money, kinsmen, friends, resources, health, beauty, strength, talent, and other things which are either of the body or external. If he had them, he made good use of them; if he had them not, he was wisely without them; if he lost them, he bore the loss moderately. Then how wisely he, whom he praises, has done or borne anything, how generously, how bravely, how justly, how splendidly, how piously, how gratefully, how kindly — what, finally, with any virtue: these and things of this kind he who wishes to praise will easily see, and he who wishes to blame, the contraries.''

\ciceroSection{47} ``Why then do you doubt,'' said Catulus, ``to make this a third kind, since it lies in the rationale of things? For not, if it is easier, must it on that account be excluded from the count.'' ``Because I do not wish,'' he said, ``to so handle all the things that fall to the orator from time to time, however small they are, as if nothing could be said without its own precepts.

\ciceroSection{48} For testimony must often be given, and sometimes even with care — as it was even necessary for me to do against Sextus Titius, a seditious and turbulent citizen. In giving that testimony I unfolded all the counsels of my consulship by which I had withstood that tribune of the plebs for the commonwealth, and I set out what I judged he had done against the commonwealth. I was held a long time, I heard much, I answered much. Does it then please us, when you give precepts about eloquence, to hand down something also about giving testimonies, as if it were part of the art?''

\ciceroSection{49} ``Not at all,'' said Catulus. ``What if — what often happens to the foremost men — instructions must be set forth, either in the senate by a general, or to a general, or to a king, or to some people from the senate; because we must use a more careful kind of speech in such cases, must this part of cases also be reckoned in, or be furnished with its own precepts?'' ``By no means,'' said Catulus. ``For an eloquent man will not lack the faculty for such matters, gathered from the rest of subjects and cases.''

\ciceroSection{50} ``Therefore,'' said Antonius, ``those things, which often have to be eloquently performed, and which I said a little before, when I was praising eloquence, are the orator's: they neither have any place of their own in the division of parts, nor a fixed kind of precepts, and yet must be performed no less eloquently than what is said in a suit — chiding, exhortation, consolation, none of which is without need of the highest ornaments of speech; but those things require no precepts from the art.''

\ciceroSection{51} ``I plainly agree,'' said Catulus. ``Come now,'' said Antonius, ``of what kind of orator and of how great a man in speaking do you reckon to write history?'' ``If as the Greeks have written it, of the highest,'' said Catulus; ``if as our own, no orator is needed: it is enough not to be a liar.'' ``And yet — that you may not despise our own — the Greeks themselves at the start wrote so, like our Cato, like Pictor, like Piso.

\ciceroSection{52} For history was nothing else but the making up of annals: for the keeping of the matter and of the public memory of it, from the beginning of Roman affairs until Publius Mucius the Pontifex Maximus, the Pontifex Maximus committed all the affairs of each year to writing and entered them on a tablet and set the tablet up in his house, that there might be the power of knowing for the people; and these things even now are called the \textit{annales maximi}.

\ciceroSection{53} This way of writing was followed by many, who without any ornaments left only records of times, of men, of places, and of things done. So such as among the Greeks Pherecydes, Hellanicus, Acusilas were, and very many others, such are our Cato and Pictor and Piso, who do not hold the things by which speech is adorned (for those things have only just been imported here), and, provided what they say is understood, think the one praise of writing to be brevity.

\ciceroSection{54} A little above himself raised, and added a louder note of history, that excellent man, Crassus's friend, Antipater. The rest were not adorners of subjects but only narrators.'' ``It is,'' said Catulus, ``as you say. But that very Caelius neither distinguished history by variety of colours, nor by the arrangement of words and a smooth and even drawing-out of speech polished that work; but as a man neither learned nor most fitted to speaking, hewed it as he could. Yet he surpassed, as you say, his predecessors.''

\ciceroSection{55} ``It is hardly to be wondered at,'' said Antonius, ``that this matter has not yet been made bright in our tongue. For no one of our men is zealous for eloquence except to shine forth in cases and the Forum; while among the Greeks the most eloquent men, withdrawn from forensic causes, applied themselves both to the rest of distinguished subjects and most of all to the writing of history. For we are told that Herodotus, who first adorned this kind, did not move at all in cases. And yet his eloquence is so great that — as far as I can understand Greek writings — he greatly delights me. After him, in my judgment, Thucydides easily surpassed all in the art of speech.

\ciceroSection{56} He is so dense in the abundance of subjects that he matches the number of words to the number of thoughts; in words so apt and compressed that you do not know whether the matters are illustrated by the speech, or the words by the thoughts. And yet not even him, although he had a part in public life, do we have on record as one of those who pleaded cases. He is said to have written these very books at the time when he had been removed from public affairs and (as has often happened to the best men at Athens) had been driven into exile.

\ciceroSection{57} Him followed Philistus of Syracuse, who, although he was the closest friend of the tyrant Dionysius, spent his leisure in writing history, and most of all imitated, as it seems to me, Thucydides. Then afterwards from the most distinguished workshop, as it were, of the rhetorician — two men of outstanding talent, Theopompus and Ephorus, urged on by Isocrates as their master, betook themselves to history. Cases they did not touch at all.

\ciceroSection{58} At last even from philosophy: that famous Xenophon, the Socratic, set out, and afterwards Callisthenes, the companion of Alexander, sent on by Aristotle, wrote history; and the latter indeed almost in the rhetorical manner. The former used a certain gentler note, and not having that impulse of the orator — perhaps less vehement, but yet, as it seems to me, somewhat sweeter. The youngest of all of these, Timaeus, was, as far as I can judge, by far the most learned, both in copiousness of subjects and most rich in variety of thoughts, and not unpolished in the very arrangement of words; he brought great eloquence to writing — but no forensic experience.''

\ciceroSection{59} When he had said this, ``What is this, Catulus?'' said Caesar; ``where are those who say Antonius does not know Greek? How many historians he has named! How learnedly, how aptly, he has spoken about each one!'' ``By Hercules, while wondering at this,'' said Catulus, ``I now leave off wondering at what I had been more wondering at: that this man, when he was not knowing these things, could so much in speaking.'' ``And yet, Catulus,'' said Antonius, ``I am not hunting for some usefulness in speaking when I am wont to read those books and not a few others, but only for the sake of pleasure, when there is leisure.

\ciceroSection{60} What then? There is, I will confess, some good in it nevertheless. Just as when I walk in the sun, even if I walk for some other reason, it happens by nature that I am tanned, so when I have read those books — at Misenum, for at Rome it is hardly possible — more eagerly, I feel that my speech is, as it were, coloured by their touch. But that this may not seem too widely true for you, I understand only those things in Greek which they themselves who wrote wished to be commonly understood.

\ciceroSection{61} If at any time I happen on your philosophers — deceived by the indices of their books which are inscribed mostly on subjects known and brilliant: about virtue, about justice, about honour, about pleasure — I understand absolutely no word: so are they bound up in narrow and clipped disputations. As for the poets — speaking, as it were, in some other tongue — I do not even attempt to touch them. With those, as I said, I delight myself who have written either histories or their own speeches, or who speak in such a way that they seem to have wished to be intimate with us, who are not the most learned.

\ciceroSection{62} But to come back: do you see what a great office of the orator history is? I am inclined to think it the greatest both for the river of speech and for variety; nor do I find it set out anywhere separately by the precepts of the rhetoricians. The principles are set before our eyes. For who does not know that the first law of history is that he should not dare to say anything false? Then, that he should dare nothing of truth not to say? That there should be no suspicion of favour in his writing, none of enmity?

\ciceroSection{63} These foundations are, of course, known to all; but the very building rests on the matter and on the words: the rationale of the matter requires the order of times, the description of regions; it requires also — since in great matters worthy of memory the counsels are first looked for, then the deeds, afterwards the outcomes — both that of counsels it should be signified what the writer approves, and in things done it should be made clear not only what was done or said, but also in what manner; and when the outcome is being told, that all the causes should be unfolded — whether of chance, or of wisdom, or of recklessness — and not only the deeds of the men themselves, but also of those who excel by report and name, the life and nature of each.

\ciceroSection{64} The plan of words and the kind of speech, easy and drawn out and flowing evenly with a certain gentleness, must be pursued without this judicial harshness and without the stings of forensic thoughts. Of these so many and so great matters, do you see that there are no precepts which are found in the arts of the rhetoricians? In the same silence have lain many other duties of orators: exhortations, precepts, consolations, admonitions, all of which must be handled most eloquently, but have no place of their own in those arts which have been handed down.

\ciceroSection{65} And in this kind there is also that boundless wood: that to the orator most have given two kinds of speaking — as Crassus also showed — one of a fixed and definite case, such as those which lie in suits, in deliberations (let one add, if he wishes, also laudations); the other, which all writers nearly call but no one explains, an unbounded question of a kind, without time and without person. What this is and how great, when they speak of it, they do not seem to me to understand.

\ciceroSection{66} For if it belongs to the orator, on whatever subject is laid down without bounds, to be able to speak, he will have to speak of how great is the size of the sun, what is the shape of the earth; on mathematics and on music he will not be able to refuse, having undertaken this burden. Finally, to him who professes that it is his to speak not only on those disputes which are marked by times and persons, that is, all forensic ones, but also on the unbounded questions of kinds, no kind of speech can be excepted.

\ciceroSection{67} But if we wish also to add that part of questions to the orator — vagrant and free and far reaching — that he must speak on things good or evil, things to be sought or shunned, honourable or base, useful or useless; on virtue, justice, self-control, prudence, greatness of soul, generosity, piety, friendship, duty, faith, the rest of the virtues and contrary vices; and likewise on the commonwealth, empire, military matters, the discipline of the state, the morals of men — let us take that part also, but in such a way that it is bounded by moderate regions.

\ciceroSection{68} For my part I think that everything which pertains to the practice of citizens, the morals of men, that lies in the habit of life, in the rationale of the commonwealth, in this civil society, in the common feeling of man, in nature, in morals, must be grasped by the orator: if not so as to give back about these subjects, in the philosophers' way, separately, then surely so as to be able prudently to weave them into a case. And about these very subjects, that he should speak as those who instituted laws, statutes, states have spoken — simply and splendidly, without any train of disputations and without a meagre wrestling of words.

\ciceroSection{69} Lest at this point there be any wonder, if no precepts are laid down by me on so many and so great matters, I lay down this: as in the other arts, when the most difficult things of each art have been handed down, the rest, because they are either easier or like, need not be handed down. As in painting, the man who has fully learned to paint one form of a man can paint any form or age, even if he has not been taught it; nor is there danger that one who paints a lion or a bull excellently cannot do the same in many other quadrupeds; nor is there any art at all in which all that can be made by it is handed down by the master. Rather, those who have learned the very kinds of the first and fixed things, will attain the rest passably by themselves.

\ciceroSection{70} Similarly, I think in this — call it rationale or practice of speaking — that he who has gained that force, that he can move at his pleasure the minds of those who hear, with some power of decision, either about the commonwealth or about his own affairs or about those against or for whom he is speaking — that man will not have to inquire of the whole class of the rest of speeches what he should say, more than that famous Polyclitus, when he was modelling Hercules, asked himself how to model the lion-skin or the hydra, even if he had never separately learned how to do these.''

\ciceroSection{71} Then Catulus: ``You seem to me, Antonius,'' he said, ``to have set splendidly before our eyes what one ought to learn who is to be an orator, and what — even if he has not learned it — he will get from what he has learned. For you have brought the whole man down into two kinds of cases only, and have left the countless rest to practice and analogy. But see that, in those two kinds, the hydra and the skin be yours, while Hercules and other greater works be left in the things you pass over. For it does not seem to me a less work to speak about the universal kinds of things than about cases of single persons, and much greater to speak about the nature of the gods than about men's quarrels.'' ``It is not so,'' said Antonius;

\ciceroSection{72} ``for I shall tell you, Catulus, not so much from learning as (what is greater) from experience: the speech about all other matters, believe me, is play to a man not dull, not unpractised, not without share of common letters and the more polished culture. But in the contention of cases there is a great work, and I do not know but that, of all human works, by far the greatest. In which the orator's force is mostly judged by the unskilled by the outcome and the victory; where there is an armed adversary present who must be both struck and beaten back; where often the man who is to be lord of the matter is unfamiliar and angry, or even a friend of the adversary and an enemy to you; when he must be either taught or untaught, or restrained or stirred up, or in every way moderated by the speech to the time, to the case (in which often goodwill must be led to hatred, and hatred to goodwill); or, as if by some engine, must be twisted now to severity, now to relaxation of mind; now to gloom, now to merriment.

\ciceroSection{73} The weight of every thought must be used, the burdens of every word; there must come in besides a delivery various, vehement, full of mind, full of breath, full of grief, full of truth. In these works, if anyone has grasped that art, so as to be able like Phidias to make Minerva's image, he will not labour, in any case, how to learn — as the same artist did on the shield — to make the smaller works.''

\ciceroSection{74} Then Catulus: ``The greater and more wonderful you have made these things, the greater the expectation that holds me of by what reasonings, by what precepts that great force is got. Not that it now matters to me — for my age does not require it; we have followed another kind of speaking, who have never wrested the votes from the judges' hands by any force of speech, but, with their minds appeased, have got only what they themselves were prepared to grant — but yet your matters I ask, drawn on by zeal of knowing, not for any practical use of mine.

\ciceroSection{75} Nor have I need of some Greek teacher to chant me common precepts, when he himself never looked upon the Forum, never any judgment. Like that Peripatetic Phormio of whom this is told: when Hannibal, expelled from Carthage, had come to Antiochus at Ephesus as an exile and, because his name was held in great glory among all, was invited by his hosts to hear, if he wished, the man I have spoken of; when he had said he was not unwilling, the copious man is said to have spoken several hours about the duty of a general and about all military matters. Then, when the rest who had heard him were vehemently delighted, they asked Hannibal what he himself thought about that philosopher. The Phoenician, in not the best Greek but yet freely, is reported to have answered that he had often seen many doting old men, but had seen no one more doting than Phormio.

\ciceroSection{76} Nor by Hercules unjustly. For what could be either more arrogant or more talkative than that, to Hannibal, who for so many years had contended with the Roman people, the conqueror of all peoples, for empire, a Greek who had never seen an enemy, never a camp, who had never even touched the smallest part of any public office, should give precepts about military matters? This all those men seem to me to do who give precepts about the art of speaking. For what they themselves have not experienced, they teach others. But in this they perhaps err the less: that they do not, like that man with Hannibal, try to teach you, but boys or young lads.''

\ciceroSection{77} ``You are mistaken, Catulus,'' said Antonius; ``for I myself have fallen in with many Phormios already. For who is there of those Greeks who thinks any of us understands anything? And to me they are not so troublesome: I bear all of them easily and patiently. For they either bring something which does not displease me, or they make me regret less that I have not learned. I dismiss them, however, not so insultingly as that famous Hannibal dismissed the philosopher; and so I have, perhaps, more bother. But still their teaching, as far as I can judge, is utterly ridiculous.

\ciceroSection{78} For they divide the whole subject into two parts: into the dispute of the case and into that of the question. They call \textit{cause} the matter set in the wrangling and dispute of the parties, and \textit{question} the matter set in unbounded doubt. About the cause they give precepts; about the other part of speaking there is wonderful silence.

\ciceroSection{79} Then they make as it were five members of eloquence: to find what to say, to arrange what is found, then to adorn with words, afterwards to commit to memory, then at the end to deliver and pronounce; a matter not, indeed, recondite. For who would not see of his own accord that no one can speak unless he both has what to say and in what words and in what order to say it, and remembers it? And these things I do not blame, but I say they are set before our eyes — as those four, five, or six (or even seven, since they are arranged differently by different men) parts, into which by these men all speech is distributed.

\ciceroSection{80} For they bid us begin so as to make him who hears us favourable to us, and teachable, and attentive; then to narrate the matter, and so that the narration be plausible, open, brief; afterwards to divide the case or set it forth; to confirm our own with arguments and reasonings; then to refute the contrary; then some place the conclusion of the speech and as it were the peroration; others bid us, before we make the peroration, to make a digression for the sake of adornment or amplification, then to conclude and peroration.

\ciceroSection{81} I do not blame these things either; for they are neatly distributed. But still — what was bound to happen with men without experience of the truth — they have done it not skilfully. For what they would have to be precepts of openings and narrations are to be observed in whole speeches.

\ciceroSection{82} For I can more easily make a judge favourable to me when I am in the course of my speech than when everything is unheard; and teachable, not when I am promising I shall demonstrate, but when I am teaching and explaining; attentive, in fact, by often rousing the minds of the judges throughout the action — not by the first announcement.

\ciceroSection{83} Indeed, when they bid that the narration be plausible and open and brief, they admonish us rightly. But that they think these belong more to the narration than to the whole speech, in this they seem to me greatly to err. And the whole error is in this: that they suppose this art is some kind of thing not unlike the rest, of the kind which yesterday Crassus was saying could be set out about civil law itself: that the kinds of things should be set out first (in which it is a fault if any kind is omitted); then the parts of each kind (in which both for any part to be lacking and for any to be redundant is faulty); then the definitions of all the words, in which neither is anything fittingly missing nor in excess.

\ciceroSection{84} But if in civil law, if also in small or middling matters, the more learned can attain this, I do not feel the same can be done in this so great and so immense a matter. If anyone thinks otherwise, he must be brought to those who teach these things; he will find them all unfolded and polished. For there are countless books about these subjects, neither hidden nor obscure. But let them see what they want — whether it is for play or for fight that they are about to take up arms; for one thing the battle and the line, and another our play and the field, demand. And yet the very playing art of arms is of some use both to the gladiator and to the soldier. But a keen and ready and sharp and yet versatile mind makes men unconquered no less easily, even with art joined to it.

\ciceroSection{85} Therefore I shall train an orator for you, if I can, in such a way that I first see what he can do. Let him be touched with letters; let him have heard something, read; let him have got those very precepts. I shall test what becomes him; what he can accomplish with voice, strength, breath, tongue. If I understand he can attain to the highest, not only shall I exhort him to labour but shall even — if he seems to me also a good man — beseech him: so much do I set in an outstanding orator and a good man as well, of ornament for the whole state. But if it appear that, even when he has done everything to the highest, he is yet to come only to middling orators, I shall leave it to himself what he wishes; I shall not greatly trouble him. If he plainly recoils and is absurd, I shall warn him to keep himself in or to turn to another study.

\ciceroSection{86} For neither he who can do best is in any way to be deserted by our exhortation, nor he who can do something to be deterred. Of which the first seems to me a kind of divinity; the second — either not to do what you cannot do best, or to do what you do not do worst — humanity; but the third — to shout against what becomes you and what you can do — is, as you, Catulus, said of a certain shouter, the part of a man, by his own household crier, gathering as many witnesses as possible to his own foolishness.

\ciceroSection{87} Of him therefore who shall be such that he must be exhorted and helped, let us so speak as to hand to him only what experience has taught us, that with us as guides he may come to that whither without a guide we ourselves came — since we cannot teach better.

\ciceroSection{88} And to begin from a friend of mine, this Sulpicius, Catulus, when he was a young man, I first heard in a small case: with voice and form and movement of body and the rest of the things apt for this office of which we are inquiring; with a swift and animated speech, which was of his talent, and with words boiling over and a little overflowing, which was of his age. I did not despise him: for I wish fertility to bring itself out in a young man. For more easily, as in vines, are those things called back which have spread themselves too much, than, if the stock is feeble, are new shoots roused by cultivation. Likewise I wish there to be in the young man something for me to prune. For sap cannot be lasting in him who has run too quickly into ripeness.

\ciceroSection{89} I saw at once his disposition, and I did not let the time go, and I exhorted him to think of the Forum as his school for learning, but to choose what master he wished. If he should listen to me, then Lucius Crassus. He snatched up this advice and confirmed that he would do so, and even added, by way of grace, that I too should be his master. Hardly a year had gone by since this exhortation of mine when he prosecuted Gaius Norbanus, with me defending. It is not credible what difference it seemed to me there was between him as he then was and as he had been a year before. Wholly into that splendid and brilliant kind of Crassus's nature itself led him; but it could not have made progress enough unless he had bent himself with the same zeal and imitation, and so accustomed himself to speak as to fix his eyes on Crassus with all his mind and soul.

\ciceroSection{90} Therefore let this be the first thing in my precepts: that we point out whom one should imitate, and so as to follow most diligently those things which most stand out in him whom he shall imitate. Then let practice come on, by which he should mould and express the man he has chosen by imitating — not as I have often known many imitators do, who pursue either those things which are easy or even those which are striking and almost faulty.

\ciceroSection{91} Nothing is easier than to imitate someone's dress or attitude or motion; and if there is something even faulty, to take that and to be like in that fault is no great thing. As that Fufius, who, even now, his voice gone, rages in public — he does not attain the sinews in speaking of Gaius Fimbria (which Fimbria did have); he imitates the crookedness of mouth and the broadness of words. Fufius did not even know how to choose whom most he should be like, and in him whom he had chosen even wished to imitate the faults.

\ciceroSection{92} The man who shall do this as he ought, must first be wakeful in choosing; then in him whom he has approved, what most stands out, that he must most diligently follow. For what reason do you suppose there is that ages have brought forth, each its own kind of speaking? Which we cannot so easily judge in our own orators, since they have left no great quantity of writings from which judgment could be made, as in the Greeks, from whose writings we can understand what method and inclination of speaking belonged to each age.

\ciceroSection{93} The most ancient, of whose writings, that is, anything is established, are Pericles and Alcibiades, and of the same age Thucydides — subtle, keen, brief, abounding in thoughts more than in words. It could not have been by chance that the kind of all was one, unless they were setting some one before themselves to imitate. After them came Critias, Theramenes, Lysias. Many of Lysias's writings are extant; some of Critias's; about Theramenes we hear. All of these still kept that sap of Pericles, but were of a fuller fibre.

\ciceroSection{94} Lo, Isocrates rises up, the master of all those, from whose school as from a Trojan horse came forth pure princes; but of these some wished to be illustrious in the procession, others in the line of battle. And those — the Theopompi, Ephori, Philisti, Naucratae, and many others — differ in their natures, but are alike in their inclination both with one another and with their master. And those who betook themselves to cases — Demosthenes, Hyperides, Lycurgus, Aeschines, Dinarchus, and many others — although they were not equal to one another, all yet moved in the same kind of imitation of truth. As long as the imitation of these men remained, so long did that kind of speaking and zeal live;

\ciceroSection{95} but after these were extinguished and all memory of them gradually obscured and faded, certain other softer and more relaxed kinds of speaking flourished. Then Demochares, whom they say was the son of Demosthenes's sister; then that famous Demetrius of Phalerum, in my view of all those the most polished, and others like these stood forth. If we should wish to follow these down to this time, we should understand: as today all Asia imitates that Menecles of Alabanda and his brother Hierocles, whom I myself heard, so there has always been someone whom most wished to be like.

\ciceroSection{96} Therefore he who would attain this likeness by imitation, let him both pursue it with frequent and great exercises, and especially with writing. If our Sulpicius did this, his speech would be much closer-pressed: in which now there is sometimes — as the country folk say of plants in the highest fertility — a kind of luxuriance, which must be eaten down by the pen.''

\ciceroSection{97} Here Sulpicius: ``You advise me rightly,'' he said, ``and I am grateful for it; but, Antonius, I do not think that you yourself have written much.'' Then he: ``As if I should not give to others precepts of which I myself fall short. But I am thought not to keep accounts. Yet both in this, from my private affairs, and in that, from what I say (whatever it is), it can be judged what I do.

\ciceroSection{98} And yet we see there are many who imitate no one, and by their own nature attain what they wish without being like anyone. Which can be rightly noted in you, Caesar and Cotta; of whom the one has gained a kind of charm and salt unusual in our orators, the other a most keen and most subtle kind of speaking. Nor does your contemporary Gaius Curio — whose father, in my view, was perhaps the most eloquent man of those times — seem to me to imitate anyone greatly. Yet by the weight, elegance and abundance of his words he has expressed a certain form and figure of his own speaking; which I could most judge in that case which he pleaded against me before the Hundred Men on behalf of the Cossi brothers — in which nothing was wanting in him which the not only copious but also wise orator should have.

\ciceroSection{99} But that we may at last bring down to cases the man we have set up, and to those, in fact, in which there is somewhat more business — judgments and suits — someone perhaps will laugh at this precept, for it is not so much keen as necessary, and rather of an ordinary monitor than of a learned master. We shall first lay this down for him: that whatever cases he is to handle, he should learn them diligently and through and through.

\ciceroSection{100} This is not taught in school. For easy cases are brought to boys: ``A law forbids a foreigner to climb the wall: he climbed it; he repulsed an enemy; he is accused.'' There is no business in learning such a case. Rightly, then, no precepts on learning the case are given. For this is well-nigh the formula for cases in school. But indeed in the Forum tablets, testimonies, agreements, conventions, stipulations, kindreds, in-laws, decrees, opinions, and at last the whole life of those who are involved in the case must be learned. By neglect of which we see most cases — and most of all the private (which often are much more obscure) — lost.

\ciceroSection{101} So some, while they wish their own pains to be much esteemed, that they may be seen to flit through the whole Forum and to go from case to case, plead cases not learned. In which there is the great offence either of negligence in things undertaken, or of perfidy in things accepted; but a yet greater than they think, in that no one can speak about a matter he does not know without the greatest disgrace. So while they despise the reproach of laziness — which is the larger — they attain that which they themselves more flee, of slowness.

\ciceroSection{102} For my part I am wont to take pains that each man should himself teach me about his own affair, and that no one else should be present, that he may speak more freely; and to plead the adversary's case, that he may plead his own and bring out into the open whatever he has thought about his own affair. So when he has gone, I uphold three persons in one with the greatest equity of mind: my own, the adversary's, the judge's. Whatever ground is such as to have more help than harm, I judge ought to be spoken; where I find more bad than good, I reject it altogether and throw it out.

\ciceroSection{103} So I attain that I think at one time what I am to say, and at another say it: which two things most, trusting in their talent, do at once. But certainly those very same men would speak somewhat better, if they thought one time should be taken for thinking, and another for speaking.

\ciceroSection{104} When I have learned the matter and the case thoroughly, it occurs at once to my mind what the cause of disagreement is. For there is nothing on which men disagree — whether the case rest on charge (as of crime), or on dispute (as of an inheritance), or on deliberation (as of war), or on person (as on praise), or on disputation (as on the rationale of living) — in which there is not asked either what has been done or is being done or will be done; or of what kind it is; or what it is to be called.

\ciceroSection{105} And our cases — those which concern criminal charges — are mostly defended by denial. For both about cases of extortion (which are the greatest), almost everything must be denied; and in cases of bribery it is rarely granted that you can separate generosity and kindness from bribery and largesse; in cases of assassination, of poisoning, of embezzlement, denial is necessary. So this is the first kind of cases in courts: from a controversy of fact. In deliberations it is mostly from the future, rarely from the present or past.

\ciceroSection{106} Often again it is asked not whether the matter is or is not, but of what kind it is. As when, with me listening, the consul Gaius Carbo was defending the case of L. Opimius before the people, he denied nothing of Gaius Gracchus's killing, but said it had been done justly for the safety of the country. As to the same Carbo, when as tribune of the plebs he was managing public affairs in another mind, P. Africanus, when he was asking about Tiberius Gracchus, had answered that he seemed to him justly slain. And by ``justly'' all things are defended that are of such a kind that they appear either right, or lawful, or necessary, or done by accident or chance.

\ciceroSection{107} And what it is to be called is asked when there is contention with what word something is to be named. As I myself had with this Sulpicius the highest contention in Norbanus's case. For although I confessed many of those things which were charged by him, yet I denied that ``majesty'' had been impaired by him — on which word, by the Appuleian law, the whole case hung.

\ciceroSection{108} And in this kind of cases some prescribe that each side should briefly define that word which makes the case — which to me always seems most childish. For the definition of words is one thing when, among learned men, dispute is about the very subjects which lie in the arts (as when it is asked what art is, what law is, what state is, in which reason and instruction prescribes that the force of the thing you define should be so expressed that nothing be lacking nor in excess);

\ciceroSection{109} which neither Sulpicius did in that case nor I tried to do. For each of us, as much as he could, dilated with all the abundance of speaking what to ``impair majesty'' was. For a definition is often wrenched out of the hands by reproaching one word, or by adding or removing one; then by the very kind it smells of teaching and almost childish exercise. Then it can enter neither into the sense nor into the mind of the judge: for it slips by before it has been grasped.

\ciceroSection{110} But in that kind in which it is asked of what kind something is, contention often arises from the interpretation of writing — in which there can be no controversy save from the ambiguous. For the very thing which has been written when it disagrees with the intent has a kind of ambiguity, which is then unfolded when those words which are lacking have been supplied; with which added it is defended that the intent of the writing was clear. From contrary written texts, if anything is in dispute, no new kind arises but the case of the previous kind is doubled. And it will either never be able to be decided, or will be so decided that, by referring to the omitted words, whatever writing we shall defend may be supplied. So it happens that in those cases which are disputed because of writing one kind is left, if anything has been written ambiguously.

\ciceroSection{111} Of ambiguities, however, since there are several kinds, which seem to me better known to those who are called dialecticians (whom these our men are ignorant of, and ought to know no less): yet that is most frequent in every habit either of speech or writing — that something is ambiguous because either a word or words have been omitted.

\ciceroSection{112} Again they err when they separate this kind of cases, which lies in the interpretation of writing, from those cases in which it is disputed of what kind a thing is. For nowhere is it so much asked of what kind the very thing is as in writing, which is wholly separated from a controversy of fact.

\ciceroSection{113} So there are altogether three kinds which can fall into dispute and controversy: what has been done, is being done, or will be done; or of what kind it is; or how it is to be named. For the kind which some Greeks add — ``whether it has been done rightly'' — is wholly contained in ``of what kind it is.'' But now I shall come back to my subject.

\ciceroSection{114} When, therefore, the case has been received and the kind known, and I begin to handle the matter, I lay down nothing first save what that is to which all the speech which is proper to the question and judgment must be referred. Then I consider most diligently those two things, of which one carries the recommendation of ourselves or of those whom we are defending; the other is adapted to the moving of the minds of those before whom we speak, towards what we wish.

\ciceroSection{115} So all the rationale of speaking rests upon three things for the persuading: that we should prove what we maintain to be true; that we should win those who hear us to ourselves; that we should call the minds of those who hear us to whatever movement the case will require.

\ciceroSection{116} For the proving, the orator has a twofold subject-matter: the one of those things which are not invented by the orator, but lie in the matter and are handled by his rationale — such as tablets, testimonies, conventions and agreements, examinations, laws, decrees of the senate, things adjudicated, edicts, opinions, and the rest if there are any which are not discovered by the orator but brought to the orator from the case and from the matter; the other is what is wholly set in the disputation and argumentation of the orator.

\ciceroSection{117} So in the former kind we must think about the handling of arguments; in this we must think also about their inventing. Those who teach, when they have cut up cases into more kinds, supply each kind with a stock of arguments. Which, even if it is more apt for training the young (so that, as soon as the case has been laid down, they have where to refer themselves, whence they can at once draw out ready arguments), yet it is the part of dull talent to follow the little streams and not see the springs of things; and it now belongs to our age and practice to fetch what we wish from the source and see whence everything flows.

\ciceroSection{118} And the first kind of those things which are brought to the orator must be considered by us in advance, fitted for every use of similar matters. For both for tablets and against tablets, for witnesses and against witnesses, for examinations and against examinations, and likewise about the other things of the same kind, we are wont to speak either separately on the universal kind or, with definition, on the single times, men, cases. These topics — I say this to you, Cotta and Sulpicius — you must have prepared and ready by much reflection and meditation.

\ciceroSection{119} It is too long now for me to set out by what method one should establish or discredit witnesses, tablets, examinations. These all require even moderate talent, but the greatest exercise. Art and precepts they require thus far at most: that they should be adorned with certain lights of speaking.

\ciceroSection{120} Likewise those things of the second kind, which are wholly produced by the orator, do not have a difficult invention; they require rather a more brilliant and polished setting forth. So when these two things must be sought by us in cases — first what, then how we are to say — the latter, although it seems wholly stained with art and although it requires art, yet it is of moderate prudence to see what should be said. The other is that in which the orator's divine force and virtue is discerned: to say what must be said ornately, copiously, variously.

\ciceroSection{121} So that earlier part — since once you have so wished — I shall not refuse to polish and finish (how far I shall succeed, you shall judge): from what topics speech may be drawn to those three things which alone avail to produce conviction, that minds may both be conciliated and taught and moved. For these are the three things. But how they may be set in light, the man is at hand who can teach all things, who first introduced this into our usage, who most increased it, who alone has perfected it.

\ciceroSection{122} For I, Catulus — I shall speak without fear of any suspicion of flattery — think no orator a little more distinguished, neither Greek nor Latin, whom our age has produced, that I have not heard often and diligently. So if there is anything in me — which I now seem to dare to hope, since you, men of this talent, give me so much chance for hearing — it comes from this: that no orator has ever pleaded with me listening, what has not deeply settled in my memory. So I, who am as much as I am for judging, having heard all orators, without any doubt do thus lay down and judge: that no man of all has had so many and so great ornaments of speaking as are in Crassus.

\ciceroSection{123} Therefore, if you too think this same, it will not be — as I think — an unequal partition, if when I have created, nourished, confirmed this orator whom I am now fashioning as I have begun, I should hand him over to Crassus to be clothed and adorned.''

\ciceroSection{124} Then Crassus: ``You go on, Antonius,'' he said, ``as you have begun. For it is the part neither of a good nor of a generous parent not to clothe and adorn him whom you have begotten and brought up — especially since you cannot deny that you are rich. For what ornament, what force, what spirit, what dignity was wanting to that orator, who, in summing up a case, did not hesitate to rouse a consular defendant and tear open his tunic and show to the judges the scars on the front of the old general? — who likewise, with this Sulpicius accusing, when he was defending a seditious and frenzied man, did not hesitate to ornate the seditions themselves and to show with the weightiest words that many of the people's onsets had often been not unjust, which no one could check; that many seditions had often been carried through for the commonwealth — as when the kings were driven out, as when the tribunician power had been established; that that famous sedition of Norbanus, born of the citizens' grief and from hatred for Caepio, who had lost an army, could neither have been put down nor was it brought together unjustly?

\ciceroSection{125} Could this so doubtful, so unheard of, so slippery, so new ground have been handled without some incredible force and faculty of speaking? What shall I say of the pity for Cn. Manlius, of that for Q. Rex? What of the others without number? In which not so much that thing all give to you — a kind of singular keenness — has shone forth, but those very things which now you wish to assign to me have always been outstanding and supreme in you.''

\ciceroSection{126} Then Catulus: ``In you both,'' he said, ``I most marvel at this: that, although you are most unlike each other in speaking, yet each of you so speaks that nothing seems either denied to him by nature or not granted by learning. Therefore, Crassus, you shall not deprive us of your sweetness, that, if anything has been omitted or left by Antonius, you should set it out; nor shall we, Antonius, judge that, if you have not said something, you could not have said it, but rather that you preferred it to be said by Crassus.''

\ciceroSection{127} Here Crassus: ``Why,'' he said, ``do you not, Antonius, drop those things you have proposed, which none of these here desire — from what topics those things are discovered which must be said in cases? They are spoken by you in a kind of new and brilliant way, but they are easier in the matter and well-worn in the precepts. Bring out for us those things from which you draw, which most often you handle and always divinely.''

\ciceroSection{128} ``I shall bring them out,'' he said, ``and that I may the more easily exact from you what I ask, I shall refuse you nothing of what you require. Of all my method in speaking, and of that very faculty which Crassus has just lifted to the heavens with words, there are three matters, as I said before: one of conciliating men, the second of teaching them, the third of stirring them up.

\ciceroSection{129} Of these three parts, the first requires gentleness of speech; the second, keenness; the third, force. For this is necessary: that he who is to adjudicate the case for us either is inclined by the bent of his will toward us, or is led by the arguments of the defence, or is forced by the moving of his soul. But since that part in which the explanation and defence of the matters themselves is set seems to contain all the doctrine of this kind, about it we shall first speak; and we shall say a few things — for there are few which we can seem already to have, treated by use and as it were noted in the mind.

\ciceroSection{130} And we shall gladly agree with you, who wisely advise, Crassus, that we leave the defences of single cases which masters are wont to hand down to boys, and that we open up those heads from which all discussion of every cause and speech is drawn. For nor, as often as a word must be written, must we hunt out the letters of that word in our thought; nor, as often as a case must be pleaded, must we be obliged to come back to that case's separate arguments. We must rather have certain topics which, as letters for writing a word, so they may at once spring up to set out the case.

\ciceroSection{131} But these topics can avail only that orator who has been versed in either the practice of things, which years at last bring; or in the hearing and reflection which by zeal and diligence has run ahead of years. For if you bring me ever so learned a man, ever so keen and sharp in thinking, ever so ready in pronouncing, if he is yet a stranger to the custom of the state, to the examples, the institutions, the morals and wishes of his fellow citizens, those topics from which arguments are drawn will not much profit him. I have need of a softened soil, like a field not once ploughed but turned and ploughed again, that it may be able to bring forth better and larger fruits. The softening is practice, hearing, reading, letters.

\ciceroSection{132} Let him first see the nature of the case, which is never hidden — whether it is asked has it been done, or of what kind it is, or what name it bears. When this has been seen, by a kind of natural prudence (and not by these subtle subtractions those men teach), there at once occurs what makes the case — that is, what, taken away, the controversy could not stand. Then what comes into judgment — which they bid us inquire as follows: ``Opimius killed Gracchus. What makes the case? That for the sake of the commonwealth, when by senatorial decree he had called men to arms. Take this away — there will be no case. But Decius denies that this was lawful against the laws. Therefore there will come into judgment whether it was lawful, by senatorial decree, for the saving of the commonwealth.'' These things are perspicuous and set in common prudence; but those things must be sought, what arguments — both by the prosecutor and by the defender — looking to what comes into judgment, must be brought.

\ciceroSection{133} And here this is to be seen, in which is the chief error of those masters to whom we send our children — not that this matters greatly to speaking, but that you may see what kind these men are who think themselves learned. It is dull and unpolished. For in dividing the kinds of speeches they set up two kinds of cases: one they call that in which the inquiry is, without persons and times, of the universal kind; the other, what is defined by fixed persons and times — being ignorant that all controversies are referred to the force and nature of the universal kind.

\ciceroSection{134} For in that very case of which I spoke before, nothing of Opimius's or of Decius's person matters to the orator's topics. For the question is unbounded, on the universal kind itself: whether one ought to be punished who, by senatorial decree, has killed a citizen for the saving of the country, when this was not lawful by the laws? Lastly, there is no case in which what comes into judgment is asked from the persons of the parties, and not from the universal doubting of the kinds themselves.

\ciceroSection{135} Indeed, even in those very cases where there is dispute about the fact (as did P. Decius take money against the laws?), the arguments both of the charges and of the defence must be brought back to the kind and to the universal nature: that he was extravagant, on luxury; that he was greedy of others' goods, on avarice; that he was seditious, on turbulent and bad citizens; that he is accused by many, on the kind of witnesses; and on the contrary, what shall be said for the defendant must necessarily, from time and man, be referred to the common totals of things and kinds.

\ciceroSection{136} And these things, perhaps to the man who does not quickly grasp in mind everything in the nature of things, may seem to be very many that come into judgment when there is inquiry about a fact. But still the multitude is one of charges and defences, not of unbounded topics.

\ciceroSection{137} As for those things which, when there is no dispute about the fact, are asked of what kind they are — if you reckon them by the parties, they are countless and obscure; if by the things themselves, very moderate and bright. For if we lay Mancinus's case in Mancinus alone, as often as a man whom the \textit{pater patratus} has handed over has not been received, so often will a new case arise. But if that controversy makes the case — whether to him whom the \textit{pater patratus} has handed over, if he has not been received, there is \textit{postliminium} — Mancinus's name does not matter to the art of speaking nor to the arguments of the defence;

\ciceroSection{138} and if the dignity or unworthiness of the man brings anything besides, it is outside the question, and yet that very speech must necessarily be referred to the disputation of the universal kind. These things I do not argue with the design of refuting learned men. Although those are to be reproved who, in defining the kind, describe these cases as set in persons and in times. For although times and persons run in, still it must be understood that the cases hang not from these but from the kind of question.

\ciceroSection{139} But this matters nothing to me; for we ought to have no contest with those men. Enough that it is understood that they did not attain even this, which they could have brought about even without this forensic exercise in so much leisure: to distinguish the kinds of things and to set them out a little more subtly.

\ciceroSection{140} But this, as I said, matters nothing to me. This matters to me, and even much more to you, our Cotta and Sulpicius: as those men's arts now stand, the multitude of cases is to be feared. For it is unbounded if it is set in persons: as many men, so many cases. But if they are referred to the universal questions of kinds, they are so moderate and few that all diligent and mindful and sober orators ought to have them run through in mind, and (so to speak) recited by heart. Unless perhaps you suppose that L. Crassus learned his case from M.' Curius, and on that account brought forth many things why, when no posthumous son was born, Curius nevertheless ought to be Coponius's heir.

\ciceroSection{141} Coponius's name or Curius's mattered nothing to the abundance of arguments or to the force and nature of the case. The whole question lay in the universal kind of the matter and business, not in the time and names. When the writing was, ``\textit{If a son is born to me, and dies before...,} etc., \textit{then let that man be my heir},'' should it be supposed that the man instituted heir, when no son was born, is heir? — a question of perpetual law and of universal kind requires not the names of men but the principle of speaking and the springs of arguments.

\ciceroSection{142} In which our jurisconsults too hinder us and deter us from learning. For I see in Cato's and Brutus's books that what each man answered to anyone about a question of law is recorded almost by name; I think — that we should suppose there was some cause for the consultation or doubt in the persons, not in the matter. So that, since there are countless men, broken by learning the law, we should give up the will to learn together with the hope of mastering it. But these things Crassus will sometime expound and set out for us, distributed by classes. For he, lest you do not know, Catulus, promised us yesterday that the civil law, which now is diffused and scattered, he would compress into fixed kinds and reduce to an easy art.''

\ciceroSection{143} ``And indeed,'' said Catulus, ``that is by no means difficult for Crassus, who has both learned what could be learned of law, and what those who taught him lacked, he himself will bring; so that he can both fitly distribute and ornately illustrate the things in law.'' ``Then,'' said Antonius, ``we shall learn those things from Crassus when he has betaken himself from the throng and the benches into leisure (as he plans) and to his throne.''

\ciceroSection{144} ``Often I have heard him say it,'' said Catulus, ``that he was now resolved to depart from trials and cases. But, as I am wont to say to him myself, he will not be allowed: for he will not endure that his help be often implored in vain by good men, nor will the state bear it with equanimity, which, if it lacks Lucius Crassus's voice, will think itself stripped of a certain ornament.'' ``By Hercules,'' said Antonius, ``if these things have been truly said by Catulus, you must live with me in the same mill, Crassus; and let us yield that yawning and dozing wisdom to the leisure of the Scaevolas and the rest of the blessed.''

\ciceroSection{145} Crassus smiled gently here and: ``Only weave through, Antonius,'' he said, ``what you have begun. Yet that yawning wisdom, as soon as I have fled to it, will set me free.'' ``Of this topic which I have just begun, this is the end,'' said Antonius. ``Since it is understood that not in the countless persons of men nor in the unbounded variety of times, but in the cases and natures of kinds, are set all things which come into doubt; and the kinds, moreover, are defined not only by number but by paucity, so that students of speaking should grasp the matter of speech belonging to each kind, distributed in all its topics, equipped and adorned — I mean in matters and in thoughts.

\ciceroSection{146} These by their own force will produce words — which always seem to me sufficiently adorned, if they are of such a kind that the matter itself seems to have produced them. And if you ask the truth, what at least seems to me — for I can affirm nothing else save my own opinion and judgment — we ought to bring this equipment of cases and of universal kinds into the Forum, not, as each matter is brought to us, then at last to scrutinize the topics from which we may draw arguments. These can have been pursued by all who have only moderately considered them, with zeal and use applied. Still, the mind must be referred to those heads and to those topics, which I have often called by name, from which all things invented for every speech are drawn. And this is the whole — whether of art or of observation or of habit — to know the regions within which you hunt out and trace through what you seek.

\ciceroSection{147} When you have hedged in the whole region by your thought, if only you have got practice of things, nothing will escape you, and everything that is in the matter will occur and fall in. So, since for finding in speaking there are three things — keenness, then reasoning, which we may call (if we wish) art, third diligence — I cannot but concede the first place to talent. Yet diligence rouses the very talent even out of slowness;

\ciceroSection{148} diligence, I say, which is of greatest worth in all things and especially in defending cases. This must be cultivated by us above all; this must always be brought to bear; there is nothing this does not attain. That the case be deeply known, as I said at first, is by diligence; that we listen attentively to the adversary, and that we catch not only his thoughts but all his words, that we look at all his looks (which mostly indicate the senses of the soul), is by diligence;

\ciceroSection{149} but to do this dissemblingly, lest he seem to be making any progress for himself, is prudence; then that the soul should turn over among those topics which I shall propose a little later, that it may slip deep into the case, that it may be intent in care and thought, is by diligence; that to these matters one should bring memory as a kind of light, voice, and

\ciceroSection{150} strength, is diligence. Between talent and diligence indeed there is very little place left for art. Art shows only where you should seek and where the thing lies that you wish to find; the rest is in care, in attention of mind, in thought, in vigilance, in assiduity, in labour. I shall embrace it in a single word, which I have often used: diligence. By this one virtue all the other virtues are contained.

\ciceroSection{151} For we see that the philosophers abound in plenty of speech — who, as I think (but you, Catulus, know this better), give no precepts of speaking, and not for that less do they undertake whatever subject is set down, on which they may copiously and abundantly speak.''

\ciceroSection{152} Then Catulus: ``It is, Antonius, as you say. Most philosophers hand down no precepts of speaking, and yet have ready what they should say on each subject. But Aristotle — whom I most admire — set forth certain topics, from which the whole way of every argument may be discovered, not only for the philosophers' disputation, but also for this speech which we use in cases. From which man, Antonius, your speech has long been not far astray — whether by likeness with that divine talent you light upon the same footsteps, or whether you have read and learned those very things (which to me does seem more probable). For I see that you have given more pains to Greek matters than we had supposed.''

\ciceroSection{153} Then he: ``You will hear the truth from me, Catulus. I have always thought the orator who first showed as little as possible of any artistry, and then no sign of Greek matters, would be more pleasing and credible to this people. And yet I likewise judged it the part of a beast, not of a man, when the Greeks were undertaking, professing, doing such great things, and were promising men that they would give them the rationale both of seeing into the most obscure subjects, and of living well, and of speaking copiously, not to lend an ear; and, if you did not dare openly to hear them, lest you diminish your authority among your fellow citizens, yet at least to catch their voices by listening in, and from afar to attend to what they were telling. So I have done, Catulus; and of all those men I have summarily tasted the cases and the kinds themselves.''

\ciceroSection{154} ``Most timidly, by Hercules,'' said Catulus, ``as if to some rock of lust you have brought your mind close to philosophy — which this state has never spurned. For Italy was once full of Pythagoreans, when in this country was that great Greece. Some report that Numa Pompilius, our king, was a Pythagorean — who lived very many years before Pythagoras himself; he therefore must be reckoned a greater man, since he came to know that wisdom for setting up the state nearly two ages before the Greeks felt that it was born. And surely this state has produced none more distinguished or weightier in authority or more polished in culture than Publius Africanus, Gaius Laelius, Lucius Furius — who openly always had with them the most learned men from Greece.

\ciceroSection{155} I have often heard from them that they said the Athenians had done a most pleasing thing both to themselves and to many leading men of the state, in that, when they were sending envoys to the senate about their greatest matters, they had sent the three most noble philosophers of that age — Carneades, Critolaus, and Diogenes — and that, while they were at Rome, those men were heard often by themselves and by others. With these as your authorities, Antonius, I wonder why you have all but declared war on philosophy, like that Zethus in Pacuvius.''

\ciceroSection{156} ``By no means,'' said Antonius. ``I have so resolved to philosophise rather, as Neoptolemus in Ennius — `\textit{a little: for, all in all, I do not like it}.' But this is my judgment, which I seemed to have set out: I do not condemn those studies, only let them be moderate. The opinion of those studies, and the suspicion of artistry before those who judge cases, I think is contrary to the orator: it diminishes both the orator's authority and the credit of his speech.

\ciceroSection{157} But that the speech may be recalled to the place from which it has slipped: of those three most distinguished philosophers, whom you said came to Rome, do you see that Diogenes was the man who said he was handing down the art of arguing well and judging the true and the false, which by the Greek word he called \textit{dialectic}? In this art, if it is an art, there is no precept of how the truth may be discovered, but only how it may be judged.

\ciceroSection{158} For everything we say, in such a way that we either say it is or is not — the dialecticians, even if it be said simply, take to judge whether it is true or false; and if it is uttered conjointly, and other things adjoined, they judge whether they are rightly adjoined and whether the sum of each reasoning is true. At the last they prick themselves with their own keenness, and by much asking find many things which not only they themselves can no longer untie, but by which the things even unwoven and rather laid open before are well-nigh unwoven again.

\ciceroSection{159} So this Stoic of yours helps us in nothing, since he does not teach how I may discover what to say. The same man even hinders us, in that both he finds many things which he says cannot in any way be untied, and he brings a kind of speech which is not clear, not flowing forth, but meagre, dry, clipped, and minute. If anyone shall approve it, he will so approve it that he yet confesses it not to be apt for the orator. For this our speech must be adapted to the ears of the multitude, for soothing minds, for impelling them, for proving those things which are not weighed by a goldsmith's scales, but on a kind of popular balance.

\ciceroSection{160} Therefore we wholly send away that art, which is too dumb in inventing arguments, too talkative in judging them. That Critolaus, whom you record came with Diogenes, I think could have helped this our pursuit more. For he was from that Aristotle from whose discoveries I seem to you to wander not far. And between this Aristotle (whose book I have read in which he set out the art of speaking of all his predecessors, and those in which he himself said certain things of his own about that same art) and these full-blooded masters of this art, this seems to me the difference: that he, with the same edge of mind by which he had seen the force and nature of all things, also looked upon those things which pertained to the art of speaking which he despised; while those, who held this alone should be cultivated, lived in the handling of this one rationale, not with the same prudence as he, but with greater practice in this one kind and study.

\ciceroSection{161} Carneades' incredible force of speaking, indeed, and his variety, would be very desirable to us — who in those his disputations defended no thing he did not prove, attacked nothing he did not overturn. But that is something greater than what should be demanded of those who hand down and teach these things.

\ciceroSection{162} For my part, if I should now wish to instruct in speaking some quite raw man, I should rather hand him over to those constant teachers who pound the same anvil with one labour day and night, and who put all the most slender particles and all the smallest morsels chewed up, like nurses to suckling boys, into his mouth. But if he should be one who seemed to me both liberally instructed in learning, and now imbued with some practice, and of sufficiently keen talent — to that place I should hurry him where he should be held by no separate little stream, but whence the universal river breaks forth; the man who should point out to him the seats and as it were homes of all arguments, and briefly illustrate them, and define them with words.

\ciceroSection{163} For what is there in which he should hesitate, who has seen that everything which is taken in a speech, either for proving or for refuting, is either taken from the matter itself by force and nature, or assumed from outside? From its own force, when it is asked either what the whole matter is, or its part, or what name it has, or anything which touches that matter; from outside, when those things are gathered which are outside and do not cling to the nature of the matter.

\ciceroSection{164} If the whole matter is asked, the universal force must be set out by definition, thus: ``If majesty is the greatness and dignity of the state, he diminishes it who has handed an army over to the Roman people's enemies, not he who has handed over to the Roman people's power him who had

\ciceroSection{165} done so.'' If the part — by partition, in this way: ``Either the senate had to be obeyed about the saving of the commonwealth, or another counsel taken, or it had to be done of one's own accord. Another counsel — proud; one's own — arrogant; we had therefore to use the senate's counsel.'' If from a name — as Carbo: ``If the consul is one who consults the country,

\ciceroSection{166} what else has Opimius done?'' But if from what touches the matter, there are more seats and topics of arguments. We shall seek both conjuncts and kinds and parts subjected to kinds, and likenesses and unlikenesses, and contraries and consequents and concomitants and as it were forerunners and reluctancies, and we shall trace the causes of things, and what has arisen from causes, and we shall seek greater, equal, less.

\ciceroSection{167} From conjuncts arguments are drawn thus: ``If the highest praise must be given to piety, you ought to be moved when you see Q. Metellus mourning so piously.'' From the kind: ``If the magistrates ought to be in the power of the Roman people, why do you accuse Norbanus, whose tribunate

\ciceroSection{168} obeyed the will of the state?'' From a part subjected to a kind: ``If all who consult for the commonwealth ought to be dear to us, surely the generals first of all, by whose counsels, valour, and dangers we keep both our safety and the dignity of our empire.'' From a likeness: ``If wild beasts love their offspring, with what fondness ought we to be

\ciceroSection{169} towards our own children?'' From an unlikeness: ``If it is the part of barbarians to live for the day, our own counsels ought to look to everlasting time.'' In each kind, both of likeness and of unlikeness, examples are to be drawn from the deeds, sayings, and outcomes of others, and fictitious narrations must often be set down. Now from a contrary:

\ciceroSection{170} ``If Gracchus did wickedly, Opimius brilliantly.'' From consequents: ``If he was killed by the sword, and you, his enemy, were caught in that very place with a bloody sword, and no one besides you was seen there, and to none was there a motive, and you were always reckless — what is there that we can doubt about the deed?'' From concomitants and forerunners and reluctancies — as Crassus once when a young man: ``Not, Carbo, if you defended Opimius, will those people on that account think you a good citizen. It is plain that you played a part and were aiming at something else: because you have often lamented Tiberius Gracchus's death in public meetings, because you were the partner of Publius Africanus's death, because you carried that law in your tribunate, because

\ciceroSection{171} you have always been in disagreement with the good men.'' From the causes of things, thus: ``If you wish to remove avarice, its mother must be removed: luxury.'' From those things which have arisen from causes: ``If we use the resources of the treasury both for the supports of war and for the ornaments of peace, let us serve the

\ciceroSection{172} taxes.'' Greater, smaller, and equal we shall compare thus: from the greater: ``If a good reputation surpasses riches and money is so eagerly sought after, how much more should glory be sought!'' From the smaller: ``This man bears the death of this man, of slight acquaintance, so familiarly: what if he himself had loved him? What will he do for me, his father?'' From the equal: ``It is the same man's task both to seize and against the commonwealth

\ciceroSection{173} to lavish money.'' From outside there are taken those things which are propped up not by their own force but by external: as: ``This is true, for Q. Lutatius said it. This is false, for the examination has been held. This must be followed: for I am reading the tablets.'' Of which whole kind I have spoken a little before.

\ciceroSection{174} These things, as briefly as could be, have been said by me. For just as if I should wish to point out gold to anyone which had been buried in many places, it should be enough to show the signs and marks of the places, knowing which he might dig for himself and find what he wished with little labour and no error — so I have given these marks of arguments, which point out to the seeker where they are; the rest is brought out by care and thought.

\ciceroSection{175} What kind of argument best suits each kind of cases — it is not for an exquisite art to prescribe, but for moderate talent to judge. For we are not now setting out some art of speaking, but only handing down to most learned men, as it were, certain monitions of our practice. With these topics, then, fixed in the mind and thought, and roused in every matter set down for speaking, there will be nothing that can escape the orator, not only in forensic disputes, but in any kind of speaking at all.

\ciceroSection{176} If indeed he attain that he seem such as he wishes himself to seem, and so affect the minds of those before whom he is to plead that he can either lead or carry them off whithersoever he wishes — there will surely be nothing else he shall require for speaking. Now we see that it is by no means enough to find what to say, unless one can handle what one has found.

\ciceroSection{177} The handling, however, must be varied — lest he who hears either recognise the art or be wearied by satiety of likeness. One ought to set forth what one brings, and to show why it is so; and from those same topics sometimes to conclude, sometimes to leave them and pass to others; often not to set forth and, by the very reasoning of what has been brought forward, to declare what would have to be set forth; if you say something is like, first to confirm that you say it as a likeness, then to add what is at hand. The pauses between arguments must mostly be hidden, lest anyone be able to count them, so that they may be distinguished by the matter, and seem confused in the words.

\ciceroSection{178} These things I run through, hastening (both before learned men, being half-learned myself), that we may at last come to those greater things. For there is nothing greater in speaking, Catulus, than that he who shall hear the orator should favour him, and that he himself should be so moved as to be ruled by some impulse of soul and disturbance, more than by judgment or counsel. For men judge much more from hatred or love or desire or anger or grief or joy or hope or fear or error or some movement of the mind, than from truth or precept or any norm of right or formula of judgment or the laws.

\ciceroSection{179} Therefore, unless something else pleases you, let us pass to those things.'' ``A little,'' said Catulus, ``something still seems lacking, Antonius, in those matters which you have set out. There is something you should set out before you go to where you say you intend.'' ``What?'' he said. ``What order pleases you,'' said Catulus, ``and what arrangement of arguments — in which you always seem to me a god.''

\ciceroSection{180} ``See how I am a god in that kind, Catulus,'' he said. ``By Hercules it would not have come into my mind unless I had been reminded; that you may judge that, in those things in which I sometimes seem to do something, I am wont in speaking to fall on it by use, or rather by chance. And that very thing — which, because I did not know it, I was passing over as one passes by an unknown man — has so much power in speaking that nothing has more for the winning. But still you seem to me to have asked of me before time the rationale of order and arrangement of subjects.

\ciceroSection{181} For if I had set the orator's whole force in arguments and in proving the matter itself by itself, it would now be time to say something about the order of arguments and about the arrangement of matters. But since three things have been proposed by me, and one has been spoken of, when I shall have spoken of the two remaining, then at last we must inquire about the disposing of the whole speech.

\ciceroSection{182} It is therefore of much weight for victory that the morals and habits and deeds and life of those who plead the cases, and of those for whom they plead, should be approved; and likewise that the morals of the adversaries should be disapproved; and that the minds of those before whom the case will be pleaded should be conciliated as much as possible to good will both toward the orator and toward him for whom the orator is to speak. Minds are conciliated by the dignity of the man, by his deeds, by the esteem of his life: which can be more easily adorned, if they are at all, than feigned, if there are none. But these things help in the orator: gentleness of voice, a sign of modesty in the face, kindness in his words; if you pursue anything more keenly, that you seem to do it unwillingly and under compulsion. To bring forth signs of his easiness, generosity, mildness, piety, gratitude, of being not appetitive, not greedy — is most useful. And all those things which belong to upright, modest, not keen, not pertinacious, not litigious, not bitter men greatly conciliate goodwill, and turn it away from those in whom these are not. So those same things must be applied conversely to the adversaries.

\ciceroSection{183} But this whole kind of speech excels in those cases where the judge's mind cannot greatly be inflamed by some keen and vehement excitement. For not always is a strong speech sought, but often a placid, modest, gentle one, which most commends the parties. The parties I call not only those who are accused but all whose case is in dispute: for so they used to speak in old time.

\ciceroSection{184} It is wondrously powerful to express in speech the morals of these — as just, upright, religious, modest, bearing wrongs. And this, whether in the openings, or in narrating the matter, or in the peroration, has so great a force, if it has been handled sweetly and with feeling, that often it avails more than the case. So much is brought about by a kind of feeling and rationale of speaking, that the speech, as it were, models the morals of the orator. For by a certain kind of thoughts and a certain kind of words, with a delivery also gentle and signifying easiness brought to bear, it is brought about that we seem to be upright men, of good morals, good men.

\ciceroSection{185} To this, joined as its opposite, is the other rationale of speech, which by a certain other kind moves and impels the minds of judges so that they either hate, or are fond, or envy, or wish for safety, or fear, or hope, or desire, or shrink, or rejoice, or grieve, or pity, or wish to punish, or are led to those motions which are neighbouring or close to such movements of the soul.

\ciceroSection{186} And this is to be wished for the orator: that the judges of their own accord bring to the case some movement of mind suited to that which the orator's interest will bear. For it is easier, as they say, to spur on a man already running than to rouse the slack. But if either there is no such movement, or it is too obscure — as a diligent doctor, before he tries to apply medicine to the sick man, must learn not only the disease of him whom he wishes to heal but also the habit of the man's health and the nature of his body — so I, when I begin in a doubtful and weighty case to handle the minds of the judges, dwell with all my mind on this thought and care: to scent out, as keenly as I can, what they feel, what they suppose, what they expect, what they wish, whither they seem most easily able to be led by speech.

\ciceroSection{187} If they yield and, as I said before, of their own accord lean and incline to where I am pushing, I receive what is given, and I set my sails to where some breeze is showing itself. But if the judge is unmoved and at rest, there is more work; for everything must be roused by speaking, with no help of nature. But that thing has so great a force which has rightly been called by a good poet ``mind-bending speech, queen of all things'' — that it can not only catch the leaning man and bend the standing one, but also, like a brave and good general, can capture even the resisting and reluctant.

\ciceroSection{188} These are the things which Crassus, jesting, was just now requiring of me, when he said I was wont to handle them divinely, and was praising as it were splendidly performed in the case of M.' Aquilius and Gaius Norbanus and some others — which, by Hercules, when they are handled by you, Crassus, in cases, I am wont to shudder at: so great is the force of mind, so great the impulse, so great the grief that is wont to be signified by your eyes, face, gesture, and at last that finger of yours; so great a river of the weightiest and best words, so upright thoughts, so true, so new, so without childish paint and rouge, that you seem to me not only to set the judge on fire but yourself to burn.

\ciceroSection{189} Nor can it happen that he who hears should grieve, hate, envy, fear, or be drawn to weeping and pity, unless all those motions which the orator wishes to apply to the judge are seen to be impressed and burnt into the orator himself. If some feigned grief had to be undertaken, and if in this kind of speech there were nothing but false and feigned by imitation, perhaps a greater art would have to be sought. Now I do not know, Crassus, what happens to you and the rest. As for myself, there is no reason to lie among most prudent and most friendly men: by Hercules I have never wished to rouse before the judges, by speaking, either grief or pity or envy or hatred, without being myself in the moving of the judges moved by those very feelings to which I wished to bring them.

\ciceroSection{190} For it is not easy to bring about that a judge be angry at the man you wish, if you yourself seem to bear it slowly; nor that he hate the man you wish, unless he has first seen you yourself blazing with hatred; nor will he be drawn to pity, unless you have shown him signs of your grief — by words, thoughts, voice, face, even by tears. For as no material is so easy to kindle that it can catch fire without fire put to it, so no mind is so prepared to grasp the orator's force that it can be inflamed unless he himself comes to it inflamed and burning.

\ciceroSection{191} And lest perhaps it seem a great and wonderful thing for a man so often to be angry, so often to grieve, so often to be moved by every motion of mind, especially in others' affairs — there is a great force in those thoughts and topics which one pleads and handles in speaking, so that there is no need of feigning and pretence. For the very nature of that speech which is undertaken to move others' minds moves the orator himself even more than any of those who hear him.

\ciceroSection{192} Lest we wonder that this happens in cases, in trials, in our friends' dangers, in concourse of men, in the state, in the Forum, when at stake is not only the esteem of our talent (for that would be lighter — although, when you have professed yourself able to do what few can, that is not to be neglected) — but other things much greater are at stake: faith, duty, diligence; led by which, even when we defend men most foreign to us, we cannot still — if we ourselves wish to be reckoned good men — judge

\ciceroSection{193} them as foreign to us. But, as I said, lest this seem wonderful in us — what can be so feigned as a verse, as a stage, as a play? Yet in this kind I myself have often seen the actor's eyes seem to burn through the mask, when he was saying the famous lines: \textit{`Did you dare to set yourself apart from him, or to enter Salamis without him? And not fear your father's face?'} — never did he say ``the face,'' but Telamon angry seemed to me to rage with grief for his son. And when the same man, with his voice bent to a piteous note, said, \textit{`In the worn-out years of life you tore me, bereft me, extinguished me, neither did the brother's death...'} — and ``the death of his small son, who was given over to your guardianship,'' — he seemed to be speaking weeping and lamenting. If that actor, when he played daily, could not still play rightly without grief, do you suppose Pacuvius was, in writing, of a calm and slack soul?

\ciceroSection{194} It could in no way be so. For I have often heard that no one is a good poet — what Democritus and Plato are said to have left in their writings — without inflammation of mind, and without a kind of breath, as it were, of madness. Therefore do not think that I myself, who do not wish in speaking to imitate and shadow forth the old falls of heroes and feigned griefs, and am no actor of another's part but the author of my own — when M.' Aquilius had to be kept by me in the state, did what I did in the peroration of that case without great grief.

\ciceroSection{195} The man whom I remembered to have been consul, a general adorned by the senate, to have ascended in ovation to the Capitol — this man when I saw stricken, broken, mourning, brought to the highest peril, I did not first try to stir pity in others until I myself was caught by pity. I felt then indeed that the judges were greatly moved, when I roused that mournful and squalid old man, and when I did those things which, Crassus, you praise — not by art (about which I do not know what to say), but by great motion of soul and grief: that I tore his tunic, that I showed his scars.

\ciceroSection{196} When Gaius Marius, present and seated, much aided the grief of my speech with his own tears, and when I, often calling him by name, commended his colleague to him and called the man himself as advocate to defend the common fortune of generals, this pity and imploring of all the gods and men and citizens and allies was not without my tears, not without great grief. And of all those words which were uttered by me then, if my own grief had been absent, my speech would not only have not been pitiable, but even ridiculous. Therefore this I teach you, Sulpicius — I, the good and learned master, of course — that in speaking you may be able to be angry, to grieve, to weep.

\ciceroSection{197} Although why should I teach you this — you who, in accusing my friend, had, by the force, grief, and ardour of your soul rather than by speech alone, kindled so great a fire that I scarcely tried to come near to put it out? For you had everything in the case as your superior: violence, flight, stoning, the cruelty of the tribune in Caepio's heavy and pitiable fall, you brought to judgment; then it was agreed that the leader both of the senate and of the state, M. Aemilius, had been struck with a stone; that L. Cotta and T. Didius had been driven by violence from the temple, when they wished to interpose against the bill, no one could deny.

\ciceroSection{198} It was added that you, a young man, were thought to make these complaints for the commonwealth with the highest dignity; while I, a man of censorial rank, hardly seemed honourable enough to be able to defend a seditious citizen who had been cruel in the misfortune of a consular man. The judges were the best citizens, the Forum was full of good men, so that some thin pardon could scarcely be given me for excusing myself, in that I yet defended one who had been my quaestor. Here why should I say I used some art? What I did, I shall tell. If you please, you shall set my defence in some place of the art.

\ciceroSection{199} I gathered all the kinds, faults, and dangers of seditions; and I drew that speech back from every variety of times of our commonwealth, and I closed it so that I should say: although seditions had always been all troublesome, some had still been just and almost necessary. Then I argued those things which Crassus has just been recalling: that the kings could not have been driven out from this state, nor the tribunes of the plebs created, nor the consular power so often diminished by plebiscites, nor the right of appeal — that patroness of the state and avenger of liberty — given to the Roman people, without a disagreement of the nobles. And, if those seditions had been to the safety of this state, it was not at once, if any popular movement had been made, to be put in the case of Gaius Norbanus as a wicked charge and capital fraud. If it had ever been granted to the Roman people that they seemed justly stirred (which, I was teaching, had often been granted), no case had been more just than that one. Then I turned all my speech and converted it to rebuking Caepio's flight, to lamenting the destruction of the army. So both I rubbed afresh, by my speech, the grief of those who were mourning their own; and I renewed the spirit of the Roman knights, before whom as judges the case was being pleaded, to a hatred of Q. Caepio, by whom they themselves had been alienated on account of the courts.

\ciceroSection{200} When I felt that I had taken hold of the judgment and of my defence — for I had won the people's good will to myself, whose right I had defended even with the tagging on of sedition; and the minds of the judges I had wholly turned to our case, both by the calamity of the state, and by their grief and longing for kinsfolk, and by their own hatred for Caepio — then I began to mix into this vehement and grim kind of speech that other kind of which I argued before, of gentleness and mildness; that for my friend, who, by ancestral custom, ought to be in the place of a child to me, and for my own whole reputation and almost my fortunes I was contending; that nothing more shameful for esteem, nothing more bitter for grief, could befall me than this — that he who was reckoned often to have been the salvation of men most foreign to me but yet my fellow-citizens, could not bring help to my friend.

\ciceroSection{201} I asked the judges that they should grant this to my age, to my offices, to my deeds, if they saw me affected by just and pious grief — especially if they had understood in other cases that I had always sought everything for the dangers of friends, never anything for myself. So in that whole defence and case, what seemed to lie in art — to speak about the Appuleian law, to set out what it was to ``impair majesty'' — I most briefly grazed and touched. With these two parts of speech (one of which has the conciliation, the other the rousing), which are least polished by the precepts of arts, the whole case was handled by me, that I might seem most keen in renewing hatred for Caepio and most mild in declaring my morals towards my own intimates. So with the minds of the judges affected rather than instructed, your accusation, Sulpicius, was overcome by us at that time.''

\ciceroSection{202} Here Sulpicius: ``Truly, by Hercules, you recall those things, Antonius. For I never saw anything that so slipped from the hands as that very case slipped from me then. For when, as you said, I had handed over to you not a trial but a conflagration — what an opening of yours, immortal gods! What fear! what hesitation! What halting and dragging of words! How you held that one ground at the start, which men gave you alone for pardon, that you spoke for a man so close to you, your former quaestor! How you first made for yourself a way to be heard.

\ciceroSection{203} But behold, when I thought you had achieved nothing else save that men should think you must be pardoned for defending a wicked citizen because of intimacy, you began to creep on secretly, while none of the others as yet suspected, but I now thoroughly afraid — that you should defend that not as Norbanus's sedition but as the Roman people's anger; and that not unjust but deserved and owing. Then what topic was passed by you against Caepio? How did you mix all those things — hatred, envy, pity! Nor only these in the defence, but also in Scaurus and the rest of my witnesses, whose testimonies you refuted not by refuting them but by fleeing back to that same impulse of the people.

\ciceroSection{204} When these things were just being recalled by you, indeed I asked for no precepts. Yet that very setting forth of your defences, recalled by you yourself, I think to be a teaching of no slight kind.'' ``Well, then,'' said Antonius, ``if it pleases you so, we shall hand down also what we are wont to follow in speaking and what most we are wont to look at. For long life and the practice of the greatest matters has now taught us so that we hold what things move men's minds.

\ciceroSection{205} For my part I am wont first to consider whether the case requires it. For these torches of speaking are not to be applied either in small matters or to men so disposed that we can do nothing to bend their minds by speech — lest we be thought worthy of derision or hatred, if we either play tragedy in trifles, or attack what cannot be moved.

\ciceroSection{206} Now since these are mostly to be wrought by speech in the minds of judges — or whoever they be before whom we shall plead — love, hatred, anger, envy, pity, hope, joy, fear, distress: we feel love is conciliated if we seem rightly to defend what is useful to those very men before whom we plead, or if we labour for good men, or at least for those who are good and useful to them. For the latter conciliates love more, the former (the defence of virtue) regard. And it profits more if there is set forth a hope of future use than the recalling of past kindness.

\ciceroSection{207} You must strive to show in the matter you defend that there is either dignity or usefulness; and you must signify that the man for whom you would conciliate this love has referred nothing to his own use and has done nothing wholly for his own sake. For envy is felt at men's own advantages; while their zeal in serving others wins favour.

\ciceroSection{208} And here it must be looked to that we do not, for those whom we wish to be loved on account of good deeds, lift up too much their praise and glory (against which envy is most felt). And from these same topics we shall learn both to build hatred against others, and to remove it from us and ours. The same kinds must be handled in rousing or calming anger. For if you should magnify what is pernicious or useless to those who hear, hatred is created. But if it has been done either against good men, or against those whom no one ought least to have offended, or against the commonwealth, then is roused, if not so bitter a hatred, yet a slight not unlike envy or hatred;

\ciceroSection{209} likewise fear is struck either from their own dangers or from common: the inner is the more proper, but this common one too must be brought back to the same likeness. The rationale of hope, joy, distress is equal and one. But I do not know whether of all motions the keenest by far is envy, and whether no less force is to be expended in repressing than in rousing it. Men envy most those equal or below them, when they feel themselves left behind, but grieve that those have flown up. But also superiors are often vehemently envied, and the more so if they bear themselves the more intolerantly and pass the equality of common right by the excellence of their dignity or fortune. If these things are to be inflamed, it must especially be said that they were not won by virtue; then by vices and faults; finally, if they will be more honourable and weighty, that those merits are still not so great as the man's insolence and disdain;

\ciceroSection{210} but for the calming, that those things were got by great labour, by great dangers, and not for their own use but for others'; and that the man, even if he seems to have produced some glory, although that be not unjust pay for the danger, is yet not delighted by it and casts it wholly away and lays it down. And on the whole it must be brought about — since most are envious, and this is the chief and lasting common vice, while the outstanding and flourishing fortune is what is envied — that this opinion should be lessened, and that fortune, outstanding in repute, should seem mixed with toils and miseries.

\ciceroSection{211} Pity is moved if he who hears can be drawn so as to call back to his own affairs (which he has either borne bitterly or fears) those things which are lamented about another, so that, looking at another, he often returns to himself. And while single chances of human miseries are taken hard if they are spoken with grief, broken and prostrate virtue is the most mournful. And as that other part of speech, which by the recommendation of uprightness ought to keep up the appearance of the good man, ought (as I have often said) to be gentle and modest, so this, which is undertaken by the orator for changing minds and bending them by every reasoning, ought to be intent and vehement.

\ciceroSection{212} But there is in these two kinds — of which we wish the one to be gentle, the other vehement — a likeness hard to distinguish. For from that gentleness, by which we are conciliated to those who hear us, something must flow into this keenest force, by which we rouse those same; and from this force sometimes something of mind must be breathed into that gentleness. And there is no more tempered speech than that in which the harshness of contention is seasoned by the orator's own kindness, and the relaxation of gentleness is firmed up by a certain weight and contention.

\ciceroSection{213} In each kind of speaking — both that in which force and contention are sought, and this which is fitted to life and morals — both the openings should be slow, and the endings likewise close-set and prolonged. For neither should one leap at once into that kind of speech (it is wholly absent from the case, and men first wish to hear that very thing which is proper to the trial); nor, when you have entered into that line, should one quickly depart.

\ciceroSection{214} For not as an argument, as soon as it has been laid down, is grasped, and a second and third demanded — so pity or envy or anger, as soon as you have brought them in, can you stir up. Reason itself confirms the argument, and this, as soon as it has been let go, sticks. But that kind of speech requires not the judge's understanding but rather his disturbance, which no one can attain save by much, varied, copious speech and equal contention of action.

\ciceroSection{215} Therefore those who speak either briefly or modestly can teach the judge but cannot move him — in which is everything. Now this is plain, that the faculty for opposite parts is supplied from the same topics for all subjects. But an argument must either be resisted by reproving those things which are taken for proving it, or by showing that what they wish to conclude does not follow from what is laid down nor is consequent; or, if you do not so refute it, something must be brought on the contrary side which is either heavier or equally heavy.

\ciceroSection{216} Those things, however, which are done either gently for the sake of conciliation or vehemently for stirring, must be carried off by contrary movements: hatred by goodwill, pity by envy. Sweet and often vehemently useful is jest and wit. Even if all other things can be handed down by art, these are surely proper to nature and require no art — in which you, Caesar, by far excel others in my opinion. So you can rather be a witness for me that there is no art of salt, or, if there is, you above all are to teach it to us.''

\ciceroSection{217} ``For my part,'' he said, ``I think it easier for a man not without urbanity to argue about every other matter than about wit itself. So when I had seen certain Greek books inscribed `On Laughable Things,' I had come into some hope that I could learn something from them. I found indeed many laughable and salty things of the Greeks; for the Sicilians excel in this kind, and the Rhodians and Byzantines, and beyond the rest the Athenians. But those who tried to hand down some method and art of this matter, stood out so dull that nothing of theirs is laughed at except their very dullness;

\ciceroSection{218} therefore this matter does not seem to me in any way to be possible to hand down by teaching. Since there are two kinds of wit, the one diffused evenly through the whole speech, the other most keen and brief — that former by the ancients was called \textit{cavillatio} (banter), the latter \textit{dicacitas} (smartness). Each thing has a light name. Of course; for the whole business of stirring laughter is light.

\ciceroSection{219} But yet, as you say, Antonius, I have very often seen much in cases got by charm and wit. But while in that kind of unbroken pleasantness no art is needed (for nature shapes men and creates them as imitators and witty narrators, with face and voice and the very kind of speech helping), what room has art in this other of \textit{dicacitas}, since the witty saying ought to stick before it has been let go, sooner than it could seem to have been thought of?

\ciceroSection{220} For what could my brother here have been helped by art, when, asked by Philippus what he was barking at, he answered that he saw a thief? What in any speech of Crassus, either before the Hundred Men against Scaevola or against the prosecutor Brutus, when he was speaking for Cn. Plancius? For what you assign to me, Antonius, must by everyone's vote be conceded to Crassus. For hardly anyone besides him will be found outstanding in both kinds of charm: both the one which is in the continuity of speech, and this which is in swiftness and saying.

\ciceroSection{221} For that whole continuous defence in the Curian case against Scaevola overflowed with cheerfulness and jest. Brief sayings it had not, for he spared the dignity of his adversary, in which he was preserving his own. This is most difficult to witty and quick-tongued men: to have regard for men and times, and to hold back things which present themselves when they could be said most pungently. So some witty men interpret this not unwittily;

\ciceroSection{222} for they say Ennius said the wise man could more easily smother flame in his burning mouth than hold back \textit{good sayings} — meaning, of course, those good sayings which are salty; for these are now called by the proper name of \textit{sayings}. But just as in the Scaevola case Crassus held them in, and in that other kind, in which there were no points of insult, he played out the case and the disputation, so in the Brutus case, whom he hated and judged worthy of insult, he fought with both kinds.

\ciceroSection{223} How many things did he say about the baths Brutus had lately sold, how many about his lost patrimony! And the brief sayings: when he said he was sweating without cause, ``Not at all,'' said Crassus, ``you have just come out of the baths.'' Innumerable were the others; but no less pleasant was that continuous portion: for when Brutus had set up two readers and given to one Crassus's speech ``On the Narbonese colony'' to read, to the other ``On the Servilian Law,'' and had compared between them contradictory passages from the commonwealth — our friend most wittily gave his three little books on the civil law of Brutus's father to three others to read.

\ciceroSection{224} From the first book: ``It happened that we were on my Privernum estate.'' ``Brutus, your father bears witness that he left you the Privernum farm.'' Then from the second book: ``We were at the Alban estate, I and my son Marcus.'' ``A wise man indeed of the first rank in our state, he knew this gulf of a son. He feared lest, when this man had nothing, it should be thought he had been left nothing.'' Then from the third book — which he made the end of his writing (for so many, as I have heard Scaevola say, are the genuine books of Brutus): ``We sat at our Tibur estate, by chance, I and my son Marcus.'' ``Where are these farms, Brutus, which your father left you signed for in his public commentaries? But for the fact, he says, that he had you yet but underage, he would have written a fourth book and would have left it written that he himself

\ciceroSection{225} bathed at the baths with his son.'' Who, then, would not confess that by this charm and these witticisms Brutus was no less refuted than by those tragic strokes the same Crassus performed when, by chance in that same trial, the old Junia woman was being carried out for burial? By the immortal gods, what a force was that! how unexpected! how sudden! when, with eyes fixed, with every gesture pressing upon him, with the highest weight and swiftness of words: ``Brutus, why do you sit? What do you wish that old woman to announce to your father? What to all those whose images you see led? What to your ancestors? What to L. Brutus, who freed this people from kingly mastery? What to do with you? In what affair, in what glory, in what virtue do you study? In increasing your patrimony? But that is not of the nobility.

\ciceroSection{226} But suppose it were — there is nothing left; lusts have wholly squandered it. In civil law? It is a paternal study. But it will be said that, when you sold your house, you did not even keep back among the fixtures and movables your father's bath-tub. In military matters? You who have never seen a camp! In eloquence? Which is not in you, and whatever there is of voice and tongue, all you have brought to that most disgraceful trade of false accusation. Do you dare to look on the day? Dare to look at these men? In the Forum, in the city, in the sight of citizens? Are you not in awe of the dead woman, of the very ancestors' images? Whom you have left yourself no place not only to imitate, but not even to set up.

\ciceroSection{227} Yet these are tragic and divine; the witty and urbane sayings without number, even out of one assembly, you remember; for there has never been a greater contest, nor weightier speech before the people, than his lately against his colleague in the censorship, nor any seasoned with more charm and merriment. Therefore, Antonius, I agree with you on both points: that wit much profits in speaking, and that it can in no way be handed down by art. This indeed I wonder at: that you have assigned so much to us in that kind, and not given to Crassus the palm of this matter too, as of the rest.''

\ciceroSection{228} Then Antonius: ``I should have done so,'' he said, ``had I not sometimes envied this Crassus a little. To be ever so witty and salty is not in itself too much to be envied; but to be at once the most graceful and the most urbane of all, and at the same time the most weighty and the most severe of all, both to be and to seem — which has fallen to him alone — that seemed to me hardly to be borne.''

\ciceroSection{229} Here, when Crassus himself had smiled, ``And yet,'' said Antonius, ``Iulius, when you were denying that there is any art of wit, you opened up something which seemed worthy of precept. For you said that account must be had of men, of the matter, of the time, lest the jest take anything from one's gravity. This indeed Crassus is wont chiefly to observe. But this is the precept of leaving aside witticisms when there is no need of them; we, however, ask how to use them when there is need: as against an adversary, and especially if his folly can be played on; against a witness foolish, greedy, frivolous, if men shall seem ready to hear it.

\ciceroSection{230} On the whole the things we say when we are challenged are more probable than what we say first; for both the swiftness of talent is the greater which appears in answering, and the answer is of human kindness. For we seem to have been ready to keep silent, had we not been challenged. As in that very assembly nothing was said by him which seemed wittier, that he did not answer to a challenge. There was, however, in Domitius such weight, such authority, that what was thrown out by him seemed rather to be lightened by charm than broken by contention.''

\ciceroSection{231} Then Sulpicius: ``What then? Shall we suffer Caesar,'' he said, ``— who though he yields wit to M. Crassus, yet himself works much more in this study — not to set out for us this whole kind of joking: what it is and whence it is drawn? — especially since he confesses that the force and usefulness of salt and urbanity are so great?'' ``What if,'' said Iulius, ``I agree with Antonius when he says there is no art of salt?''

\ciceroSection{232} Here, when Sulpicius had fallen silent, ``As if,'' said Crassus, ``of these things themselves of which Antonius has long been speaking there were any art. There is a certain observation, as he himself said, of those things which avail in speaking. If this could make men eloquent, who would not be eloquent? For who could not learn these things either easily or surely in some way? But I think there is in these precepts this force and this usefulness: not that we should be led by art to discovering what to say, but that those things which we attain by nature, by zeal, by exercise, we may either trust to be right or understand to be faulty, when we have learned with what they should be referred.

\ciceroSection{233} Therefore, Caesar, I too ask you: that, if you please, you should argue about this whole kind of joking, what you feel — lest some part of speaking, since you have so wished, in this such gathering and in so careful a conversation, seem to have been passed over.'' ``For my part,'' said he, ``since you exact the contribution from a guest, Crassus, I shall not bring it about that, by fleeing, I should give you any reason for refusing. Although I am wont often to wonder at the impudence of those who play their part on the stage when Roscius is looking on; for who can move himself, whose faults that man does not see? So I now, with Crassus listening, shall first speak about wit, and shall teach as it were the sow how to talk to the orator — whom, when Catulus had lately heard, he said the others ought to be grass.''

\ciceroSection{234} Then he: ``Catulus was joking,'' he said, ``especially since he himself so speaks that he seems should be fed on ambrosia. But let us hear you, Caesar, that we may see Antonius's remaining matter.'' And Antonius: ``Very few things remain to me,'' he said, ``but yet, weary now from my labour and journey of disputation, I shall rest in Caesar's speech as in some most opportune wayside inn.'' ``Yet,'' said Iulius, ``you will say my hospitality is not too generous; for as soon as you have tasted ever so little I shall thrust you out and cast you on the road.

\ciceroSection{235} And lest I keep you longer, I shall set out very briefly what I think of this whole subject. About laughter there are five things to be asked: first, what it is; second, whence it is; third, whether it is the orator's part to wish to stir laughter; fourth, how far; fifth, what the kinds of laughable things are. As for the first — what laughter itself is, by what means it is roused, where it lies, in what way it arises and so suddenly bursts out that we cannot hold it back when we wish, and how at once it seizes the sides, the mouth, the veins, the eyes, the face — let Democritus see to it. For this does not concern this conversation; and even if it concerned it, I should not be ashamed of not knowing what not even those knew who promised it.

\ciceroSection{236} The seat and as it were the region of the laughable — this is what we ask next — is contained in a certain ugliness and deformity. For these things are laughed at, alone or chiefly, which mark and point out some ugliness, but not in an ugly way. As for the third question — it is plainly the orator's part to stir laughter, either because cheerfulness itself wins goodwill for the man through whom it has been roused; or because all admire the keenness often set in one word, especially of him who answers, sometimes also of him who challenges; or because it breaks the adversary, hinders him, lightens him, deters him, refutes him; or because it shows that the orator himself is a polished man, learned, urbane, and most of all because it softens and relaxes gloom and severity, and often dissolves by jest and laughter hateful matters which are not easily to be dissolved by arguments.

\ciceroSection{237} How far the laughable is to be handled by the orator must be looked at most diligently — which we put in the fourth place. For neither marked wickedness joined with crime, nor in turn marked misery, is the subject of laughter. For the criminal men prefer to be wounded by some greater force than that of the laughable; the wretched they are unwilling to mock, unless they happen to flaunt themselves. We must spare especially the affection in which men are held, lest we speak rashly against those who are loved.

\ciceroSection{238} This moderation, then, must first be applied in jesting; therefore those things are most easily made play of which deserve neither great hatred nor the greatest pity. For which reason all the matter of the laughable is in those vices which are in the life of men neither dear nor calamitous, nor of those who seem to be hurried off to punishment for crime; and these, prettily handled, are laughed at.

\ciceroSection{239} The deformity and faults of the body are also a sufficiently pretty matter for jesting; but we ask the same as in other matters most must be asked: how far. In which not only is this prescribed — that nothing should be done dully — but also, even if you can speak most laughably, both must be avoided by the orator: that the jest be neither buffoonish nor mimic. What sort these are, we shall now more easily understand when we come to the very kinds of the laughable. For there are two kinds of wit, one of which is handled in a thing, the other in a saying:

\ciceroSection{240} \textit{in a thing}, whenever something is told as in a kind of little story — as you once, Crassus, against Memmius: that he had bitten the arm of Largus when at Tarracina they had quarrelled with him about a girl. The whole story was salty, and yet wholly invented by yourself. You added a closing piece: that throughout Tarracina then on all the walls had been inscribed the letters ``L. L. L. M. M.''; and when you asked what this was, some old townsman had told you,

\ciceroSection{241} ``\textit{Lacerat lacertum Largi mordax Memmius}'' — ``Snappy Memmius lacerates the arm of Largus.'' You see how witty this kind is, how elegant, how oratorical — whether you have something true to tell, which is yet to be sprinkled with little fictions, or whether you invent. The virtue of this kind is that you should so set forth what was done that the morals of the man you are telling about, his speech, his face, all things, may be expressed; that to those who hear, those things may seem to be then and there happening and being done.

\ciceroSection{242} The laughable is also \textit{in the thing} which is wont to be drawn from a certain corrupted imitation, as the same Crassus: ``By your nobility, by your family!'' What else made the assembly laugh, but that imitation of his face and voice? But when he said, ``By your statues!'' and added a little even of gesture by stretching out his arm, we laughed more vehemently. Of this kind is that famous Roscian imitation of an old man: ``\textit{For you, Antipho, do I sow these things,} he says; \textit{it is doting old age, when I hear.}'' Yet this whole kind is so laughable that it must be handled most cautiously. For it belongs to mimes and to character-actors, if the imitation is too much — like obscenity. The orator should snatch the imitation, that he who hears may think more than he sees; he should also keep up his free-born manner and modesty by avoiding ugliness of words and obscenity of things.

\ciceroSection{243} These are the two kinds of the laughable which is set in the thing, which are proper to continuous wit, in which men's morals are described and so shaped that — either by some matter narrated, of what kind they are, may be understood, or by an imitation briefly thrown in upon some marked vice — they may be detected for derision.

\ciceroSection{244} \textit{In the saying}, however, the laughable is what is moved by some keenness of word or thought; but as in that earlier kind of either narration or imitation the likeness of mimes and character-actors must be avoided, so in this the buffoonish quickness of tongue must be most greatly fled by the orator. How then shall we distinguish from Crassus, from Catulus, from the rest, your friend Granius, or my friend Vargula? It does not, by Hercules, come into my mind. For they are quick-tongued; than Granius, indeed, none was quicker. This, I think, first: that we should not feel bound to say a saying every time one can be said.

\ciceroSection{245} A tiny witness came forward. ``May I question him?'' said Philippus. Then the questor, hurrying: ``Only briefly.'' Hereupon he: ``You will not accuse me — I shall question him very tinily.'' Witty. But the judge sitting was L. Aurifex, shorter himself than even the witness; the whole laugh was turned upon the judge: the witticism seemed wholly buffoonish. Therefore those things which can fall on those whom you would not wish, however pretty they are, are still by their kind buffoonish;

\ciceroSection{246} as that fellow who wishes to be quick-tongued — and by Hercules he is — Appius; but sometimes he slips into this buffoonish fault. ``I will dine with you,'' he said to my friend Gaius Sextius, who has only one eye, ``for I see there is room for one only.'' This is buffoonish, both because he challenged without cause, and because he yet said what would fit all one-eyed men. These things, because they are thought premeditated, are less laughed at. That outstanding answer of Sextius, on the spot: ``Wash your hand,'' he said,

\ciceroSection{247} ``and dine.'' So observation of the time and the moderation of the wit itself, and the temperance and rarity of sayings, will distinguish the orator from the buffoon — also that we speak with a cause, not so as to seem laughable, but so as to gain something; while they speak the whole day and without cause. For what did Vargula gain when, when the candidate A. Sempronius had embraced him with his brother Marcus: ``Boy, drive off the flies''? He was hunting laughter — which is, in my judgment, the slenderest fruit of talent. The time of speaking, then, we shall measure by prudence and gravity. Would that we had some art for these! But mistress is nature.

\ciceroSection{248} Now let us set out summarily the kinds themselves which most stir laughter. Let this then be the first division: that what is wittily said has its wit at one time in the matter, at another in the word; men are most delighted, however, if at any time the laughter is moved jointly by matter and word together. But remember this: whatever topics I shall touch, from which the laughable is drawn — from those same topics weighty thoughts can also be drawn. So great is the difference, that gravity is set in honest and severe matters, jest in slightly base and as it were misshapen ones — as we may both praise an honest slave with the same words, and, if he is worthless, joke. Laughable is that old saying of Nero about a thievish slave: that he was the only one in whose case nothing at home was either sealed or shut up — which same thing is wont to be said of a good slave.

\ciceroSection{249} But this is by the same words; from the same topics arise all things. For when Sp. Carvilius — limping heavily from a wound received for the commonwealth, and on that account ashamed to come out in public — was told by his mother: ``Why don't you come out, my Spurius? As often as you take a step, may the memory of your virtues come to you'' — that is splendid and weighty. But when Glaucia said to Calvinus, who was limping: ``Where is that old saw, `Surely he is not limping (\textit{claudicat})?' But here he is limping (\textit{clodicat})!'' — that is laughable. And both are drawn from what could be observed in the limping. ``What is more lazy (\textit{ignavius}) than this Navius?'' Scipio said sternly. But to a foul-smelling man Philippus said, ``I see I am being beset (\textit{circumveniri}) by you'' — somewhat laughably. But each kind contains the likeness of a word changed by a single letter.

\ciceroSection{250} Sayings from ambiguity are reckoned the most clever, but they do not always lie in jest, often also in gravity. To the elder Africanus, when he was fitting on a wreath for his head at a banquet, when it kept breaking, P. Licinius Varus said: ``Don't be surprised if it doesn't fit; for the head is great.'' Both praiseworthy and honourable. But of the same kind is: ``It is enough for Calvus that he says little.'' In short: there is no kind of jest from the same of which severe and weighty things may not be drawn.

\ciceroSection{251} This too must be observed: not all laughable things are witty. For what can be so laughable as a clown? But it is by mouth, by face, by imitating manners, by voice, finally by his very body that he is laughed at. I can call him salty, and so, not so that I should wish such an orator, but so that I should wish such a mime. Therefore this first kind, which most of all moves laughter, is not for us: the morose, the superstitious, the suspicious, the boastful, the fool — natures themselves are laughed at, characters which we are wont to set in motion, not to act.

\ciceroSection{252} The second kind is in imitation, very laughable, but only by stealth, if at all, and quickly may we use it; otherwise it is least free; the third, the distortion of the face, is not worthy of us; the fourth, obscenity, is not only not worthy of the Forum but scarcely worthy of a free man's banquet. So many things being subtracted from this oratorical topic, the witty things which remain are seen to be set either in the matter, as I divided before, or in the word. For what is witty in whatever words you use, is contained in the matter; what loses its salt when the words are changed, has all its charm in the words.

\ciceroSection{253} Ambiguous sayings are first of all keen and set in the word, not in the matter; but they do not often move great laughter — they are praised rather as prettily, as elegantly, said. As of that Titius, who, when he was busy playing ball, and was thought also to break sacred images by night, and his comrades were asking after him when he had not come to the field, Vespa Terentius excused him, saying that he had broken his arm. Or that of Africanus in Lucilius: ``\textit{What of Decius? Do you wish to make him a little nut, or pierced through?}''

\ciceroSection{254} Or, Crassus, that friend of yours, Granius: ``He's not worth a sixth.'' If you ask, he who is called \textit{quick-tongued} most excels in this kind; but other things move greater laughter. Ambiguity by itself is approved indeed, as I said, very greatly; for it seems the part of an ingenious man to be able to draw the force of a word into another sense than the rest take it. But it stirs admiration rather than laughter, unless it happens to fall in upon some other kind of laughable. Of which kinds I shall now run through.

\ciceroSection{255} You know that this is the most familiar kind of the laughable: when we expect one thing and another is said. Here our own error stirs the laugh in us. If ambiguity is also mixed in, it becomes saltier. As in Novius the man seems compassionate, who sees one being led off after judgment; he asks: ``For how much was he assigned?'' ``A thousand sesterces.'' If he had only added ``you may take him,'' that would be the kind of laughable beyond expectation; but because he added ``I add nothing — you may take him,'' with the ambiguous added, the other kind of laughable, the saying was, in my view, the saltiest. This is graceful when in altercation a word is snatched from the adversary and from it, as Catulus did to Philippus, something is hurled at the very man who challenged.

\ciceroSection{256} Since there are several kinds of ambiguity, about which there is some more subtle teaching, we shall have to be on the watch and to fowl for words; in which, if we avoid those that are colder (we must beware lest the saying seem fetched), we shall yet say very many things keenly. The second kind is that which has a slight change of word, what — set in a single letter — the Greeks call \textit{paronomasia}: as Cato's ``\textit{nobiliorem mobiliorem}''; or, when he had said to someone, ``Let us go for a walk (\textit{deambulatum}),'' and the other said, ``What was the need of \textit{de}?'' he answered: ``On the contrary — what was the need of you?'' Or that answer of his

\ciceroSection{257} ``If you are unchaste both \textit{adversus} (front) and \textit{aversus} (back).'' The interpretation of a name has also keenness, when you turn to the laughable why anyone is so called: as I lately, of Nummius the briber — that as Neoptolemus took his name at Troy, so this man took his name on the Campus Martius. All these are contained in the word. Often a verse is also wittily inserted, either as it is or with a little change, or some part of a verse — as that of Statius's by Scaurus when he was angered, from which there are some who say your law on citizenship was born, Crassus: ``\textit{St, be silent — what is this clamour? Whose neither mother nor father — such confidence? Take that pride away.}'' Now in Caelius's case, Antonius, what you said was useful even to the matter: when the witness was saying that the money had been put forward by him and he had a more dainty son, with the man already going off, ``Don't you feel, old man, that you have been touched for thirty minae?''

\ciceroSection{258} Into this kind are also thrown proverbs, as that of Scipio when Asellus was boasting that he had traversed all the provinces in his service: ``\textit{Drive the little ass}'' and the rest. So these too — since by changing the words they cannot keep the same grace — should be drawn as set not in the matter, but in the words.

\ciceroSection{259} Set also in the word is a not dull kind which arises from this — when you seem to take the matter according to the word, not according to the meaning. Of this one kind is wholly composed the old mime \textit{Tutor}, very laughable indeed. But I leave aside the mimes; only I wish this kind of laughable to be marked with some signal and noted matter. Of this kind is what you, Crassus, said to him who lately asked you whether you would be put out if he had come to you well before light: ``You will not be a bother to me,'' you said. ``You will then have me wakened?'' said he. And you: ``Surely I had said that you would not

\ciceroSection{260} be a bother.'' Of this same kind is that old saying: that Maluginensis, M. Scipio, when from his century he was returning Acidinus as consul, and the herald said, ``Pronounce on L. Manlius'': ``A good man,'' said he, ``and an outstanding citizen, I judge him.'' Laughable too is that of L. Porcius Nasica to Cato the censor, when the latter asked: ``According to the conviction of your soul, do you have a wife?'' ``By Hercules,'' he said, ``not according to my soul's conviction.'' These things are either cold or salty when something else has been expected. For nature, as I said before, delights us with our own error; from which, when we are as it were deceived in our expectation, we laugh.

\ciceroSection{261} In words also are those which are drawn either from speech changed, or from the transferred meaning of a single word, or from the inversion of words. From change: as once when Rusca was carrying a law on age, his opponent M. Servilius said: ``Tell me, M. Pinarius, if I shall speak against you, will you abuse me as you have done the others?'' He said: ``\textit{As you sow, so shall you reap.}''

\ciceroSection{262} From metaphor: as when the elder Scipio, when the Corinthians were promising him a statue in the place where there were others of generals, said the cavalry kind would not please. Words are inverted, as: when Crassus was speaking before the judge M. Perperna for Aculeo, on the other side, against Aculeo and for Gratidianus, was L. Aelius Lamia, deformed, as you know; when he was hatefully interrupting, ``Let us hear,'' said Crassus, ``the pretty boy''; when laughter had broken out, ``I could not,'' said Lamia, ``shape my own form, but my talent I could''; then Crassus: ``Let us hear the eloquent one,'' he said. There was a much more vehement laugh. These too are graceful — as in weighty thoughts, so in witticisms; for I said long ago that the rationale of jest is one and of severity another,

\ciceroSection{263} but that of weighty things and jests there is one matter. Words referred contrarily then most of all adorn a speech, the same kind being often also witty, as that famous Servius Galba: when he was nominating his own friends as judges to the tribune of the plebs L. Scribonius, and Libo said: ``When at last, Galba, will you come out of your dining-room?'' he said: ``When you, out of someone else's bedroom.'' Not far from this kind is that of Glaucia to Metellus: ``You have a villa at Tibur, a sheepfold (\textit{cors}) on the Palatine.'' I think I have said the kinds of words which are witty.

\ciceroSection{264} Of matters there are more, and they, as I said before, are more laughed at; in which is narration, a thing surely difficult; for those things must be expressed and set before the eyes which are likely and like the truth (which is proper to narration), and which are, what is proper to the laughable, somewhat ugly. An example, that it may be most brief, is that of Crassus on Memmius, which I set down before. To this kind let us also add the narrations of fables.

\ciceroSection{265} Something is also drawn from history, as when Sex. Titius said he was Cassandra: ``I,'' said Antonius, ``can name many of your Aiaxes son of Oileus.'' From likeness is also: which has either a comparison or as it were an image. A comparison: as that Gallus once, a witness against Piso, when he had said that countless money had been given to the prefect Magius, and Scaurus had refuted this by Magius's slenderness: ``You err, Scaurus,'' said he. ``For I do not say that Magius preserved it, but that as a man naked might gather nuts, so he carried it off in his belly.'' Or that of M. Cicero the elder, my friend's father, this excellent man's father: ``Our men are like Syrian slaves: as

\ciceroSection{266} each one knew Greek best, so was he the most worthless.'' Greatly laughed at also are images which are mostly drawn into deformity or some bodily fault, with the likeness of something more shameful: as that one of mine against Helvius Mancia: ``Now I shall show what kind of man you are.'' When he said, ``Show, please,'' I pointed to a Gaul painted on the Marian Cimbric shield under the Novae Tabernae, twisted, with tongue thrust out, with cheeks streaming. There was a laugh; nothing seemed so like Mancia. Or to Titus Pinarius, who twisted his chin in speaking: ``Then might he speak, if he wished, when he had broken a nut.''

\ciceroSection{267} Also those which, for the sake of diminishing or amplifying, are brought to incredible wonder; as you, Crassus, in the assembly: ``So great did Memmius appear to himself that, as he came down to the Forum, he ducked his head at the Fabian arch.'' Of which kind is also that which Scipio at Numantia, when he was angered with C. Metellus, is said to have said: ``If his mother

\ciceroSection{268} were to bear a fifth, she would bear an ass.'' Clever signification, too, is when by some small thing — and often by a single word — an obscure and hidden matter is illustrated; as when, when P. Cornelius — a man, as was thought, greedy and thievish but outstandingly brave and a good general — was thanking C. Fabricius because he, an enemy, had made him consul, especially in a great and serious war: ``There is no reason,'' said he, ``why you should thank me, if I preferred to be despoiled rather than sold.'' Or as when Asellus reproached Africanus with that unhappy lustrum, he said: ``Don't be surprised; for he who removed you from the rank of \textit{aerarii} closed the lustrum and sacrificed the bull.'' There is a silent suspicion that Mummius bound the state by religion, in that he eased Asellus from disgrace. There is also urbane dissimulation, when other things are said than what you feel — not in that kind of which I spoke before, when you say contraries, as Crassus to Lamia, but when in your whole kind of speech you sport with severity, when you feel one thing and speak another;

\ciceroSection{269} as our Scaevola, when Septumuleius of Anagnia (for whose handing over of Gaius Gracchus's head gold was weighed in payment) asked that he should take him as his prefect to Asia: ``What do you wish, you mad man?'' said he. ``So great is the multitude of bad citizens that I assure you of this: if you stay at Rome, you will come within a few years to the greatest fortune.''

\ciceroSection{270} Of this kind Fannius in his annals says this Africanus Aemilianus was outstanding, and calls him by the Greek word \textit{eiron}; but, as those who know these things better tell us, Socrates I think in this irony and dissimulation by far excelled all in charm and humanity. The kind is most elegant and salty with weight, and apt both for oratorical pleadings and for urbane conversations.

\ciceroSection{271} And by Hercules all these things which I am arguing about wit are not greater seasonings of forensic actions than of all conversations. For just as that saying in Cato — who reported many things, of which I set down some by way of example — seems to me very well noted: that C. Publicius was wont to say P. Mummius was a man for any time. So in fact the matter stands: that there is no time of life in which it is not becoming for charm and humanity to move.

\ciceroSection{272} But I come back to the rest. Akin to this dissimulation is when a faulty thing is called by an honourable word: as when Africanus the censor was moving from his tribe a certain centurion who had not been at the battle of Paulus, and the man said he had remained in camp for the sake of the watch and asked why he was being marked by him: ``I do not love,'' said he,

\ciceroSection{273} ``the over-diligent.'' Keen too is when from another's speech you take a meaning different from his. As Maximus to Salinator: when Tarentum had been lost and yet Livius had kept the citadel and from it had made many distinguished engagements, and some years later Maximus had recovered that town and Salinator asked him to remember that by his work Tarentum had been recovered: ``Why,'' said he, ``should I not remember? For I should never have

\ciceroSection{274} recovered it, if you had not lost it.'' There are also those somewhat absurd things, but for that very reason often laughable, very apt not only for mimes but in some way also for us: ``A foolish fellow, after he had begun to have property, died.'' ``And what is that woman of yours?'' ``A wife.'' ``Like, by Hercules.'' ``And while he was at the waters, he never died.'' This kind is lighter and, as I said, mimic, but has sometimes a place even with us, that even one not foolish, as if foolishly, may say something with salt: as Mancia to you, Antonius, when he had heard you, the censor, were summoned by M. Duronius for bribery,

\ciceroSection{275} ``Now at last,'' said he, ``you may handle your own business.'' Such things are greatly laughed at, and by Hercules everything that is said by prudent men with the dissimulation, as it were, of not understanding, somewhat absurdly and with salt. Of which kind is also not seeming to understand what you do understand: as Pontidius — ``What do you think of one who is caught in adultery?'' — ``Slow!'' Or I myself, when I would not accept Metellus's excuse for my eyes in his levy

\ciceroSection{276} and he had said: ``You then see nothing?'' ``Yes I do,'' said I; ``from the Esquiline gate I see your villa.'' Or that of Nasica, who, when he had come to the poet Ennius and was asking for him at the door, was told by the maid he was not at home — Nasica felt she had said it by her master's order, and that he was inside. A few days afterwards, when Ennius came to Nasica's and was asking for him at the door, Nasica shouted out that he was not at home. Ennius said: ``What? Do I not know your voice?'' Nasica: ``You are an impudent man: when I was looking for you, I believed your maid that you were not at home; you do not believe me myself?'' Pretty too is that, from which the man who has spoken is laughed at in the very kind in which he has spoken;

\ciceroSection{277} as when Q. Opimius the consular, who as a young man had been ill spoken of, said to the festive Egilius — who seemed somewhat soft but was not — ``What now, my Egilia? When are you coming to me with your distaff and wool?'' ``By Pollux, I do not dare,'' he said; ``for my mother forbade me to go near disreputable women.''

\ciceroSection{278} Salty too are those things which have a hidden suspicion of the laughable; of which kind is that of the Sicilian, when an acquaintance was complaining to him that he said his wife had hanged herself on a fig-tree: ``Please,'' he said, ``give me some shoots of that tree to plant.'' Of the same kind is what Catulus said to a bad orator who, in the epilogue, thought he had stirred pity, and after he had sat down asked Catulus whether he seemed to have moved pity: ``Greatly indeed,'' said he, ``for I think there is no one so hard to whom your speech does not seem worthy

\ciceroSection{279} of pity.'' By Hercules, there are also the angry and as it were grumpy laughable things, not when they are said by a grumpy man (for then it is not salt but nature that is laughed at); in which to me that of Novius seems most salty: ``Why are you weeping, father?'' ``It would be strange if I sang: I have been condemned.'' To this kind is opposed as it were the patient and slow kind of the laughable, as when Cato had been struck by a man who was carrying a chest, and the man said ``Take care,'' Cato asked: ``Are you carrying anything else

\ciceroSection{280} besides the chest?'' Salty too is the rebuke of stupidity, as that Sicilian, to whom the praetor Scipio was giving as patron of the case his guest, a noble man but utterly stupid: ``Please, praetor,'' said he, ``give that patron to my adversary, then give me none.'' Those things move us also which by conjecture are explained far otherwise than they are, but keenly and neatly: as when Scaurus was accusing Rutilius of bribery — when he himself had been made consul and Rutilius had failed — and showed in his accounts the letters A. F. P. R., and said it meant ``Acted, by faith of Publius Rutilius'' (whereas Rutilius himself said ``Made before, reported afterwards''); C. Canius, the Roman knight, when assisting Rufus, exclaimed that neither was meant by those letters. ``What then?'' said Scaurus. ``Aemilius did it, Rutilius is punished.'' Discrepancies are also laughed at:

\ciceroSection{281} ``What is missing in this man, except wealth and virtue?'' Pretty also is the friendly rebuke as of one going wrong: as when Granius rebuked Albius for rejoicing over Scaevola's complete acquittal, when something seemed to have been proved against him from his own tablets, and he did not understand that judgment had been given against his tablets.

\ciceroSection{282} Like to this is also a friendly admonition in giving counsel: as when Granius advised a bad orator who had bruised his voice in speaking, to drink cold honey-wine as soon as he should have come home — ``I shall ruin my voice,'' he said,

\ciceroSection{283} ``if I do that.'' ``Better,'' said Granius, ``than your client.'' Pretty too is when something is said as suited to each: as Scaurus had some envy from this — that he had taken the goods of Phrygio Pompeius, a rich man, without a will, and he was sitting as advocate for the defendant Bestia, when a certain funeral was being led, the prosecutor C. Memmius said: ``See, Scaurus — a dead man is being snatched away. If you can be the possessor.''

\ciceroSection{284} But of all these, nothing is more laughable than what is beyond expectation. Of which there are countless examples — as that of the elder Appius, who, in the senate, when business was on the public lands and the Thorian law, and Lucullus was being attacked by those who said that the public lands were being grazed off by his cattle: ``That is not Lucullus's flock,'' he said, ``you err.'' He seemed to be defending Lucullus. ``I take it to be free; for it grazes wherever it likes.''

\ciceroSection{285} I like also that of the famous Scipio who struck down Tiberius Gracchus: when many reproaches had been thrown against him by M. Flaccus, who put forward P. Mucius as judge: ``I reject him as biased,'' he said. There was a murmur. ``Ah,'' he said, ``conscript fathers, I am not rejecting him as biased against me, but against all.'' From this Crassus, however, comes nothing wittier: when the witness Silus had hurt Piso, by saying that he had heard a thing about him, ``It can happen, Silus,'' said he, ``that he, from whom you say you heard it, said it in anger.'' Silus nodded. ``It can also be that you did not rightly understand him.'' That too he nodded with his whole head, that he might give himself to Crassus. ``It can also happen,'' he said, ``that what you say you heard, you never heard at all.'' This came so beyond expectation that the laughter of all overwhelmed the witness. Of this kind Novius is full, of whom is the familiar joke

\ciceroSection{286} ``\textit{The wise man, if you shall be cold, shall tremble},'' and very many others. Often it is also witty to grant the adversary the very thing he is taking away from you: as Gaius Laelius, when one ill-born said to him that he was unworthy of his ancestors: ``By Hercules,'' said he, ``you are worthy of yours.'' Often too laughable things are said sentenciously, as M. Cincius, on the day he carried his law on gifts and presents, when C. Cento came forward and contemptuously enough asked, ``What are you proposing, Cinciole?'' said: ``So that you may buy, Gaius, if you wish to use anything.''

\ciceroSection{287} Often too one wittily wishes things which cannot happen: as M. Lepidus, when the others were exercising in the Campus and he himself had reclined on the grass, said: ``I should wish this were labouring.'' Salty too is to answer slowly to those asking and as it were inquiring something they would not wish: as the censor Lepidus, when he had taken away M. Antistius of Pyrgi's horse and his friends were shouting and asking what he should answer his father, why his horse was being taken from him when he was the best farmer, the most thrifty, most modest, most frugal: ``I,'' he said,

\ciceroSection{288} ``believe none of those things.'' Other things are gathered by the Greeks — execrations, expressions of wonder, threatenings; but these very things I seem to myself to have distributed into too many kinds. For those which are contained in the rationale and force of a word are mostly fixed and defined, which (as I said before) are wont rather to be praised than laughed at;

\ciceroSection{289} but these which are in the very matter and thought are countless in their parts, few in their kinds. For laughter is moved by the deceiving of expectations, by deriding others' natures, by indicating one's own laughably, by the likeness of something more shameful, by dissimulation, by saying somewhat absurd things, and by reproving foolish ones. Therefore he who would speak jocosely should be steeped, as it were, in nature suited for these kinds and in habits — that to each kind of laughable his very face be also fitted. The more severe and gloomier this is — as in you, Crassus — the saltier the things said are wont to seem.

\ciceroSection{290} But now you, Antonius — who said you would gladly rest in this wayside inn of my speech, as though you had turned aside into the Pomptine territory, neither pleasant nor healthful place — I think you should reckon you have rested long enough, and continue the rest of your journey.'' ``I have indeed,'' he said, ``and most cheerfully entertained by you, both more learned through you, and now also bolder for jesting. For I do not fear that anyone will think me lighter in this kind, since you have brought forward to me as my authorities Fabricii, Africani, Maximi, Catones, Lepidi.

\ciceroSection{291} But you have what you wish to hear from me, of which indeed I had to speak more carefully and to think; for the rest are easier, and from what has been said all the rest are born. For when I have approached the case and pursued everything in thought as far as I could; when I have seen and recognised both the arguments of the case, and those topics by which the minds of the judges are conciliated, and those by which they are moved — then I lay down what each case has of good, what of bad. For there is hardly any matter that can be brought into the speaker's dispute or controversy that does not have both, but how much of each, this matters.

\ciceroSection{292} My own method in speaking is wont to be this: that what good I have I embrace, ornate, exaggerate; there I stop, there I dwell, there I cling. From a bad and faulty thing of the case I so withdraw — not so that I should appear to flee it, but so that, by ornating and amplifying that good thing, the whole may be buried hidden. And if the case lies in arguments, I most strongly maintain whichever are firmest, whether they are several or some one; but if the case is in conciliation or in moving, I bring myself most to that part which most can move men's minds.

\ciceroSection{293} The summary of this kind, in fine, is this: that if my speech in refuting the adversary can be firmer than in establishing my own things, I bring all my weapons against him; if our things can be more easily approved than his refuted, I try to draw their minds away from the contrary defence and to lead them to ours.

\ciceroSection{294} Two things finally — which seem easiest, since I cannot do those that are harder — I claim for myself by my own right. One: that to a troublesome or difficult argument or topic I sometimes answer nothing at all — at which perhaps someone will rightly laugh; for who is there who cannot do that? But I am now arguing about my own faculty, not about others'; and I confess that, if anything presses too vehemently, I am wont so to yield that I seem to flee not only with my shield not cast away but not even shifted to the back; but to apply in speaking a kind of show and pomp like a battle's flight; and to take stand in my place of refuge so that I seem to have given way not for fleeing the enemy but for taking up a position.

\ciceroSection{295} The second is what I think most to be guarded and provided against by the orator and what is wont to vex me most: I am wont to labour not so much that I may help the case as that I may not hurt it. Not that one should not strive in each, but it is much more disgraceful for the orator to seem to have hurt the case than to have not helped it. But what is this, Catulus, you and your neighbours together? Do you, as is fit, despise this matter?'' ``By no means,'' said he; ``but Caesar seemed to wish to say something on this very point.'' ``Most willingly,'' said Antonius, ``let him say it, either to refute or to inquire.''

\ciceroSection{296} Then Iulius: ``By Hercules, Antonius,'' he said, ``I have always been one to say of you as orator that you, of all in speaking, seem to me the most guarded; and that this very thing is the proper of your praise — that nothing has ever been said by you that hurt the man for whom you spoke. And I remember, when conversation had been begun by Crassus and me, with many listening, and Crassus was praising your eloquence with many words, that I said, with the rest of your praises, that this was the greatest: that you both said what was needed and did not say what was not needed.

\ciceroSection{297} Then I remember he answered me that other things were highly to be praised in you, but that this — to say what was alien and would hurt the man for whom each man was speaking — was the part of a wicked and treacherous man; therefore he who did not do this seemed to him not eloquent but only an upright man, while he who did this — wicked. Now if you please, Antonius, I would have you show why you think this so great a thing — to do nothing harmful in the case — that nothing seems to you greater in the orator.''

\ciceroSection{298} ``I shall say, Caesar,'' he said, ``what I understand. But you and you all,'' he said, ``remember this: that I am not speaking of the divinity of the perfect orator, but of the moderate measure of my own exercise and habit. Crassus's answer is of an outstanding and singular talent, to whom indeed it has seemed something prodigious that any orator could be found who, by speaking, could do something hurtful and harm him whom he was defending.

\ciceroSection{299} For he conjectures from himself, whose force of talent is so great that he thinks no one says anything against himself save knowingly. But I am now arguing not about an outstanding and exceptional, but about a well-nigh common and shared force. So among the Greeks it is reported that the famous Athenian Themistocles was of incredible greatness of counsel and talent. To him a learned man and chief among scholars is said to have come up, and to have promised that he would hand down to him the art of memory, which was then for the first time being brought forward; when Themistocles asked what that art could effect, the teacher said, ``so that you may remember everything''; and Themistocles answered that the man would do him a greater favour if he taught him to forget what he wished, than if he taught him to remember.

\ciceroSection{300} Do you see what force in a man of the keenest talent, what kind and how great a mind there was? Who so answered that we may understand nothing from his soul, once it had been poured in, ever could have flowed out — when indeed it had been more desirable to him to be able to forget what he wished not to remember, than to remember what he had once heard or seen. But neither, on account of this answer of Themistocles, must we not give pains to memory, nor on account of Crassus's outstanding prudence is that caution and timidity of mine in cases to be neglected. For neither of those men brought any faculty for me, but only signified his own.

\ciceroSection{301} For very many things in cases must be looked round at in every part of the speech, lest you should offend or strike on something. Often some witness either does no harm, or less harm, unless he is challenged. The defendant prays, the advocates press that we should attack, that we should abuse, finally that we should question; I am not moved, I do not obey, I do not satisfy them. Nor yet do I gain any praise; for the unskilled can more easily reproach what you have foolishly said than praise what you have wisely been silent about.

\ciceroSection{302} How much harm is here, if you have hurt an angry, not foolish, not light witness! For he has both will to harm in his anger, and force in his talent, and weight in his life. Nor, even though Crassus does not commit this, do many therefore not commit it often. Than which nothing is wont to seem to me more disgraceful, than that out of the orator's saying or answer or question this conversation follows: ``He has killed.'' ``The adversary?'' ``On the contrary,'' they say, ``himself and the man whom

\ciceroSection{303} he is defending.'' Crassus does not think this can happen save by perfidy. But I most often see in cases men least bad doing something bad. What of that which I said before, that I am wont to give way and (to speak more plainly) to flee those things which most heavily press my case — when others do not do this and dwell in the camp of the enemy and dismiss their own garrisons, do they hurt their cases moderately when they either confirm the adversaries' supports or open up things they cannot heal?

\ciceroSection{304} What if they have no regard for the persons whom they defend — if those things which are odious in them they do not soften by lessening, but make more odious by praising and lifting them up; how much harm is in that? What if you should inveigh more bitterly and contumeliously against men dear and pleasing to the judges, with no preliminary fortification of speech: would you not estrange the judges from yourself?

\ciceroSection{305} What if those vices or disadvantages which are in some one or several judges, you, in reproaching them in adversaries, do not understand that you are inveighing against the judges — is that a slight fault? What if, when you speak for another, you make the case your own, or, hurt, are carried away by anger and abandon the case, do you do no harm? In which I, not because I gladly hear myself ill spoken of, but because I do not gladly leave the case, am thought too patient and slow — as when I rebuked you yourself, Sulpicius, that you were attacking the assistant, not the adversary. From this also I gain that, if anyone speaks ill of me, he seems either insolent or plainly insane.

\ciceroSection{306} In the arguments themselves, if you have laid down anything either openly false, or contrary to what you have said or are about to say, or by its very kind removed from the practice of trials and the Forum — do you do no harm? In short: all my care is wont to be in this — for I shall say it again — that, if I can, I should accomplish something good by speaking; if not, that at least I should do nothing bad.

\ciceroSection{307} So now I come back, Catulus, to that on which you praised me a little before — to the order and arrangement of subjects and topics. Of which there is a twofold rationale: one which the nature of cases brings; the other which is gathered by the orator's judgment and prudence. For that we should say something before the matter, then set the matter forth, afterwards prove it by confirming our supports and refuting contraries, then conclude and so peroration — this the very nature of speaking prescribes.

\ciceroSection{308} But how we should arrange those things which must be said for proving and teaching, that is most especially proper to the orator's prudence. For many arguments occur, many which seem profitable in speaking; but of these some are so light that they should be despised; some, even if they have any help, are sometimes such that there is some fault in them, and that which seems to profit is not worth being joined with some bad thing.

\ciceroSection{309} Of those, however, which are useful and firm — if they are yet, as often happens, very many — those which are either lightest or like others heavier, I think should be set apart and removed from the speech. For my part, when I gather the arguments of cases, I am wont not so much to count them as to weigh them.

\ciceroSection{310} And since, as I have often said, by three things we lead men to our opinion — either by teaching, or conciliating, or moving — one of these three things is to be put in front of us, that we may seem to wish nothing else save to teach. The other two, like blood in bodies, must be diffused through continuous speeches. For both the openings and the rest of the parts of speech, of which we shall presently say a few things, ought greatly to have this force: that they may be able to bear

\ciceroSection{311} on the moving of the minds before whom one is to plead. But to those parts of speech which, even if they teach nothing by arguing, yet by persuading and moving accomplish very much — although the chief and proper place is in the opening and in the peroration — to digress from what you have proposed and are pleading is yet often useful, for the moving of minds.

\ciceroSection{312} So either when the matter has been narrated and set forth, often a place is given for digressing for moving minds, or after our arguments have been confirmed, or contraries refuted, or in each place, or in all of them, if the case has the dignity and abundance, this can rightly be done. And those cases are weightiest and fullest for amplifying and adorning, which give the most outlets for such digressions, that we may use those topics by which the impulses of the minds of those who hear are either driven on or turned back.

\ciceroSection{313} I also rebuke those who put first what is least firm; in which I think those err too who, when they bring in several patrons (which has never pleased me), in proportion as they think there is least power in any of them, so they wish him to speak first. The matter requires this — that the expectation of those who hear should be met as quickly as possible. If at the beginning satisfaction has not been given to it, much more must be laboured at in the rest of the case. For a case is in a bad way which does not seem at once, as it begins to be spoken, to become better.

\ciceroSection{314} So as in the orator each best one, so in the speech let each firmest matter be first; provided in each it is held that what is excellent is also kept for the peroration. If there are middling things (for the faulty ought to have no place anywhere), let them be cast into the middle throng and herd.

\ciceroSection{315} All these things being considered, then at last what is to be said first, I am wont to think of last: what opening to use. For if I have ever wished to find this first, none has come to me save either meagre or trifling or vulgar or common. The openings of speaking ought always to be careful and keen and equipped with thoughts, fitting in words, and indeed proper to the cases. For first, as it were, comes the recognition and commendation of the speech in the opening, which ought immediately to soothe and entice him who hears.

\ciceroSection{316} In which I am wont to wonder not at those who have given no pains to this matter, but at a man chiefly eloquent and learned, Philippus, who is wont so to rise to speak, that he does not know what first word he is to have. He even says that, when he has warmed up his arm, he is then wont to fight; and he does not consider that those very ones from whom he draws this likeness throw their first spears so gently as to serve grace very greatly and to spare their remaining strength.

\ciceroSection{317} Nor is there doubt that the opening of speaking ought not often to be vehement and pugnacious; but if even in that very gladiatorial contest of life — by sword decided — many things happen before the engagement that seem to avail not for wounding but for show, how much more must this be looked at in speaking, in which not so much force as delight is required! Lastly, there is nothing in the nature of all things that pours itself out wholly and unrolls itself entirely on a sudden. Thus all things which happen and which are most keenly performed, nature herself has fringed about with gentler beginnings.

\ciceroSection{318} These things in speaking are not to be sought from outside somewhere, but to be taken from the very inwards of the case. So when the whole case has been thoroughly tested and seen through, and all the topics found and equipped, one must consider what opening to use.

\ciceroSection{319} Thus they will be both easily found — they will be taken from those things which will be most copious either in the arguments or in those parts to which I said one ought often to digress. So they will bring some weight, since they will have been drawn nearly from the inmost defence; and it will be plain that they are not common nor capable of being transferred to other cases, but that they have flowered out from deep in the case being then pleaded.

\ciceroSection{320} Every opening will have to have either a signification of the whole matter to be pleaded, or an entrance to the case and a way of being fortified, or some kind of ornament and dignity. But as we set vestibules and approaches before houses and temples, so we must set openings before cases in proportion to their size. So in small and infrequent cases it is often more advantageous to begin from the matter itself.

\ciceroSection{321} But when an opening must be used, as for the most part it must, thoughts may be drawn either from the defendant or from the adversary, or from the matter, or from those before whom one is to plead. From the defendant (I call defendants those whose case is at stake): things which signify a good man, generous, calamity-stricken, worthy of pity; things which avail against false accusation; from the adversary, contraries from the same topics in the main;

\ciceroSection{322} from the matter — if cruel, if wicked, if unexpected, if undeserved, if pitiful, if ungrateful, if unworthy, if new, if such cannot be restored or healed; from those before whom one is to plead — that we should make them favourable and well-thinking, which is more readily brought about by acting than by asking. This indeed must be diffused through the whole speech and not least at the end. But yet many openings are born from this kind.

\ciceroSection{323} For the Greeks teach that we should, in the opening, make the judge attentive and teachable. Which are useful — but no more proper to the opening than to the rest of the parts. They are easier in the openings — both because men are then most attentive when they expect everything, and they can more readily be made teachable in the beginnings; for what is said in the openings is brighter than what is said in the middles of cases by way of arguing or refuting.

\ciceroSection{324} The greatest abundance of openings, for either luring or rousing the judge, we shall draw from those topics for moving the soul which will lie in the case. These yet should not be wholly unfolded in the opening, but the judge should only be lightly impelled at first, so that the rest of the speech may rush upon him already inclined.

\ciceroSection{325} But the opening should be so connected with the following speech that it may seem not, as a citharist's prelude, something tacked on, but a member cohering with the whole body. For some, when they have produced their meditated piece, so pass to the rest as if they did not wish a hearing for themselves. And that prelude must be of such a kind, not as that of the Samnites, who brandish their spears before the fight, with which in the fighting they make no use, but with the very thoughts at which they have prelude they may also be able to fight.

\ciceroSection{326} As for narrating the matter — that they bid be brief, if brevity must be called the case where no word overflows: brief is L. Crassus's speech. But if brevity is then, when there is just so much of words as is needed, sometimes it is needed; but often it is harmful, especially in narrating, not only because it brings obscurity, but because it removes that virtue which is the greatest of narration: that it be pleasing and apt for persuading. Look at Terence's: ``\textit{For after he came forth from boyhood} . . .''

\ciceroSection{327} How long the narration is! The morals of the youth himself, the slave's questioning, the death of Chrysis, the look and form and lamenting of his sister — the rest is most variously and pleasantly told. But if he had sought this brevity: ``She is carried out, we go, we come to the tomb, she is laid on the fire'' — he could have got the whole done in some ten short verses. Although that very ``she is carried out, we go'' is so clipped that it is not brevity that has been served but rather grace.

\ciceroSection{328} For if there had been nothing but ``she is laid on the fire,'' yet the whole matter could easily have been known. But narration distinguished by persons and punctuated by speeches has merriment; and what you say has been done is more probable, when you set out the way it has been done; and it is much plainer to understand, if it is constituted at some length and not run over with that brevity.

\ciceroSection{329} For narration must be as plain as the rest. But this must be more laboured at, because it is harder not to be obscure in narrating a matter than either in the opening or in arguing or in perorating. With even greater danger this part of speech is obscure than the rest: either because, if anything has been said somewhat obscurely in some other place, only that perishes which has been so said, while an obscure narration darkens the whole speech; or because other things, if you have said them somewhat obscurely once, you can say more plainly in another place, but the narration has only one place in the case. The narration will be perspicuous if usual words are used, if the order of times is preserved, if it is told without interruption.

\ciceroSection{330} But when narration must be used or not, that is of judgment. For neither, if the matter is known and there is no doubt what has been done, must one narrate; nor, if the adversary has narrated, unless we are to refute. And if it is to be narrated sometimes, neither shall we keenly pursue those things which would create suspicion and charge and be against us, and we shall draw off whatever we can; that what Crassus thinks, if it ever happens, happens by perfidy, not by stupidity — that we hurt our case — may not happen. For it pertains to the summary of the whole case whether the matter has been set out cautiously or the contrary; for the narration is the spring of all the rest of the speech.

\ciceroSection{331} Next, the case must be set down, in which we must see what comes into controversy. Then the supports of the case must be supplied jointly, both for refuting the contraries and for confirming yours. For one method in cases is of that speech which avails for proving the argument; and that asks both for confirmation and for reprehension. But because neither can what is said against be reproved unless you confirm yours, nor can yours be confirmed unless you reprove theirs, on that account these are joined both in nature and in usefulness and in handling.

\ciceroSection{332} All things must mostly be concluded with amplifying matters either by inflaming the judge or by softening him. And all things, both with the previous topics of speech and especially in the last, must be brought together for moving the minds of the judges as much as possible and calling them to our advantage.

\ciceroSection{333} Nor does there now seem to be any reason why we should set apart those precepts which are to be handed down on hortations or laudations: for most are common. Yet still to advise something or dissuade something seems to me of the weightiest character. For it belongs to a wise and to an honourable and to an eloquent man to set out his counsel on the greatest matters — that you may foresee in your mind, prove by authority, persuade by speech. And these things must be carried on in the senate with less display: for it is a wise body, and place must be left for many others to speak; even the suspicion of showing off talent must be avoided.

\ciceroSection{334} A public meeting holds all the force of speech and asks for weight and variety. So in advising nothing is more to be desired than dignity. For he who pursues utility looks not at what the persuader most wishes, but at what is sometimes more attainable. There is no one, especially in so renowned a state, who does not think dignity above all is to be sought; but utility for the most part wins, when there is the underlying fear that, if it be neglected, even dignity cannot be retained.

\ciceroSection{335} The controversy among men's opinions is either on which is more useful, or, even when it is agreed which is, the contention is whether thought should rather be had for honour or for utility. These often seem to fight together. He who shall defend utility will count up the advantages of peace, of resources, of power, of revenues, of military protection, of the rest by which we measure fruit by utility, and likewise the disadvantages of the contraries. He who urges to dignity will gather the examples of the ancestors, which were glorious even with danger; he will increase the immortal memory of posterity; he will defend that utility is born of praise and is always joined with dignity.

\ciceroSection{336} But what can or cannot be done, and what is even necessary or not, in either thing must most be inquired. For all deliberation is cut off if it is understood that something cannot be done, or if necessity is brought; and he who has shown this when others did not see it has seen most.

\ciceroSection{337} For giving counsel about the commonwealth the chief thing is to know the commonwealth; for speaking probably, to know the morals of the state, which, because they often change, the kind of speech must often be changed too. And although the force of eloquence is one in the main, yet because the highest dignity is of the people, the weightiest case is of the commonwealth, the greatest movements are of the multitude, the kind of speech also must seem to be applied as somewhat grander and brighter; and the greatest part of the speech must be applied to the movements of minds — sometimes by exhortation or by some recalling, to be roused either to hope or to fear or to desire or to glory; often also to be called back from rashness, anger, hope, injury, envy, cruelty.

\ciceroSection{338} It happens that, because the assembly seems to be the orator's greatest stage, we are roused by nature herself to a more ornate kind of speaking. For the multitude has a kind of force such that, just as a piper without his pipes, so an orator without a multitude listening cannot be eloquent.

\ciceroSection{339} And since there are many and various lapses among the people, the contrary acclamation of the people is to be avoided. This is roused either by some fault of speech — if anything seems to have been said roughly, arrogantly, basely, sordidly, by some fault of mind — or by an offence at the men or by envy (which is either just, or comes from accusation and rumour); or if the matter is displeasing; or if the multitude is in some movement of its own desire or fear. To these four causes are opposed as many medicines: rebuke, if there is authority; admonition, as a gentler rebuke; promise that, if they will hear, they will approve; entreaty, which is feeble, but sometimes useful.

\ciceroSection{340} In no place do witticisms profit more, and swiftness, and some brief saying not without dignity and with charm. For nothing is so easy as for the multitude to be drawn from gloom — and often from bitterness — by something said suitably and briefly and keenly and cheerfully. I have set out for you, as I could, in each kind of cases what I am wont to follow, what to flee, what to look at, and on what method altogether to move in cases.

\ciceroSection{341} Nor is that third kind of laudation hard, which I had at first set apart, as it were, from our precepts. Both because there are many kinds of speeches both weightier and of greater abundance about which hardly anyone would give precept, and because we are not wont to use laudations very much, I was setting this whole topic aside. For the Greeks themselves wrote laudations more for the sake of reading and pleasure or of adorning some man, than for the sake of this forensic utility; whose books are extant, in which Themistocles, Aristides, Agesilaus, Epaminondas, Philippus, Alexander, and others are praised. Our laudations, which we use in the Forum, either have the bare and unornamented brevity of testimony or are written for a funeral assembly, which is least apt for the praise of speech. But still — since they must sometimes be used and sometimes even written (as Gaius Laelius wrote for Quintus Tubero when he was praising his uncle Africanus, or so that we ourselves, for ornament, in the Greek manner, may be able to praise whom we wish) — let this topic too be handled by us.

\ciceroSection{342} It is plain, then, that some things in a man are to be desired, others praised. Birth, beauty, strength, wealth, riches, and the rest which fortune gives, either externally or to the body, do not have in themselves true praise, which is thought owed only to virtue; but yet, because virtue itself is most discerned in the use and moderation of those things, these goods of nature and fortune must also be handled in laudations. In which the highest praise is: not to have lifted oneself up in power, not to have been insolent in money, not to have set oneself before others on account of the abundance of fortune, so that resources and stores may seem to have given the means and matter not to pride and lust but to goodness and moderation.

\ciceroSection{343} Virtue itself, which is praiseworthy of itself and without which nothing can be praised, has yet several parts, of which one is more apt than another for laudation. For there are some virtues which are seen as set in the morals of men and in a certain kindness and beneficence; others which are in some faculty of talent or in greatness and strength of soul. For clemency, justice, kindness, faith, courage in common dangers are pleasant to hear in laudations;

\ciceroSection{344} all these virtues are reckoned profitable not so much to those who have them, as to the human race. Wisdom, and greatness of soul (by which all human things are reckoned slight and as nothing), and a certain force of talent in inventing, and even eloquence have admiration no less, but pleasantness less. For these seem rather to adorn and protect those whom we praise, than those before whom we praise. But yet in praising these kinds of virtues too must be joined; for the ears of men bear that those things which are pleasing and pleasant, and also those which are wonderful in virtue, should be praised.

\ciceroSection{345} And since each single virtue has certain offices and duties, and to each its own praise is owed, it must be set out in the praise of justice what the man being praised did with faith, with equity, with such a kind of duty. Likewise in the rest, the deeds done will be fitted to the kind, force, and name of each virtue.

\ciceroSection{346} The most pleasing praise is held to be of those deeds which seem to have been undertaken by brave men without profit and reward; but those which were also with their own labour and danger, these have the richest abundance for praise: both because they can be most ornately said, and most easily heard. For that finally seems the virtue of an outstanding man, which is fruitful to others, but to himself either toilsome or dangerous or at any rate gratis. Great too is wont to seem that praise — to have borne adverse chances wisely, not to have been broken by fortune, to have kept dignity in harsh circumstances.

\ciceroSection{347} Nor yet do they not adorn — honours had, decreed rewards of virtue, deeds approved by men's judgments. In which it is part of laudation that even good fortune itself is assigned to the judgment of the immortal gods. The matters to be taken should be either outstanding for greatness, or first by novelty, or singular in their very kind. For neither small nor common nor vulgar things are wont to seem worthy of admiration or even of praise.

\ciceroSection{348} There is also in laudation a brilliant comparison with other outstanding men. About this kind I have wished to say a little more than I had shown — not so much for forensic use (which has been treated by me throughout this whole conversation), as that you might see this: that, if laudations are in the orator's office (which no one denies), the knowledge of all the virtues, without which laudation cannot be effected, is necessary to the orator.

\ciceroSection{349} It is now plain that the precepts of blaming must be drawn from the contrary vices. At the same time, this is before our eyes: that neither can a good man be praised properly and copiously without knowledge of the virtues, nor an evil one be marked and blamed sufficiently distinctly and harshly without knowledge of the vices. And these topics, both of praising and of blaming, must often be used by us in every kind of cases.

\ciceroSection{350} You have what I think about discovering and arranging matters. I shall add also about memory, that I may relieve Crassus of labour and leave nothing for him to discuss except those things by which these are adorned.'' ``Go on, then,'' said Crassus, ``for gladly I see you, the master once recognised and now also unrolled from those wrappings of your dissimulation, laid bare. And in that you leave me nothing or not much, you do most commodiously, and it is welcome to me.''

\ciceroSection{351} ``How much I shall have left to you,'' said Antonius, ``will be in your power. For if you wish to act truly, I leave everything to you; but if to dissemble, you must see how to satisfy these. But to come back to the matter, I am not,'' he said, ``of so great a talent as Themistocles was, that I should prefer the art of forgetting to that of memory; and I owe thanks to that famous Simonides of Ceos, who, they say, first brought forward the art of memory.

\ciceroSection{352} They say that, when Simonides was dining at Crannon in Thessaly with Scopas, a fortunate and noble man, and had sung the song which he had written about him, in which much had been written, in the way of poets, in praise of Castor and Pollux for ornament — Scopas in too sordid a manner told Simonides he would give him half of what they had agreed for that song. The rest he should seek from his Tyndaridae, whom he had praised equally, if it pleased him.

\ciceroSection{353} A little later they say it was reported to Simonides to come out: that two young men were standing at the door who were urgently calling him out. He rose, came out, saw no one. In this interval the room in which Scopas was banqueting fell down. By that ruin Scopas himself, with his kinsfolk, was crushed and killed. When his people wished to bury them and could not in any way distinguish them as crushed, Simonides is said to have been the demonstrator for burying each — from this, that he had remembered in what place each of them had reclined. Warned by this matter, he is reported to have discovered that order is what most brings light to memory.

\ciceroSection{354} So those who would exercise this part of talent must take places, and shape in their minds the things they wish to keep in memory and place them in those places. Thus it would happen that the order of the places kept the order of the things, and the things themselves were marked by the figures of the things; and we should use the places as wax, the images as letters.

\ciceroSection{355} What use of memory, what utility, what force is the orator's, why need I say? To hold what you have learned in the receiving of the case, what you have yourself thought; that all your thoughts be fixed in the soul, all the equipment of words written out; so to listen either to him from whom you learn, or to him to whom you must answer, that they may seem to you not to pour their speech into your ears, but to inscribe it on your soul. So those alone who have a vigorous memory know what and how far and how they shall speak, what they have answered, what is left. The same men also remember many things from other cases now and then handled by themselves, much they have heard from others.

\ciceroSection{356} I confess therefore that the nature of this good is principal, as of all those things of which I spoke before. But this whole art of speaking — or whatever image and likeness of art it is — has this force: not that it should beget and produce some whole of which there is no part in our talents, but that those things which have already arisen and been produced in us, it should bring out and confirm.

\ciceroSection{357} Yet hardly any one is of so keen a memory that, with the matters not arranged and noted, he could embrace the order of all words or thoughts; nor again of so dull a memory as not to be helped by this habit and exercise. For he saw this prudently, whether Simonides or some other invented it: that those things are most stamped on our minds which are handed down and impressed by the senses; and that the keenest of all our senses is that of seeing. So things which are perceived by the ears or by thought can most easily be kept in mind, if they are also handed to our minds by the recommendation of the eyes — that things blind and removed from the judgment of sight a certain shape and image and figure should mark, so that what we could scarcely embrace by thinking, we should as it were hold by gazing upon.

\ciceroSection{358} By these forms and bodies, just as by all things which come under sight, our memory is reminded and roused; but a seat is needed, since a body cannot be understood without a place. Therefore, that I may not be too long and unfamiliar in a known and well-worn matter — many places must be used, illustrious, set out, with moderate intervals; images, however, that are active, keen, marked, which can quickly come and strike the soul. This faculty both exercise will give (from which habit is born), and the noting of like words turned and changed by cases or transferred from a part to a kind, and the shaping of a whole thought by the image of one word — by the rationale and method of some chief painter who distinguishes places by variety of forms.

\ciceroSection{359} The memory of words, however, which is less necessary to us, is distinguished by greater variety of images. For there are many words which, like joints, connect the limbs of speech, which can be shaped by no likeness; for them we must invent images for ourselves, which we should always use. The memory of matters is the orator's own; this we can mark with single persons well placed, that we should grasp the thoughts by images, the order by places.

\ciceroSection{360} Nor is it true, what is said by the slothful, that memory is overwhelmed by the weight of images, and that even what nature could of itself have held is obscured. For I have seen the foremost men with well-nigh divine memory: at Athens Charmadas, in Asia (whom they say is alive today) Metrodorus of Scepsis. Each of these said that, just as one writes letters in wax, so he wrote out by images in those places he had what he wished to remember. Therefore by this exercise memory is not to be dug out, if there is none in nature; but surely, if it lurks, it must be called forth.

\ciceroSection{361} You have a man's tolerably long discourse — would that not an impudent one! That, certainly, of one not too modest; who, with you, Catulus, listening, with L. Crassus also listening, has said so much about the rationale of speaking. For the age of those over there should perhaps have moved me less. But you will surely pardon me, if you will hear what now drove me to this loquacity unusual to me.''

\ciceroSection{362} ``We do indeed,'' said Catulus, ``— for I answer this both for myself and for my brother — not only pardon you, but love you and owe you great thanks. And while we recognise your kindness and easy nature, we wonder at this knowledge and abundance. For my part I think I have gained this too, that I am freed from a great error and from that wonder which, with many others, I was always wont to feel — whence came that great divinity of yours in cases. For I did not think you had touched these things, which I now see you have most diligently learned and gathered from every side, and, taught by use, have in part corrected, in part approved.

\ciceroSection{363} Nor on that account do I admire your eloquence less, but your virtue and diligence much more. And at the same time I rejoice that the judgment of my mind is approved, by which I have always laid down that no one can attain praise of wisdom and eloquence without the highest zeal and labour and learning. But yet what is that thing you said would happen — that we should pardon you, if we knew what cause had driven you to the conversation? For what other cause is there but that you have wished to humour us and the eagerness of these young men, who have heard you most attentively?''

\ciceroSection{364} Then he: ``I wished,'' he said, ``to take away from Crassus all refusal, whom I knew a little before to come to this kind of conversation either too modestly or too unwillingly (I do not wish to say more fastidiously of so sweet a man). For what could he say? That he is a consular and a censorial man? Our case is the same. Will he plead his age? He is four years younger. That he does not know these things — which I, late, snatched up in haste, in spare hours, as they say, while he from boyhood with the highest zeal, with the highest teachers? I shall say nothing of his talent, to which no one was equal. For when I was speaking, no one ever was so contemptuous of himself that he did not hope to speak either better or in the same way; when Crassus speaks, no one is so arrogant as to trust that he would ever speak similarly. Therefore, lest these such men have come in vain, let us at last hear you, Crassus.''

\ciceroSection{365} Then he: ``Granting that those things are so,'' he said, ``Antonius — which are very much otherwise — what at last have you left to me today, or to any man, that could be said? For I shall speak truly, my dearest friends, what I feel: I have often heard learned men. Why do I say often? Rather, sometimes; for how could I do so often, who came as a boy to the Forum and never was away thence longer than my quaestorship? But yet I heard, as I was saying yesterday, when I was at Athens, the most learned men, and in Asia that very Metrodorus of Scepsis, when he was disputing about these very matters; nor has anyone seemed to me more copious or more subtle in this kind of speaking than this man has been today. But if it were otherwise and I should understand something passed over by Antonius, I should not be so uncivil and almost inhuman as to refuse what I felt you wished.''

\ciceroSection{366} Then Sulpicius: ``Have you forgotten then, Crassus,'' he said, ``that Antonius so divided with you: that he himself should set out the orator's equipment, and leave to you its distinguishing and ornament?'' Here he: ``First, who allowed Antonius,'' he said, ``both to make the parts and to take whichever he wished himself first? Then, if I have rightly understood, since I was hearing very gladly, he seemed to me to have spoken jointly on each.'' ``He, however,'' said Cotta, ``did not touch the ornaments of speech, nor that praise from which eloquence has its name.'' ``Antonius then,'' said Crassus, ``has left me the words; he himself has taken the matter.''

\ciceroSection{367} Then Caesar: ``If he has left to you what is harder,'' he said, ``we have a reason why we should desire to hear you. But if what is easier, you have no reason to refuse.'' And Catulus: ``What of what you said, Crassus,'' he said, ``that, if we should remain with you today, you would humour us — do you think this is nothing to your faith?'' Then Cotta, smiling: ``I could yield to you, Crassus,'' he said, ``but see lest Catulus has brought any religious bond. This is censorial business: how that fits a censorial man, you see.'' ``Come now,'' he said, ``as you wish. But now indeed, since the time is what it is, I think we should rise and rest. After noon, if so it is convenient to you, we shall speak about something — unless perhaps you prefer to put it off to tomorrow.'' All said they wished either at once, or, if he himself preferred, after noon, yet to hear as soon as possible.

\ciceroBook{3}

\ciceroSection{1} Quintus my brother, as I set out to recall to mind and to consign to this third book the conversation which Crassus held after Antonius's discussion, a sharp recollection has renewed my old grief and trouble of soul. For that talent worthy of immortality, that humanity, that virtue of L. Crassus was extinguished by his sudden death scarcely ten days after the day which is contained in this and the previous book.

\ciceroSection{2} For when he had returned to Rome on the last day of the stage games, deeply moved by that speech which was reported to have been made in the assembly by Philippus — who, it was agreed, had said: ``I must look for some other counsel; with that senate I cannot conduct the commonwealth'' — early on the Ides of September both he himself and a crowded senate came at Drusus's summons into the curia. There Drusus, after he had complained much about Philippus, brought the matter before the senate concerning that very thing — that the consul had inveighed so weightily in the assembly against that order.

\ciceroSection{3} Here, as I have often seen agreed among the wisest men, although it had nearly always happened to Crassus, when he had spoken anything more carefully, that he was thought never to have spoken better, yet by the consensus of all it was so judged at that time: that the rest had always been surpassed by Crassus, but on that day even he himself was surpassed by himself. For he lamented the chance and bereavement of the senate, whose patrimony of dignity was being torn away by the consul (who ought to be like a good parent or faithful guardian), as if by some criminal robber. Nor was it to be wondered at, if, when by his own counsels he had ruined the commonwealth, he should reject from the commonwealth the senate's counsel.

\ciceroSection{4} When he had brought, as it were, certain firebrands of words to the man — vehement and eloquent and chief of all in resisting, Philippus — Philippus did not bear it, and burned with grave anger, and, having seized pledges, set about coercing Crassus. In which very place many things are reported to have been said by Crassus divinely, when he denied that man could be his consul, to whom he himself was no senator. ``Or do you, when you have reckoned the whole authority of our entire order as a pledge and have cut it down before the eyes of the Roman people, suppose me to be terrified by these pledges? Not those things must you cut down, if you wish to coerce L. Crassus: this tongue must be cut out by you, by which, even if torn out, my liberty by the very breath shall confound your lust.''

\ciceroSection{5} Many things are said by him on that day, with the highest contention of soul and body, weightily and abundantly, and it was indeed his outstanding statement, his ``swan-song'' — which we, awaiting more from him, used to come, as it were, to his house and tomb after his death, that we might see at least the very place in which he had lastly stood. The pain of his side at once attacked him in saying; sweat broke out, then a chill came on. So, having returned home with a fever, on the seventh day after, he died of pleurisy.

\ciceroSection{6} O the deceitful hope of men, and the delusive fortune, and our vain studies, which, in mid-course, are often broken and dashed and shipwrecked before they could see the harbour! For while Crassus's life had been busy with the labours of canvassing, of duties, of cases, while he was outweighing all by his usefulness — towards the end, even by his authority — that age had then come which both gave him not labour and gave him commendation. As long as some, even of his close friends, did not approve those wonderful deeds of his — that he had withdrawn himself from the courts and ceased to be the leader of cases. Yet to himself, however much he relaxed of forensic labour, the auctorship of public affairs would have been increased. And for me, neither the highest dignity of his consular kind, nor his most diligent and devoted friendship to me, would have given so much pleasure as that great work of his, in the senate, on the day before his death.

\ciceroSection{7} For of his death there are many memorials and many griefs left to me. Mine indeed: for I have lost a man — distinguished, the friendliest of all to me, with whom I lived most intimately, the foremost of all in dignity, in my judgment, and the most distinguished in eloquence; whose work in cases had been most pleasant to me. Of all our citizens too. For he was, as it seemed to all wise men, the brightest light of our state, who could have raised it both by his counsel and by his speech, ablaze in either, when bent upon the help of his brother and to the safety of the commonwealth.

\ciceroSection{8} For when by Drusus's law the senatorial courts had been overthrown, and again the order of senators with the order of equestrians had been brought into contention, there was nothing of his speech in that case more vehement, more weighty, more eloquent than that, than to which I was a frequent hearer in his other cases. Yet, broken in body, with strength shattered, his soul was not bent at all, and he lifted himself up by speaking; and, the rest having departed almost from terror, he poured into the senate, as it were a will, the testimony of his counsel and judgment about the commonwealth — that for him, who was now perishing, it was easy to die.

\ciceroSection{9} I lament that great citizen, and (since I cannot go further) that I have lost so great an ornament of speech and of public counsel. Just so, surely as the most wretched among many other things, I bewail the death of L. Crassus before all. For when did I see him, but with that lessening of pain we feel for friends with whom even our closeness has been lessened? But, by Hercules, when I see his image, when I read his speeches — although the very things he said grow old in my soul — yet a certain solemn shape, even in the writing, is set before me of L. Crassus.

\ciceroSection{10} What, then, can be greater than this — that we should bear, for our orators of these times, that they should be most distinguished and most outstanding (for so they are), and yet a longing, a memory, a desire is yet excited in us, of those who have already gone before? Wherefore, do you not yet judge it to be true, what I am wont to say to you in our daily conversations: that, even now, after the death of L. Crassus, I see no orator, when I have heard him, who can satisfy my hearing — let alone my judgment.

\ciceroSection{11} He did not see the horrible and pitiful chances of those very ones who, then young men, had given themselves to Crassus. Of whom Gaius Cotta, whom he had left flourishing, a few days after Crassus's death, was driven by hatred from the tribunate, and not many months from that time was cast out of the state. Sulpicius, who had been in the same flame of hatred, made it his business in his tribunate to strip those whom, as a private man, he had lived with most closely, of all dignity. Whose life, indeed, blooming towards the highest glory of eloquence, was snatched by the sword, and the punishment of his recklessness was set, not without great evil to the commonwealth.

\ciceroSection{12} I think you, Crassus, both by the flower of your life and by the opportuneness of your death, were both adorned and extinguished by divine counsel. For you would either have had to undergo, according to your manly soul and constancy, the cruelty of the citizens' steel; or, if any fortune had freed you from the atrocity of death, that same would have forced you to be a spectator of the funerals of your country. Nor would only the dominance of the wicked have been a grief to you, but also the victory of the good — because of the slaughter of citizens mixed with it.

\ciceroSection{13} For me, indeed, my brother Quintus, when I think of both the chances of those of whom I spoke before, and what we ourselves have borne and felt for our incredible and singular love for the commonwealth, your opinion often seems to me true and wise — you who, on account of so many, so great, and so headlong falls of the most distinguished and the best of men, have always called me back from all contention and conflict of soul.

\ciceroSection{14} But since these things now cannot be undone for us, and our highest labours are softened by the great glory of compensation, let us press on to those consolations which can be both pleasant for our settled troubles, and salutary for those still pricking us; and let us hand to memory the rest, and well-nigh last, conversation of L. Crassus; and to him, if not equal to his talent, yet according to our zeal, let us pay back the merited and owed thanks.

\ciceroSection{15} For there is none of us who, when he reads Plato's wonderfully written books, in which Socrates is expressed in nearly all of them, although those things are written divinely, does not yet suspect something greater of him about whom they are written. Which we likewise ask — not from you, who give us all the highest things, but from the rest who shall take this into their hands — that they suspect something greater of L. Crassus than what shall be expressed by us.

\ciceroSection{16} For we who were not present at the conversation, and to whom Gaius Cotta had handed down only the topics and thoughts of this disputation, in what kind of speech we had recognised each orator, this very thing in their conversation we have tried to sketch out. But if anyone — led by common opinion — shall think Antonius leaner or Crassus fuller than they are introduced by us, he will be one of those who either has not heard them or cannot judge. For each was, as I set forth before, both outstanding to all in zeal, talent, and learning, and perfect in his own kind, so that this ornament of speech was neither lacking in Antonius nor overflowing in Crassus.

\ciceroSection{17} When, therefore, they had departed before noon and had rested a little, this Cotta said he had especially noticed: that all that midday time Crassus had spent in the keenest and most attentive thought; and he himself, who well knew his face when he had to speak, and the gaze of his eyes when in thought, and had often seen them in the greatest cases — then, of set purpose, while the others were resting, came into the alcove in which Crassus had reclined on a couch. When he felt that Crassus was fixed in thought, he at once withdrew, and in that silence about two hours were spent. Then, when all had come to Crassus, the day now bending to the afternoon, ``What, Crassus,'' said Iulius, ``do we go to sit? Although we have come not to demand it of you, but to remind you.''

\ciceroSection{18} Then Crassus: ``Do you think me so impudent that I should suppose I can owe you this duty especially any longer?'' ``Where then will the place be?'' he said. ``Does the middle of the wood please? For it is the most shady and cool of places.'' ``Indeed,'' said Crassus, ``in that place there is a seat not unsuited to this our conversation.'' When this had pleased the rest also, they came into the wood, and there, with great expectation of hearing, they sit down.

\ciceroSection{19} Then Crassus: ``Both your authority and friendship,'' he said, ``and Antonius's easiness, have snatched from me, in my best case, the liberty of refusing. Although in the dividing of our disputation, when he took for himself those things which the orator ought to say, but left to me to set forth how those things ought to be ornated, he divided what cannot be separated. For since all speech consists of matter and words, neither can the words have a seat if you should subtract the matter, nor can the matter have light if you remove the words.

\ciceroSection{20} And to me, the ancients seem, having grasped something greater in their souls, to have seen far more than the edge of our talents can look upon — they who said all these things which are above and beneath are one, and bound up by one force and by one consensus of nature. For there is no kind of things which, torn from the rest, could either of itself stand, or, if the rest should be without it, could keep its own force and eternity.

\ciceroSection{21} But if this seems too great a rationale to be grasped by men's sense or thought, there is also that true voice of Plato, surely not unheard by you, Catulus: that all the doctrine of these freeborn and human arts is contained by some one chain of fellowship. For when the force of that reasoning is seen by which the causes of things and their outcomes are known, a wonderful concord, as it were, of all branches of knowledge and a harmony is found.

\ciceroSection{22} But if this also seems higher than that we, lying flat on the ground, can look up to it, yet that surely which we have embraced, which we profess, which we have undertaken, we ought to know and hold. For there is one eloquence — what both I said yesterday, and Antonius signified in several places of his morning speech — into whatever shores or regions of disputation it is brought.

\ciceroSection{23} For whether it speaks about the nature of heaven or of earth, of divine or of human force; whether from a lower place or from an equal or from a higher; whether to drive men or to teach or to deter or to rouse or to turn back or to inflame or to soothe; whether to few or to many, whether among foreigners or with one's own or with oneself — speech is led off by streams, not by springs, and wherever it goes, by the same equipment and ornament it is accompanied.

\ciceroSection{24} But since we are now overwhelmed by the opinions not only of the common people, but also of men lightly learned, who, what they cannot embrace whole, these they more easily handle torn apart and as it were rent in pieces — and who, as body from soul, so words from thoughts they separate (without the perishing of which neither can come about) — I shall not undertake more in my speech than is laid on me. I shall only signify briefly that no ornament of words can be found, save when the thoughts have been distinguished and expressed, nor is there any thought illustrious without the light of words.

\ciceroSection{25} But before I try to touch those things by which I think speech is to be ornated and lit up, I shall briefly set out what I think about the universal kind of speaking. There is no nature, as it seems to me, which does not have within its own kind several things unlike one another, which are yet judged worthy of similar praise. For we perceive by the ears many things which, although they delight us by sounds, are yet so various that what you have just heard seems most pleasant. And by the eyes are gathered well-nigh countless pleasures, which so seize us that they delight one sense in unlike kinds. The other senses too disparate pleasures delight, so that the judgment of the most outstanding sweetness is hard.

\ciceroSection{26} And the same which is in the natures of things can be transferred also to the arts. There is one art of modelling, in which Myro, Polyclitus, Lysippus were outstanding, all unlike one another, but yet in such a way that you would wish no one of them to be unlike himself. There is one art and rationale of painting, and yet most unlike one another are Zeuxis, Aglaophon, Apelles; nor is any of these one to whom anything in his own art seems wanting. And if this is wonderful in those, as it were, mute arts, and yet true, how much more wonderful in speech and in the tongue? Which, although it moves in the same thoughts and words, has the highest dissimilarities — not so that some should be blamed, but so that those who are agreed to be praised may yet in unlike kinds be praised.

\ciceroSection{27} And this is first to be seen in the poets, with whom is the closest kinship to orators: how unlike one another are Ennius, Pacuvius, and Accius; how unlike, among the Greeks, Aeschylus, Sophocles, Euripides — although nearly equal praise is assigned to all in unlike kinds of writing.

\ciceroSection{28} Look now upon those men, and gaze upon them — concerning whose faculty we are inquiring what is the difference of their pursuits and natures. Sweetness Isocrates, subtlety Lysias, keenness Hyperides, resounding Aeschines, force Demosthenes had. Which of them was not outstanding? Yet who is like anyone but himself? Weight Africanus had, gentleness Laelius, harshness Galba, a flowing and singing thing Carbo. Which of these was not chief in those times? And yet each was chief in his own kind.

\ciceroSection{29} But why should I seek out old examples, when I may use present and living ones? What pleasanter has ever fallen on our ears than the speech of this Catulus? — which is so pure, that he seems to speak Latin almost alone; and so weighty, that, in singular dignity, all human kindness and charm is yet present. In short, listening to him I am wont so to judge: that whatever you should add or change or take away, the speech would be more faulty and the worse.

\ciceroSection{30} Has not our Caesar here brought a kind of new principle of speech and introduced an almost singular kind of speaking? Who, except him, has ever handled tragic matters almost in comic vein, gloomy ones lightly, severe ones cheerfully, forensic ones with almost stage grace, and so that neither was the jest excluded by the greatness of the matters, nor was the gravity diminished by the wit?

\ciceroSection{31} Behold, present here, two nearly contemporaries, Sulpicius and Cotta. What so unlike between them? What so outstanding in his own kind? The one polished and subtle, setting forth the matter with proper and apt words; he sticks always in the case, and when he has most keenly seen what should be proved to the judge, omitting other arguments, on this he fixes his mind and speech. But Sulpicius, with the bravest impulse of soul, with the fullest and greatest voice, with the highest contention of body and dignity of motion, has also such weight and abundance of words that alone of all he seems to be by nature most equipped for speaking.

\ciceroSection{32} I come back now to ourselves, since we have always been so compared that, as it were, we have been called by the conversations of men into some judgment of contention. What so unlike as I in speaking and Antonius? — when he is such an orator that nothing can be more outstanding than he, while I, although I am dissatisfied with myself, am yet most especially joined with him in comparison. Do you see what kind of orator Antonius is? Strong, vehement, moved in pleading, fortified and hedged on every side of the case, keen, sharp, kernel-cleared, dwelling on each thing in its place, honourably yielding, keenly pursuing, terrifying, supplicating, with the highest variety of speech, with no satiety to our ears.

\ciceroSection{33} As for ourselves, whoever we are in speaking — since we seem to you to be in some account — we surely differ much from his kind. What kind that is, it is not for me to say, because each man is least known to himself, and most hardly each can judge of himself. But still the unlikeness can be understood from the moderate measure of my own motion, and from this — that the footsteps in which I have first set out, in those I am wont mostly to peroration, and that somewhat greater labour and care vexes me in choosing words than thoughts, fearing lest, if my speech is a little obsolete, it should not seem to have been worthy of expectation and silence.

\ciceroSection{34} Now if among us who are present there are such great unlikenesses, such fixed proper qualities of each, and in that variety it is rather the better faculty than the kind that distinguishes the better from the worse, and everything is praised which is perfect in its own kind — what do you suppose, if we should wish to embrace all who anywhere are or have been orators, will not the result be that, as many as the orators, almost so many kinds of speaking are found? From which discussion of mine perhaps that may occur — that, if there are well-nigh countless, as it were, forms and figures of speaking, unlike in appearance, praiseworthy in kind, things which differ among themselves cannot be shaped by the same precepts and one institution.

\ciceroSection{35} Which is not so. And this is most diligently to be looked at by those who train and instruct anyone: whither each man's nature seems most to bear him. For we see that from the same school, as it were, of the foremost masters and teachers in each one's kind, have come pupils unlike one another, and yet praiseworthy, when the teacher's training was adapted to each one's nature.

\ciceroSection{36} Of which especially marked is that example — to leave the rest of arts aside — which Isocrates, the singular teacher, used to say: that he was wont to use spurs on Ephorus, but, on the other hand, the bridle on Theopompus. For the one, leaping with audacity of words, he restrained; the other, hesitating and as it were bashful, he urged on. Nor did he make them like one another, but only added to one and shaved from the other, so that he shaped in each what each one's nature would bear.

\ciceroSection{37} These things had to be foretold by me, that, if not all the things proposed by me are stuck to the zeal of all of you and to that kind which each of you approves in speaking, you might feel that that kind is being expressed by me which most is approved by myself. Therefore those things must both be done by the orator which Antonius set out, and be said in some way. What better way of speaking, then — for of delivery I shall see later — than that we speak Latin, plainly, ornately, aptly and congruently to whatever is being treated?

\ciceroSection{38} As to the first two of these I have spoken of, I do not think the rationale of pure and clear speech is expected from me. For we do not try to teach speaking him who does not know to talk; nor to hope that he who cannot speak Latin will speak ornately; nor that he who does not say what we may understand can say something we shall admire. Let us therefore leave aside these things which have an easy understanding and a necessary use. For the one is handed down by letters and a boy's instruction; the other is brought to bear so that what each one says may be understood — which we see is so necessary that yet there can be nothing less.

\ciceroSection{39} But all elegance of speaking, although it is polished by the science of letters, is yet increased by reading the orators and poets. For those ancients, who could not yet adorn what they said, all spoke nearly splendidly. Those who shall have grown accustomed to their conversation cannot, even if they wish, speak otherwise than in Latin. Yet we shall not have to use those words which our usage no longer uses, save sometimes for the sake of ornament sparingly, as I shall show; but the man who has been zealously and much turned over in old writings will be able so to use accustomed words as to use the choicest.

\ciceroSection{40} And in order to speak Latin, not only must we look that we should bring forth those words which no one would justly reproach, and so preserve them in their cases and times and gender and number that nothing be confused or out of agreement or backward — but the very tongue and breath and sound of voice itself is to be regulated.

\ciceroSection{41} I do not wish letters to be expressed too pedantically, nor to be obscured too negligently; I do not wish words to come out meagrely lifeless, nor swelled and as it were heavily panted. For of voice I do not yet say those things which belong to delivery, but this — what seems to me as it were joined with conversation. For there are certain faults which there is no one but desires to flee: a soft voice or womanly, or as it were unmeasured, off-tune, absurd.

\ciceroSection{42} There is, however, a fault which some pursue of set purpose: a rustic and rural voice delights some, that the more, if it sound thus, their conversation may seem to retain antiquity. As your friend, Catulus, L. Cotta, seems to me to delight in the heaviness of his tongue and the rural sound of his voice; and he thinks what he says will seem ancient if it is plainly rustic. But your sound and that subtlety of yours delights me — I leave aside that of words, although it is the chief thing — but that, reasoning brings, letters teach, the habit of reading and of speaking confirms. I mean this sweetness which comes out of the mouth: which, just as among the Greeks the speech of the Athenians, so in Latin speech is most proper to this city.

\ciceroSection{43} At Athens long ago the learning of the Athenians themselves perished; only the home of pursuits remains in that city, in which the citizens are at leisure, and foreigners enjoy them, captured in some way by the name of the city and its authority. Yet any unlearned Athenian will easily surpass the most learned Asiatic men — not in words, but in the sound of voice, and not so much by speaking well as sweetly. Our men devote themselves less to letters than the Latins; yet of those city-bred you know — among whom is the least of letters — there is no one who would not easily overcome the most lettered of all togaed men, Quintus Valerius of Sora, by the gentleness of his voice and the very pressure and sound of his mouth.

\ciceroSection{44} Therefore, since there is some certain voice proper to the Roman race and city, in which nothing may offend, nothing displease, nothing be marked, nothing sound or smell foreign — let us follow that, and learn to flee not only rustic harshness but also foreign novelty.

\ciceroSection{45} For when I hear my mother-in-law Laelia — for women more easily preserve uncorrupted antiquity, who, having had no part in the conversation of many, always hold those things which they first learned — but I so hear her that I seem to be hearing Plautus or Naevius. So upright and simple is the very sound of her voice, that she seems to bring nothing of ostentation or imitation. From which I judge that her father so spoke, so her ancestors — not roughly as that man I spoke of, not vastly, not rustically, not gappingly, but pressed and evenly and gently.

\ciceroSection{46} Therefore our Cotta — whose broadness, Sulpicius, you sometimes imitate, dropping the letter \textit{i} and saying \textit{e} most fully — seems to me to imitate not the ancient orators but the harvesters.'' Here, when Sulpicius himself had laughed, ``So I shall deal with you,'' said Crassus. ``Since you wished me to speak, you shall hear something of your faults.'' ``Would that we might!'' he said. ``For that very thing we wish; and if you do it, many faults, I think, we shall lay aside today.''

\ciceroSection{47} ``But not without my own danger,'' said Crassus, ``can I rebuke you, Sulpicius, since Antonius said you seemed to him most like himself.'' Then he: ``You yourself indeed warned us that we should imitate those things which were greatest in each. From which I fear lest I be nothing of yours but the imitator of the foot-stamp and a few words and some movement, perhaps.'' ``Therefore those things which you have from me,'' said Crassus, ``I do not rebuke, lest I deride myself — and they are far more and greater than you say. Those, however, which are plainly yours, or expressed by imitation from someone, of these I shall warn you, if some place shall remind me.

\ciceroSection{48} Let us pass over then the precepts of speaking Latin, which boyish learning hands down, and which subtler knowledge and the rationale of letters nourishes, or the habit of daily and household conversation; books confirm them, and the reading of the ancient orators and poets. Nor let us linger longer on that other point — to argue by what means we may attain that what we say be understood:

\ciceroSection{49} by speaking, of course, in Latin, with words usual and pointing properly to those things we wish to be signified and made plain; without ambiguous word or speech, without too long a continuity of words, without too far drawing out things which by way of likeness are transferred from other matters; without thoughts torn apart, without backward times, without confused persons, without disturbed order. In short: the whole matter is so easy that it often seems to me very strange when it is more difficult to understand what the patron wishes to say than if he himself, who calls in the patron, were to speak about his own affair.

\ciceroSection{50} For those who bring cases to us, mostly themselves teach us so that you do not desire it to be said more plainly. But the same matters, as soon as Fufius or your contemporary Pomponius begins to plead, I do not equally understand what they say, unless I attend most carefully. So confused is the speech, so disturbed, that nothing is first, nothing second; and there is so great an unaccustomedness and a throng of words that the speech, which ought to bring light to matters, brings obscurity and shadows; and they seem in some way themselves to drown out themselves in speaking.

\ciceroSection{51} But, if it pleases — since I hope these things at any rate seem troublesome and putrid to you, the elder ones — let us pass to the rest, somewhat more burdensome.'' ``But you see,'' said Antonius, ``how unwillingly we are doing other things and how unwillingly we are listening to you, who could be drawn — for I conjecture from myself — leaving everything else to follow you. So polished is your speech on rough matters, so full on meagre, so new on most common.''

\ciceroSection{52} ``For those two parts were easy, Antonius,'' he said, ``which I have just run through or rather almost passed over: of speaking Latin and of speaking plainly. The rest are great, complex, various, weighty, in which all wonder of talent, all praise of eloquence is contained. For no one ever wondered at an orator because he spoke Latin; if it is otherwise, they laugh, and they do not think him an orator only, but even a man. No one has lifted up with words him who had so spoken that those who were present understood what he said, but they have despised him who could less do that.

\ciceroSection{53} In whom, then, do men shudder? Whom do they gaze upon stupefied as he speaks? At whom do they cry out? Whom do they think a god, so to say, among men? Those who speak distinctly, who set out clearly, who copiously, who illustriously both with matters and words, and in the very speech as it were make some rhythm and verse — that is, what I mean by ornately. Those same who so moderate themselves as the dignity of matters, of persons, demands — those are to be praised in that kind of praise which I call apt and congruent.

\ciceroSection{54} Antonius said he had not yet seen those who so spoke, and said this name of eloquence was to be assigned to them alone. Therefore deride and despise — on my counsel — all those who think they have embraced all the orator's force in the precepts of those rhetoricians as they are now called, and have not yet been able to understand what character they hold or what they profess. For everything in the life of men — since the orator moves in it and that is his subject-matter — must be sought, heard, read, argued, handled, stirred up.

\ciceroSection{55} For eloquence is one of the highest virtues; although all virtues are equal and on a par, yet there is in appearance one more comely and brighter than another, as is this faculty, which, having grasped the knowledge of things, so sets out the perceptions and counsels of mind by words that it can drive those who hear, wherever it shall lean. As its force is the greater, the more must it be joined to uprightness and the highest prudence; if to those without these virtues we should hand the abundance of speaking, we shall not have made them orators but shall have given certain weapons to madmen.

\ciceroSection{56} This rationale of thinking and pronouncing and the force of speaking, the ancient Greeks called wisdom; from this came the Lycurgi, the Pittaci, the Solons; and from this likeness our Coruncanii, Fabricii, Catones, Scipiones, perhaps not so learned, but with similar impulse and disposition of mind. Others with the same prudence, but with a different counsel about the pursuits of life, sought rest and leisure: as Pythagoras, Democritus, Anaxagoras, transferred themselves wholly from governing states to the knowledge of things; which life, on account of its tranquility and the sweetness of the science itself (than which nothing is more pleasant to men), has delighted more men than has been useful to commonwealths.

\ciceroSection{57} So when men of the most outstanding talents had given themselves to that pursuit, out of that highest faculty of free and empty time the most learned men, abounding in too much leisure and richest talents, took on themselves to handle and seek and investigate far more things than was necessary. For that old learning seems the same — both the teacher of doing rightly and of speaking well. Nor were the teachers separate, but the same were teachers of living and of speaking — as that Phoenix in Homer, who says he was given by Achilles' father Peleus as companion to the young man for war, that he might make him an orator of words and an actor of deeds.

\ciceroSection{58} But as men accustomed to constant and daily labour, when on account of weather they are kept from work, betake themselves to ball or knucklebones or dice, or even invent for themselves some new game in their leisure — so those, kept from public business as it were from work, by the times or by their own will idle, betook themselves wholly, some to poets, some to geometers, some to musicians; some, too, like the dialecticians, brought forth for themselves a new study and game; and in those arts which were discovered, that the minds of boys might be shaped to humanity and virtue, they used up all their time and ages.

\ciceroSection{59} But because there were some — and many — who flourished in public life by reason of the twin wisdom of doing and speaking which cannot be separated, like Themistocles, Pericles, Theramenes, or who, less themselves engaged in public affairs, were yet teachers of this same wisdom, like Gorgias, Thrasymachus, Isocrates — there were found those who, themselves abounding in learning and talents, but recoiling from civic affairs and public business by some judgment of mind, harassed and despised this practice of speaking;

\ciceroSection{60} of whom Socrates was the leader, who, by the testimony of all the learned and by the judgment of the whole of Greece, was easily first of all in prudence and keenness and grace and subtlety, and indeed in eloquence, variety, abundance, into whichever side he had given himself. From those who handled, did, taught these things which we are now seeking — when they were called by one name (because all knowledge of the best things and exercise in them was named philosophy) — Socrates snatched away this common name, and separated by his own disputations the science of feeling wisely and speaking ornately, although they cohered in fact. His talent and various conversations Plato handed to immortality by his writings, since Socrates himself had left no letter.

\ciceroSection{61} Hence arose that division as it were of tongue and heart — absurd surely and unprofitable and to be reproved — that some teach us to be wise, others to speak. For when several had risen mostly from Socrates, because each had grasped one thing from his various and diverse and on every side spread disputations, families were sown among one another quarrelling, much severed and unlike — when yet all the philosophers wished to be called and judged themselves to be Socratics.

\ciceroSection{62} And first from Plato himself Aristotle and Xenocrates, of whom one held the name of the Peripatetics, the other of the Academy. Then from Antisthenes, who in Socrates' conversation had most loved patience and hardness, the Cynics first, then the Stoics. Then from Aristippus, whom those rather voluptuous disputations had pleased, flowed the Cyrenaic philosophy, which he and his successors defended simply; while these, who now measure all by pleasure, while they do this more bashfully, neither satisfy dignity (which they do not despise) nor protect pleasure (which they wish to embrace). There were also other kinds of philosophers — who almost all said they were Socratics — Eretrians, Herillians, Megarians, Pyrrhonians; but these have long ago been broken and extinguished by the force and arguments of these.

\ciceroSection{63} Of those that remain, that philosophy which has undertaken the patronage of pleasure, even if it seems to anyone true, is yet far from that man whom we are seeking, and whom we wish to be the author of public counsel and the leader of governing the state, and chief of opinion and eloquence in the senate, before the people, in public causes. Yet no injury will be done by us to that philosophy. It will not be repelled from where it shall wish to enter, but will rest in its little gardens — where it wishes, where, even reclining softly and daintily, it calls us off from the Rostra, from the courts, from the curia. Wisely perhaps, especially in this commonwealth.

\ciceroSection{64} But I am not now asking what philosophy is most true, but what is most joined to the orator. So let us send those men off without any insult; for they are good men, and (since they think themselves so) blessed. Let us only warn them that, even if it is most true, they should hold it silent like a mystery — that they say it is not the wise man's part to be in public life. For if they shall persuade us and every best man of this, they cannot themselves be — what they most desire — at leisure.

\ciceroSection{65} The Stoics, however — whom I least disapprove — I yet send away, nor do I fear them angry, since they do not at all know how to be angry. To them I owe this thanks: that they alone of all said eloquence is virtue and wisdom. But there is doubtless in them what greatly recoils from the orator we are equipping. Either because they say all who are not wise are slaves, robbers, enemies, mad — and yet that no one is wise. But it is most absurd to commit an assembly or senate or any gathering of men to him to whom no one of those who are present seems to be sane, no one a citizen, no one a free man.

\ciceroSection{66} It is added that they have a kind of speech perhaps subtle and surely keen, but, for an orator, meagre, unaccustomed, recoiling from the ears of the common folk, obscure, empty, dry, and yet of such a kind that to use it before the multitude is in no way possible. For other things seem good and bad to the Stoics and to the rest of the citizens — or rather, the nations — different is the force of honour, of disgrace, of reward, of punishment; truly or falsely, nothing for the present matters; but if we should follow them, we shall never be able to bring any matter to a head by speaking.

\ciceroSection{67} The Peripatetics and the Academics remain. Although the name of the Academics is one, the opinions are two. For Speusippus, Plato's sister's son, and Xenocrates who heard Plato, and Polemo and Crantor who heard Xenocrates, did not greatly disagree from Aristotle, who together with them had heard Plato; in the abundance and variety of speaking they were perhaps not equal. Arcesilas first, who had heard Polemo, snatched up especially this from Plato's various books and Socratic conversations: that nothing certain can be perceived either by the senses or by the soul. He, having used a certain extraordinary charm of speaking, is said to have spurned all judgment of soul and sense; and he first established (although it had been chiefly Socratic) not to show what he himself thought, but to argue against what each had said he thought.

\ciceroSection{68} Hence flowed this more recent Academy, in which Carneades stood out by a kind of divine swiftness of talent and abundance of speaking. Although I have known many of his hearers at Athens, I can yet name most certain authorities: my father-in-law Scaevola, who as a young man heard him at Rome; and my friend Q. Metellus, son of Lucius, that most distinguished man, who said that Carneades had been heard by him as a young man at Athens, in his old age, for many days.

\ciceroSection{69} These — as the rivers from the Apennine, so from the common ridge of wisdom — were made the partings of doctrines: so that the philosophers as it were flow down into the upper sea, the Ionian, a Greek and harboured one, while the orators are slipped into this lower, Tuscan and barbarian one, full of rocks and dangers, in which even Ulysses himself wandered.

\ciceroSection{70} Therefore, if we are content with this eloquence and this orator, who knows either that he must deny what is charged against him, or, if he cannot, then to show that what the accused did was either rightly done, or by another's fault, or by injustice, or by law, or not against law, or imprudently, or necessarily, or that the matter must not be claimed under the name by which it is charged, or that the case is not being managed as it ought and was lawful — and if you think it enough to learn what those writers of the art teach, which yet Antonius has set out far more ornately and copiously than they say it — but if you are content with these, and with what you wished to be said by me, you compel the orator out of an enormous and unbounded field into a really narrow ring.

\ciceroSection{71} But if you wish to follow that ancient Pericles, or this man, more familiar to us because of the multitude of his writings, Demosthenes; and if you have fallen in love with that brilliant and exceptional sight of the perfect orator, and his beauty — either this Carneadean or that Aristotelian force is to be grasped by you.

\ciceroSection{72} For, as I said before, those ancients up to Socrates joined the knowledge and science of all things which pertain to the morals of men, to life, to virtue, to the commonwealth, with the rationale of speaking. Afterwards, when (as I have set out) the eloquent had been parted from the learned by Socrates, and thereafter the philosophers, like all the Socratics, despised eloquence, and the orators wisdom, neither touched anything of the other's part, save what these borrowed from those, or those from these — from which they would have drawn promiscuously, if they had wished to remain in their pristine community.

\ciceroSection{73} But as the ancient pontiffs, on account of the multitude of sacrifices, wished there to be three men as banqueters, when they themselves had been instituted by Numa to perform also that banquet sacrifice of the games — so the Socratics severed from themselves the pleaders of cases, and from the common name of philosophy, when the ancients had wished there to be a wonderful fellowship of speaking and understanding.

\ciceroSection{74} Since these things are so, I shall a little plead for myself, and I shall ask you that what I shall say you should think to be said, not about myself, but about the orator. For I am one who, though taught with the highest zeal of my father in boyhood, and brought to the Forum so much talent (as I myself feel — not so much as perhaps I seem to you), cannot say that I have learned these things which I am now embracing as I shall say they ought to be learned. For of all I came most early to public causes, and at twenty-one called into court a most noble and most eloquent man. The Forum was my school; my master was practice and the laws and institutions of the Roman people and the custom of the ancestors.

\ciceroSection{75} I tasted only a little of those arts of which I am speaking, when I was quaestor in Asia, having got the rhetorician of about my own age from the Academy, that Metrodorus of whose memory Antonius spoke; and on my way back from there, at Athens — where I should have lingered longer, had I not been angry with the Athenians, who would not put on the mysteries again, to which I had come two days too late. So this — that I embrace so great a knowledge and force of learning — is not for me but rather against me. For I am arguing not what I can do, but what the orator can; and all those who set out the rhetorical arts most ridiculously: for they write of the kinds of suits, of openings, of narrations.

\ciceroSection{76} But the force of eloquence is so great that it grasps the origin, force, and changes of all things, all virtues, duties, all nature, which contains the morals of men, the souls, the life; the same describes morals, laws, rights, governs the commonwealth, and ornately and copiously speaks of everything, to whatever matter it pertains.

\ciceroSection{77} In this kind we move as much as we can, as much as we are strong by talent, by moderate learning, by use; nor do we yield much in disputation to those who have set up their tent of life in philosophy alone.

\ciceroSection{78} For what can my friend Gaius Velleius bring why pleasure should be the highest good, that I could not more copiously either defend, if I wished, or refute, from those topics which Antonius set out, by this practice of speaking — in which Velleius is unschooled, and each of us has been versed? What is there, that either Sextus Pompeius or the two Balbi, or my friend who lived with Panaetius, M. Vigellius the Stoic, could say of the virtue of men in such an argument, in which I should have to yield to them, or any of you?

\ciceroSection{79} For philosophy is not like the rest of the arts. For what can he do in geometry who has not learned? What in music? He must either be silent or not even be reckoned sane. But these things which are in philosophy are dug out by talents to elicit what in each thing is like the truth, and they are polished by exercised speech. Our common orator here, if he is less learned, but yet exercised in speaking, by this very common exercise will beat off those of ours, nor will he allow himself to be despised and contemned by them.

\ciceroSection{80} But if anyone shall ever come forth who can in the Aristotelian manner speak on every subject on either side, and in every case set out two contrary speeches by his precepts; or who in this manner of Arcesilas and Carneades shall argue against everything that is laid down — and shall add to that rationale this rhetorical use and habit and exercise of speaking, he will be the true, the perfect, the only orator. For neither without the forensic sinews can the orator be sufficiently vehement and weighty, nor without variety of learning sufficiently polished and wise.

\ciceroSection{81} Therefore let us suffer that ancient Corax to hatch his own chicks in the nest, that they may fly out as hateful and troublesome shouters; and let us allow that I-know-not-who Pamphilus to portray, in fillets, so great a matter as if it were certain childish playthings. Let us ourselves, by this so brief disputation of yesterday's and today's day, set out the orator's whole duty — provided that thing be so great that it seems to be embraced in all the philosophers' books which none of those orators of yours has ever touched.''

\ciceroSection{82} Then Catulus: ``In no way, by Hercules,'' he said, ``Crassus, is it to be wondered at that there is in you so great a force of speaking, or sweetness, or copiousness. Of you indeed I before judged that you spoke by nature in such a way that you seemed to me not only the highest orator but also the wisest man. Now I understand that you have always reckoned even more important those things which look towards wisdom, and that out of these has flowed this copiousness of speaking. Yet, when I recall all the steps of your age, and consider your life and pursuits, I do not see at what time you learned those things, nor do I understand you to have been greatly given to those pursuits, men, books. And yet I cannot decide whether more I should wonder that you, in your such occupations, could have learned those things which you persuade me are the greatest helps; or, if you could not, that you can speak in this way.''

\ciceroSection{83} Here Crassus: ``This,'' he said, ``Catulus, I would have you persuade yourself first: that I do not, when I argue about the orator, do much otherwise than I should do, if I had to speak about an actor. For I should deny that he could satisfy in gesture, unless he had learned the wrestling-school, unless he had learned to dance. Nor when I said this would it be necessary that I be an actor, but perhaps not a foolish judge of another's art.

\ciceroSection{84} Similarly now I am speaking of the orator at your urging — of the highest, of course; for always, whatever art or faculty is asked about, the question is wont to be of the absolute and perfect. Therefore, if you wish me to be even an orator, even tolerably good, even good — I shall not resist; for why should I now be silly? I know I am so reckoned. But if it is so, I am surely not the highest. For there is no thing among men either harder or greater, or which requires more helps of learning.

\ciceroSection{85} And yet, since we must argue about the orator, I must speak about the highest orator. For the force and nature of the matter cannot be understood, of what kind and how great it is, unless the perfect is set before the eyes. As for me, Catulus, I confess that I do not now live in those books and with those men, nor indeed (which you rightly remember) ever had any time set apart for learning, nor have given to learning so much time as my boyish age and forensic holidays have allowed me.

\ciceroSection{86} And, if you ask, Catulus, what I think of that learning of yours: I think there is need of so much time not to a talented man, who has the Forum, the curia, cases, the commonwealth before him, as those have taken on themselves to whom learning has lasted to the end of life. For all arts are handled differently by those who transfer them to use, differently by those who, delighted with the handling of the arts themselves, are about to do nothing else in life. The Samnite gladiator-master here is now in extreme old age and is exercising daily, for he cares for nothing else; but Q. Velocius as a boy had learned a little, but, because he was apt for it and had wholly mastered it, was, as is said in Lucilius, ``Although a good Samnite himself in the school and with the wooden swords, severe enough to anyone''; but he gave more pains to the Forum, to friends, to private affairs. Valerius sang every day, for he was an actor: what else could he do?

\ciceroSection{87} But our friend Numerius Furius sings when it is convenient, for he is a head of household, a Roman knight; as a boy he learned what was to be learned. The same rationale is in these greatest arts. Day and night we used to see Q. Tubero, a man of the highest virtue and prudence, when he was giving pains to philosophy. But his uncle Africanus you would scarcely understand to be doing this, although he was doing it. These things are easily learned if you both take so much as is needed, and have one who can teach faithfully, and yourself know how to learn;

\ciceroSection{88} but if you wish to do nothing else in your whole life, the very handling and inquiry produces something out of itself daily, that you may track it down with idle delight. So it happens that the stirring of subjects is unbounded, the knowledge easy, if practice confirms the learning, moderate effort is given, memory and zeal remain. But it is always pleasant to learn — as if I should wish to play knucklebones excellently, or be held by zeal for the ball — even, perhaps, if I cannot attain it. But others, because they do brilliantly, are more vehemently than the case requires delighted, as Titius with the ball, Brulla with the knucklebones.

\ciceroSection{89} Therefore there is no reason why anyone should fear the magnitude of the arts because old men learn them. For either old men have approached them, or are detained in those studies right up to old age, or are most slow. The matter, in my view, is thus: that, unless what each could learn quickly, he could never altogether master.''

\ciceroSection{90} ``Now, now,'' said Catulus, ``Crassus, I understand what you say. And by Hercules I agree. I see well enough that for you, a man most keen at mastering, there has been time enough for learning those things you say.'' ``Do you persist,'' said Crassus, ``in thinking that what I say I am saying about myself, not about the matter? But now, if it pleases, let us return to our institutes.'' ``It does please me,'' said Catulus.

\ciceroSection{91} Then Crassus: ``Whither, then, does this so long and so deeply drawn-out discourse look? These two parts which remain to me, of illustrating the speech and of crowning the whole eloquence — of which one demands ornate speaking, the other apt — have this force: that the speech be as pleasing as possible, that it flow into the senses of those who hear as much as possible, and that it be furnished with as many things as possible.

\ciceroSection{92} But this forensic equipment, litigious, keen, drawn from the opinions of the common people, is small and quite mendicant. That itself, again, which those hand down who profess themselves teachers of speaking, is not much greater than that vulgar and forensic stock. We need furnishing and exquisite matters gathered, sought, and brought together from every side — as it is for you to do, Caesar, next year; as I laboured in my aedileship, because I did not think I could satisfy this people with everyday and homespun matters.

\ciceroSection{93} The rationale of choosing and arranging and concluding words, or even — without rationale — practice itself, is easy. Of matters there is a great wood. Since the Greeks no longer hold this, and on that account our youth nearly unlearns by learning, even Latins (so it pleases the gods) in this two-year stretch have arisen as teachers of speaking. Whom I, as censor, removed by my edict — not because, as some, I am told, said, I should not wish the talents of the youth to be sharpened, but on the contrary, I should not wish their talents to be dulled, their shamelessness strengthened.

\ciceroSection{94} For I saw that, among the Greeks, of whatever kind they were, there was, beyond this exercise of the tongue, some kind of doctrine and science worthy of culture; but these new masters I understood could teach nothing save daring; which, even when joined with good things, is by itself greatly to be fled. Since this alone was being handed down, and since it was a school of impudence, I thought it was the censor's part to look out lest it should creep further.

\ciceroSection{95} Although I do not so lay down and decree these things as if I were despairing that those things on which we have argued can be handed down and polished in Latin. For both our tongue and the nature of things suffers that the ancient and outstanding prudence of the Greeks be transferred to our use and custom. But there is need of learned men, of whom thus far in our country there have been none in this kind. But if any shall ever stand forth, they will be set even before the Greeks.

\ciceroSection{96} Speech then is adorned first by kind and as it were by a certain colour and sap of its own. For that it be weighty, that it be pleasing, that it be learned, that it be liberal, that it be admirable, that it be polished, that it have feeling, that it have grief as much as is needed — that does not lie in single joints. These things are seen in the whole body. Then, that it be sprinkled, as it were, with flowers of words and thoughts — that ought not to be diffused evenly through the whole speech, but so distinguished that there should be, as in ornament, certain marks set out and lights.

\ciceroSection{97} A kind of speaking, then, must be chosen which most holds those who hear, and which not only delights but delights without satiety. For I do not think it is now expected from me to warn that you should beware lest your speech be meagre, lest unkempt, lest vulgar, lest obsolete. Something else greater both your talents and your ages exhort me to.

\ciceroSection{98} For it is hard to say what cause there is why those things which most strike our senses with pleasure and most keenly stir us at first sight, we are most quickly turned away from in a kind of disgust and satiety. How much more flowery in beauty and variety of colours are most paintings new than old? Which yet, even if they have at first sight caught us, do not delight us long; whereas we are held by those very rude and obsolete features in old paintings. How much softer and more dainty in song are the inflexions and the false little notes than the fixed and severe? At which not only the austere, but, if they happen too often, the multitude itself cries out.

\ciceroSection{99} It is to be seen in the rest of the senses: that we are less long delighted by perfumes seasoned with the highest and keenest sweetness than by these moderate ones; and what seems to smell of earth is more praised than what seems to smell of saffron. In touch itself there is a measure both of softness and of smoothness. Even taste, which is the most pleasure-loving of all the senses, and which is moved by sweetness above the other senses, how quickly does it scorn and reject what is very sweet! Who can drink or eat what is sweet for long? — when in each kind those things which lightly move the sense with pleasure most easily flee satiety.

\ciceroSection{100} So in all things disgust is the neighbour of the greatest pleasures. Wherefore we should the less wonder at this in speech, in which we can judge from poets or from orators that whatever poetry or speech be neat, distinguished, ornate, festive, without intermission, without reproach, without variety, even painted with bright colours, cannot be lasting in delight. And we are the more quickly offended by the orator's or the poet's curls and rouge, because the senses by nature, not by mind, are sated in too much pleasure. In writings and in sayings the painted faults are known not only by the ears but, even more, by the judgment of the soul.

\ciceroSection{101} Therefore though to us ``well and brilliantly'' be often said, ``prettily and gaily'' I will not have too often said. Although that very exclamation, ``It cannot be better,'' I would have frequent. But yet let that admiration and highest praise in speaking have some shadow and recess, that what is illumined may seem more to stand out and shine.

\ciceroSection{102} Roscius never plays this verse with the gesture he can: ``\textit{For the wise man seeks honour as the reward of virtue, not booty}.'' But he throws it off entirely, that in the next: ``\textit{But what do I see? Hedged with steel, he holds the sacred seats}'' — he may push, may look, may admire, may stupefy. What of that other: ``\textit{What guard shall I seek?}'' — how gently, how slackly, how unanimatedly! For there presses next: ``\textit{O father, O country, O house of Priam!}'' — at which a delivery so great could not be moved, if it had been used up by the previous motion and exhausted. Nor did the actors see this before the poets themselves, and finally those too who composed the music; from each of which something is lowered, then increased, thinned, swelled, varied, distinguished.

\ciceroSection{103} So let our orator be ornate and pleasing — nor indeed can he be otherwise — but in such a way that he has a sweetness austere and solid, not sweet and over-cooked. For the precepts that are given for ornament are of such a kind that anyone, even the most faulty orator, can set them out. Therefore, as I said before, first the wood of matters and thoughts must be gathered, of which Antonius spoke. These must be shaped by the very thread and kind of speech, illustrated by words, varied by thoughts.

\ciceroSection{104} The highest praise of eloquence is to amplify a matter by ornament, which avails not only for raising something and lifting it up higher by speech, but also for thinning and casting it down. This is required in all those topics which Antonius said are applied for producing the credit of speech, either when we explain anything, or when we conciliate minds, or when we rouse them.

\ciceroSection{105} But in this last thing I have spoken of, amplification can do the most, and that is the one praise of the orator, and especially proper. Even greater is that exercise which Antonius placed in the closing speech and at first rejected — of praising and blaming. For nothing is more apt for exaggerating and amplifying speech than to be able to do each of these in the fullest way.

\ciceroSection{106} There follow also those topics which, although they ought to be proper to cases and clinging in their sinews, yet, because they are wont to be handled about the universal matter, were called by the ancients common — of which some have a kind of keen accusation or complaint of vices and faults with amplification, against which nothing is wont to be said and nothing can be: as against an embezzler, against a traitor, against a parricide. These should be used when the charges have been confirmed; otherwise they are dry and empty;

\ciceroSection{107} others have entreaty or pity; others, indeed, two-sided disputations, in which copiously on both sides the universal kind may be argued. This exercise, now thought to be proper to those two philosophies of which I spoke before, was among the ancients of those from whom the whole rationale and abundance of speaking on forensic matters was sought. For of virtue, of duty, of right and good, of dignity, utility, honour, disgrace, reward, punishment, and like matters, we ought to have both force and art for speaking on each side.

\ciceroSection{108} But since we have been driven from our possession and left in a small and litigious little farm, and we, advocates of others, could not hold and protect our own, from those who have broken into our patrimony — what is most unworthy — let us borrow what we need.

\ciceroSection{109} Those, then, who now indeed take their name from a small portion of the city and place — and are called Peripatetic or Academic philosophers — but in old time, by reason of their outstanding knowledge of the greatest things, were called by the Greeks ``political'' philosophers, by the name of the universal commonwealth — say that all civic speech moves in one of two kinds of these: either of a definite controversy in fixed times and parties (in this way: ``Should our captives be received back from the Carthaginians by returning theirs?''); or, of one who asks generally about the universal kind: ``What altogether ought to be decreed and felt about a captive?'' The earlier of these two kinds they call \textit{cause} or \textit{controversy}; and they define it by three: suit, deliberation, or laudation. The other, however, is called an unbounded \textit{question} and as it were a \textit{consultation} put forward.

\ciceroSection{110} And they so far speak, and use this division too in setting things up. But in such a way that they recover their lost possession not by right or by judgment, but, finally, by force; and that, not by civil law, but by breaking off a twig of the property, they seem to take possession by usage. For that other kind, which is defined by times, places, parties, they hold; and that itself only by a fringe — for now with Philo, who I hear is most flourishing in the Academy, the inquiry into and exercise of these very cases is being celebrated. The other only they name in the first handing-down of the art and say it is the orator's. But neither do they set out the force or the nature of it, nor its parts, nor its kinds; so that it had been better passed by altogether than to have been touched and abandoned. For now they are understood to be silent through poverty; then they would have seemed silent through judgment.

\ciceroSection{111} Every matter, therefore, has the same nature of being doubted, on which question and dispute can be made — whether it be argued in unbounded consultations, or in those cases which move in the state and forensic dispute. Nor is there any which is not referred to either the force and rationale of knowing or of doing.

\ciceroSection{112} For either the very knowledge and science of a thing is sought, as ``Is virtue to be desired for its own dignity, or for some fruit?'' Or the counsel of doing is asked, as ``Should the wise man take part in the commonwealth?''

\ciceroSection{113} There are three modes of knowing: conjecture, definition, and (so to say) consequence. For what is in a thing is asked by conjecture, as ``Is there wisdom in the human race?'' What force each thing has, definition sets out, as if it be asked, ``What is wisdom?'' Consequence is handled when it is asked what follows from each thing, as ``Is it sometimes the part of a good man to lie?''

\ciceroSection{114} They come back again to conjecture and divide it into four kinds. For either it is asked what something is, in this way: ``Is right by nature among men, or in opinions?'' Or what is the origin of each thing, as ``What is the beginning of laws or of commonwealths?'' Or the cause and reason, as if it should be asked, ``Why do most learned men disagree on the greatest matters?'' Or about a change, as if it should be argued, ``Whether virtue can perish in a man, or be turned to vice.''

\ciceroSection{115} The disputes of definition are either when it is asked what is impressed, as it were, in the common mind, as if it should be argued: ``Is that the right which is useful to the greatest part?'' Or when it is asked what is each thing's own, as: ``Is to speak ornately the orator's own, or can someone besides do it?'' Or when the matter is distributed into parts, as if it should be asked, ``How many kinds are there of things to be sought? Are there three: of the body, of the soul, of external things?'' Or when the form and as it were the natural mark of each thing is described, as if there be sought the kinds of the avaricious, the seditious, the boastful.

\ciceroSection{116} Two first kinds of questions of consequence are laid down. For either the dispute is simple, as if it be argued whether glory is to be sought; or from comparison, ``Is praise or wealth more to be sought?'' Of the simple, there are three modes: about things to be sought or fled, as ``Are honours to be sought, or poverty to be fled?''; about the just or unjust, as ``Is it just to revenge wrongs even against kinsfolk?''; about the honourable or base, as: ``Is it honourable to die for glory?''

\ciceroSection{117} Of comparison there are two modes. One when it is asked whether something is the same or has some difference: as to fear and to revere, as a king and a tyrant, as a flatterer and a friend. The other when it is asked what is preferable to what, as: ``Are wise men led by the praise of each best man, or by popular praise?'' These disputes which are referred to knowledge are about thus described by the most learned men.

\ciceroSection{118} Those, however, which are referred to action, either move in the dispute of duty (in which kind it is asked what is right and to be done, under which topic the whole wood of virtues and vices is set); or are handled in some movement of soul to be either produced or calmed or removed. To this kind are subjected exhortations, rebukes, consolations, pity, and every impulsion to every motion of soul, and, if the case so bear, mitigation.

\ciceroSection{119} These kinds and modes of all disputes being set out, it is nothing to the purpose if in any matter our partition has differed from Antonius's division: the same members are in each one's discussion, but parted and assigned a little differently by me and by him. Now I shall move on to the rest, and call myself back to my own task and weight. For from those topics which Antonius set out, all arguments are to be taken for each kind of question. But to one kind one topic will be more apt than another. About which it is unnecessary to say anything, not so much because it is long as because it is plain.

\ciceroSection{120} The most ornate, then, are those speeches which range most widely and bring and turn themselves from a private and singular controversy to setting out the force of the universal kind, so that those who hear, with the nature, kind, and universal matter known, may decide about the single parties, charges, and suits.

\ciceroSection{121} Antonius has exhorted you, young men, to this practice of exercise, and thought that you should be drawn from minute and narrow contests to the whole force and variety of arguing. Therefore this duty is not of a few little books, as those think who have written about the rationale of speaking; nor of the Tusculan villa or this morning walk or our after-noon sitting. For not only must our tongue be sharpened and forged, but the breast must be loaded and filled with the sweetness, abundance, and variety of the greatest and most numerous matters.

\ciceroSection{122} For ours — if indeed we are orators, if in the disputes of citizens, if in dangers, if in public deliberations we are to be applied as authors and chiefs — ours, I say, is all that possession of prudence and learning, into which men, abounding in leisure while we are busy, have leapt as upon something fallen and empty; and either, mocking the orator (as Socrates does in the \textit{Gorgias}), they bandy words; or they teach in a few books some of the orator's art and entitle them \textit{rhetorical}, as if those things were not the rhetoricians' own which are said by them on justice, on duty, on instituting and governing states, on every rationale of living, and indeed of nature.

\ciceroSection{123} Which since we cannot now from elsewhere — they must be taken by us from those very men by whom we have been despoiled. Provided we transfer those things to this civil science to which they pertain and at which they look; and provided, as I said before, we do not waste our whole life in learning these matters. But when we have seen the springs (which, unless one quickly come to know, he will never come to know at all), then, as often as need shall be, from them we shall draw as much as the matter shall require.

\ciceroSection{124} For there is not so keen an edge in the natures and talents of men that anyone could see such great matters unless they have been pointed out; nor yet is there so great obscurity in matters that a man of keen talent cannot deeply discern them, if only he has looked. In this so vast, so unbounded a field, when the orator may freely range, and wherever he has stood, stand on his own ground, all the equipment and ornament of speaking easily abounds.

\ciceroSection{125} For the abundance of matters begets the abundance of words; and if there is honour in the very things spoken about, there arises out of the matter a kind of natural splendour in words. Let only him who shall speak or write be liberally trained by his early education and learning, and burn with zeal and be helped by nature, and be exercised in unbounded disputations of universal kinds, and have chosen the most ornate writers and orators for knowing and imitating — he will surely not need those teachers to know how to construct and illumine words. So easily, in the abundance of matters, does nature itself, if only it has been exercised, slip without a guide to the ornaments of speech.''

\ciceroSection{126} Here Catulus: ``Immortal gods,'' he said, ``Crassus, what variety of matters, what force, what abundance you have grasped, and out of what straits have you dared to lead the orator and to set him in the kingdom of his ancestors! For we have heard that those ancient teachers and authors of speaking thought no kind of disputation alien from themselves, and were always engaged in every rationale of speech.

\ciceroSection{127} Of whom Hippias of Elis, when he had come to Olympia at that great five-yearly celebration of the games, boasted, almost with all Greece listening, that there was nothing in any art of all things that he himself did not know — and that not only did he hold those arts which contain the liberal and free-born doctrines, geometry, music, the knowledge of letters and poets, and what is said of the natures of things, of the morals of men, of commonwealths — but also that he had made by his own hand the ring he had on, the cloak he was clothed with, the slippers he had put on.

\ciceroSection{128} Of course this man went too far. But from him is easy conjecture how much those orators sought for themselves of the most distinguished arts — when they did not even shrink from baser ones. What shall I say of Prodicus of Ceos, Thrasymachus of Chalcedon, Protagoras of Abdera? Of whom each one, as much as in those times was possible, both argued and wrote about the nature of things.

\ciceroSection{129} That very Gorgias of Leontini, with whom as patron, as Plato wished, the orator yielded to the philosopher — who either was never overcome by Socrates and that conversation of Plato is not true, or, if he was overcome, then surely Socrates was more eloquent and skilled and (as you call him) more copious and a better orator. But this man, in that very book of Plato, professes that he will speak most copiously of every matter, whatever may be called into dispute and question. He first of all dared, in the assembly, to demand on what each wished to hear; to whom such honour was given by Greece, that for him alone of all there was set up at Delphi a statue not gilded but golden.

\ciceroSection{130} But these whom I have named, and very many besides, the highest teachers of speaking, were at one time. From whom it can be understood that the matter is, as you say, Crassus: that the name of orator among the ancients in Greece flourished with somewhat greater abundance and glory.

\ciceroSection{131} Wherefore I doubt the more whether more praise should be assigned to you or more blame to the Greeks: that you, born in another tongue and other customs, in a most busy state, distracted either by the affairs of nearly all private men or by the management of the world and the steerage of the highest empire, have grasped so great a force and knowledge of matters, and have joined all of it with the science and exercise of him who avails by counsel and speech in the state — while they, born in letters and burning with these studies, but flowing in leisure, not only have acquired nothing, but have not even kept what was left and handed down and their own.''

\ciceroSection{132} Then Crassus: ``Not in this matter alone, Catulus,'' he said, ``but in many others too, the magnitudes of the arts have been diminished by the distribution and separation of parts. Or do you suppose that, when that famous Hippocrates of Cos lived, there were then other doctors who healed diseases, others who healed wounds, others who healed eyes? Surely with Euclid or Archimedes, surely with Damon or Aristoxenus in music, surely with Aristophanes or Callimachus in letters, those things were not so torn apart that no one embraced the universal kind, and that one set apart for himself one part, another another, in which to labour?

\ciceroSection{133} For my part, I have often heard this from my father and my father-in-law: that our men also, who wished to excel in glory of wisdom, were wont to embrace everything which our state then knew. They remembered Sextus Aelius; M.' Manilius indeed we ourselves have seen walking across the Forum — which was a sign that he who did this offered to all his fellow-citizens the abundance of his counsel. To these, both as they walked and as they sat at home in their high seat, men used to come — not only that they might consult them about the civil law, but also about giving a daughter in marriage, about buying a farm, about cultivating a field, about every duty or business.

\ciceroSection{134} This was the wisdom of that old P. Crassus, of Tiberius Coruncanius, of my son-in-law's great-grandfather Scipio (a most prudent man) — all of whom were Pontifex Maximus, that things divine and human might be referred to them. The same gave their counsel and faith in the senate, before the people, in the cases of friends, at home and on military service.

\ciceroSection{135} For what was lacking to M. Cato besides this most polished doctrine from across the sea? Did he, because he had learned the civil law, not plead cases? Or, because he could plead, did he neglect the science of law? In each kind he both laboured and excelled. Did he, because of the favour gathered from private men's affairs, become slower in undertaking public affairs? No one was braver before the people, no one a better senator; the same easily the best general; in short, nothing in this state in those times could be known or learned that he did not both investigate and know — and even write.

\ciceroSection{136} Now, on the contrary, most come naked and unarmed to the gaining of offices and to the conduct of public affairs, equipped with no knowledge of things, with no science. But if anyone excels above the rest, he flaunts himself if he brings one thing — either martial virtue or some military experience (which now indeed have grown obsolete); or the science of law, but not even of all of it (for the pontifical law, which is joined to it, no one learns); or eloquence, which they think is set in shouting and in the rush of words. The fellowship and kinship of all the good arts, indeed of the very virtues, they do not know.

\ciceroSection{137} But to bring the speech back to the Greeks, whom in this kind of conversation we cannot do without (for as examples of virtue must be sought from ours, so of doctrine from them) — there are said to have been seven, at one time, who were both held and called wise. All these except Thales of Miletus presided over their states. Who is reported more learned in those same times than Pisistratus, or whose eloquence more equipped with letters? — who is said first to have so arranged Homer's books, before confused, as we now have them. He, indeed, was not useful to his fellow-citizens, but yet so flourished in eloquence that he excelled in letters and learning.

\ciceroSection{138} What of Pericles? About whose force of speaking we have heard this: that, when he spoke against the will of the Athenians somewhat sternly for the safety of his country, yet that very thing which he said against the popular men seemed popular to all and pleasant. On whose lips the old comic poets — even when they spoke ill of him (which was then permitted at Athens) — said that charm dwelt; and so great a force was in him that, in the minds of those who had heard him, he left as it were certain stings. But him no declaimer had taught to bark by the water-clock, but, as we have heard, that famous Anaxagoras of Clazomenae, the highest man in the science of the greatest matters. So this man, outstanding in learning, counsel, eloquence, presided over Athens for forty years, both in city and in war affairs.

\ciceroSection{139} What of Critias? what of Alcibiades? Not, indeed, good men for their states, but surely learned and eloquent — were they not trained in Socratic disputations? Who polished Dion of Syracuse with all doctrines? Was it not Plato? And the same man drove him on, equipped him, armed him, not only as master of his tongue, but of his soul and virtue, to free his country. With other arts, then, did Plato train this Dion, with others did Isocrates train that most distinguished man Timotheus, son of Conon the most outstanding general — himself the highest general and a most learned man? Or with others did that famous Pythagorean Lysis train the Theban Epaminondas, who I do not know whether he was the single highest man of all Greece? Or Xenophon Agesilaus? Or Philolaus Archytas of Tarentum? Or Pythagoras himself that whole old Greece of Italy, which once was called Magna?

\ciceroSection{140} I do not think so; for I see thus: that there was a single doctrine of all things which were worthy of a learned man, and of him who would excel in the commonwealth. Those who had received it, if they had been equally strong in talent for delivery and had given themselves also to speaking with no recoil of nature, excelled in eloquence.

\ciceroSection{141} So Aristotle himself, when he saw Isocrates flourishing through the noble birth of his pupils — because he had transferred his disputations from forensic and civic causes to an empty elegance of speech — suddenly changed nearly the whole form of his discipline, and quoted, with a slight change, that verse of Philoctetes. For he says it is base to him to be silent when barbarians; but Aristotle, when Isocrates speaks. So he ornated and lit up that whole doctrine, and joined the knowledge of matters with the practice of speech. Nor did this escape the wisest king Philip, who summoned this man as teacher for his son Alexander, that from him he might receive both the precepts of doing and of speaking.

\ciceroSection{142} Now whoever wishes — let him call by my leave a philosopher who hands down to us the abundance of matters and speech — by the name of orator. Or whether he prefer to call this orator (whom I say has wisdom joined to eloquence) a philosopher, I shall not hinder; provided this stand: that neither the speechlessness of him who knows the matter but cannot set it out by speaking, nor the ignorance of him to whom the matter is not at hand though words are not wanting, is to be praised. Of which, if one had to be wished, I should rather choose unspoken prudence than chattering folly. But if we ask what one excels of all, the palm must be given to the learned orator;

\ciceroSection{143} whom if they suffer to be the same as the philosopher, the controversy is removed. But if they part them, those will be the lower in this — that in the perfect orator there is all that science of theirs, while in the philosophers' knowledge eloquence does not necessarily lie; which, although they despise it, must yet seem to bring some heap to their arts.'' When Crassus had said this, he was silent for a little, and there was silence from the rest.

\ciceroSection{144} Then Cotta: ``For my part, Crassus,'' he said, ``I cannot complain that you seem to me to have argued something else, which you had not undertaken; for you have brought somewhat more than was assigned to you and laid down by us. But surely those parts of yours were that you should speak about illustrating the speech, and you yourself had now begun, and had described all the praise of speech in four parts. When you had spoken about the first two enough, indeed (but, as you yourself said, swiftly and briefly), you had made two for yourself remaining: how we might first speak ornately, then aptly.

\ciceroSection{145} When you had entered on this, suddenly a kind of tide of your talent snatched you far from the land and carried you out into the deep, almost out of the sight of all. For you, having grasped the whole knowledge of things, did not indeed hand it down to us; for time was not so much. But what you have profited with these, I do not know. Me indeed you have driven wholly into the Academy. In which I should wish that to be true which you often laid down: that it is not necessary to use up one's life, and that he can discern all those things who has only looked at them. But even if it is sometimes denser, or if I am slower, I shall surely never rest nor be wearied before I have grasped those two-headed paths and methods of arguing on either side and against everything.''

\ciceroSection{146} Then Caesar: ``One thing,'' he said, ``most has stirred me from your speech, Crassus: that you said that he who could not learn anything quickly could ever entirely master it. So that it is not difficult for me to make trial, and either at once to grasp those things which you have lifted with words to the heavens, or, if I cannot, not to lose time, since I can yet be content with these of ours.''

\ciceroSection{147} Here Sulpicius: ``For my part,'' said he, ``Crassus, I am wanting neither for that Aristotle nor Carneades nor any of the philosophers. Either think me to be despairing of mastering those things, or — what I do — to despise them. To me this common knowledge of forensic and common matters is great enough for that eloquence at which I aim. From which itself, however, I do not know very many things; which then at last, when some case to be pleaded by me requires it, I ask. Therefore, unless you happen to be wearied and if we are not heavy to you, come back to those things which pertain to the praise and splendour of speech itself; which I wished to hear from you, not so as to despair of attaining eloquence, but so as to learn something more.''

\ciceroSection{148} Then Crassus: ``You ask common matters,'' he said, ``and not unknown to you, Sulpicius. For who has not taught about that kind, has not trained, has not even left writings? But I shall humour you, and I shall set out for you briefly only those things which are known to me. I shall yet judge that we must come back to those who are the authors and inventors of these surely small matters.

\ciceroSection{149} All speech, then, is composed of words, of which we must first see the rationale singly, then jointly. For there is one ornament of speech which lies in single words; another which consists of words joined continuously. Therefore we use words either of those which are proper and as it were the fixed names of things, born almost together with the things themselves; or of those which are transferred and as it were placed in another's place; or of those which we make new and shape ourselves.

\ciceroSection{150} In proper words, then, the orator's praise is to flee what is cast off and obsolete and to use what is choice and bright, in which something full and resounding seems to be. But in this kind of proper words a certain choice must be had, and that must be weighed by some judgment of the ears; in which the habit of speaking well also avails most.

\ciceroSection{151} So that which the unskilled commonly say of orators — ``This man uses good words,'' or ``Some other does not'' — is not weighed by any art, but is judged by some natural sense, as it were. In which it is not great praise to avoid fault, although it is great. But this is, as it were, a single ground and foundation: the use and abundance of good words.

\ciceroSection{152} But what the orator himself builds, and where he applies art — that, it seems to me, must be sought and set out by us. There are three things, then, in single words which the orator should bring for illustrating and adorning speech: either an unaccustomed word, or a new-coined, or a transferred.

\ciceroSection{153} Unaccustomed are mostly old and through their antiquity long since dropped from the use of daily speech, which are freer for the licence of poets than ours; but rarely some poetic word has dignity also in speech. Nor would I shrink from saying, with Caelius, ``In what season the Punic came into Italy''; nor ``offspring'' or ``stock'' or ``to declare aloud'' or ``to call by name''; or, as you, Catulus, are wont, ``I did not deem'' or ``opined''; or many others, by which, when set in place, the speech is wont often to seem grander and more ancient.

\ciceroSection{154} Coined words are those which are produced and made by the speaker himself, either by joining words, as: ``\textit{Then a panic crushes from my breast all my wisdom},'' or ``Surely you do not wish his crooked-tongued malices'' — for both \textit{versutiloquas} and \textit{expectorat} are words made from joining, not born; but often even without joining words are coined, like ``that abandoned old age (\textit{senium})'', ``the gods of begetting (\textit{di genitales})'', ``to bend down (\textit{incurvescere}) with abundance of berries.''

\ciceroSection{155} The third mode of transferring a word is widely open. Necessity, forced by want and straits, gave it birth; afterwards pleasantness and delight made it celebrated. For as clothing, found first to ward off cold, was afterwards begun to be applied also to the body's adornment and dignity, so the transfer of a word was set up for the sake of poverty, made frequent for delight. For ``the vine sets gems,'' ``there is luxury in the grasses,'' ``glad crops,'' even peasants say. For what can scarcely be made plain by the proper word — when said by transfer, the likeness of the thing we have set in another's word illustrates that which we wish to be understood.

\ciceroSection{156} These transfers are as it were borrowings, when you take from another what you do not have. Those are a little bolder, which do not show poverty, but bring some splendour to speech. Of which I would not lay down for you either the reasoning of inventing or the kinds.

\ciceroSection{157} Of likeness it is the brevity drawn into a single word — which word, set in another's place as in its own, if it is recognised, delights; if it has nothing similar, is rejected. But those things ought to be transferred which either make the matter brighter, as all those: ``\textit{The sea bristles, the shadows are doubled, and the blackness of night and storm cloud darkens; flame quivers among the clouds, the heaven trembles with thunder, hail mixed with copious rain falls suddenly headlong, on every side all the winds break out, savage whirlwinds rise, the deep boils with seething}.'' All this, that it might be brighter, was said by words transferred through likeness;

\ciceroSection{158} or by which the whole thing of some deed or counsel is more signified, as he who, of one hiding it on purpose lest what was being done could be understood, by two transferred words, by likeness itself, signifies: ``\textit{Since he wraps himself round in words, fences himself with guile.}'' Sometimes brevity also is got by transfer, as that ``If the weapon flew from his hand'': the imprudence of a weapon let go could not be set out more briefly by proper words than has been signified by one transferred.

\ciceroSection{159} In this kind it often seems to me wonderful what the cause is, that all are more pleased by transferred and foreign words than by proper and their own. For if a thing has not its name and proper word — as the foot in a ship, as \textit{nexum} which is done by the scale, as in a wife divorce — necessity forces us to take from elsewhere what we do not have. But in the greatest abundance of one's own words men are yet much more delighted by foreign ones, if they are transferred with reason.

\ciceroSection{160} I think this happens, either because it is some sign of talent to leap over what lies before our feet and take other things sought from afar; or because he who hears is led by his thought elsewhere and yet does not stray, which is the greatest delight; or because in single words a thing and a whole likeness is brought about; or because every transfer, which has been taken with reason, is brought to the senses themselves, especially of the eyes — which is the keenest sense.

\ciceroSection{161} For both ``the smell of urbanity'' and ``the softness of humanity'' and ``the murmur of the sea'' and ``the sweetness of speech'' are drawn from the other senses; but those of the eyes are far keener, which place almost before the gaze of the soul what we cannot perceive and see. For there is nothing in the nature of things whose word and name we cannot use of other things. Whence a likeness can be drawn — and it can be drawn from all things — thence one word, transferred, containing the likeness, will bring light to speech.

\ciceroSection{162} In this kind unlikeness must first be fled. ``Vast vaults of heaven'': however much Ennius brought a sphere onto the stage, as is said, yet in a sphere the likeness of a vault cannot be. ``\textit{Live, Ulysses; while you may; with your eyes the last beam of light snatch up}'' — he did not say ``seek'' or ``take'' (for that would have the delay of one hoping to live longer), but ``snatch'': this word is fitted to that

\ciceroSection{163} which he had said before, ``while you may.'' Then it must be looked at lest the likeness be drawn from afar. ``Syrtis'' of patrimony — I should rather say ``rock''; ``Charybdis'' of goods — rather ``whirlpool.'' For the eyes of the mind are more easily borne to those things which are seen than to those which are heard. And since this is the very highest praise in transferring words — that what is transferred should strike the sense — all baseness must be fled in those things to which the likeness will draw the minds of those who hear.

\ciceroSection{164} I would not have said that by Africanus's death the commonwealth was ``castrated''; I would not have Glaucia called ``the dung of the curia.'' Even though the likeness be there, in each is a deformed thinking of likeness. I would not have it either greater than the matter requires (``a tempest of revelry''), or less (``a revelry of a tempest''). I would not have the word that is transferred narrower than that proper one and its own would have been: ``What is it, pray? What do you nod me away from approach?'' It would be better, ``you forbid, prohibit, frighten away'' — since he had said: ``\textit{Get back here, lest my touch or shadow harm the good.}''

\ciceroSection{165} If you fear that the transfer may seem somewhat too harsh, it must often be softened by a word set before. As if, when M. Cato had died, anyone should say the senate had been left ``orphaned,'' that would be a little harsh; but ``so to say, orphaned'' is somewhat milder. For a transfer ought to be modest, so that it may seem brought to a foreign place, not to have broken in; and to have come on sufferance, not by force.

\ciceroSection{166} But there is no more flowery way in single words, nor any which brings more light to speech. For that which flows from this kind is not in one word transferred, but is connected from several joined, so that one thing is said, another to be understood: ``\textit{Nor shall I let myself again on the one rock, as once the Achaean fleet, dash myself.}'' And: ``\textit{You err, you err. For when you leap and trust yourself, the strong reins of the laws will rein you in, and the yoke of empire will press you down.}''

\ciceroSection{167} A like matter being taken, the words proper to it are then transferred to another matter, as I said. This is a great ornament of speech, in which obscurity must be fled. For these things, when carried too far in this kind, become what are called \textit{enigmas}. The way is not in a word but in a speech, that is, in the continuity of words. Nor again does that turning and changing have a kind of fashioning in a word but in speech: ``\textit{Africa trembles with terrible tumult, the rough land}.'' For the Africans is taken Africa: nor was the word made (as ``the rock-breaking sea''), nor transferred (as ``the sea is softened''), but for ornament's sake one proper word is changed for another. ``\textit{Cease, Rome, your enemies — and the great plains are witnesses.}'' This way is grave in adornment of speech and often to be taken; of this kind: that Mars of war is common, to call Ceres for crops, Liber for wine, Neptune for the sea, the curia for the senate, the Campus for the assemblies, the toga for peace, arms and weapons for war.

\ciceroSection{168} Likewise in the same kind, both virtues and vices are named for those very persons in whom they reside: ``Luxury, when it broke into the house''; ``Avarice penetrated where''; or ``Faith availed, justice carried it through.'' You see this whole kind: when, by an inflected and changed word, the same matter is uttered more ornately. To which are neighbouring those less ornate but yet not to be unknown — when we wish something to be understood: either the whole from a part (as for buildings when we say walls or roofs); or a part from the whole (as when we say one squadron for the cavalry of the Roman people); or several from one: ``\textit{But the Roman, even though the matter is well done, in his heart trembles}''; or when from many one is understood: ``\textit{We are the Romans, who before were Rudini}'' (Rudine, Ennius's home). Or in whatever way — not as it is said, in this kind is it understood, but as it is felt.

\ciceroSection{169} We often also misuse a word, not so elegantly as in transferring, but, even if more freely, sometimes not impudently — as when we say ``grand speech'' for ``long,'' ``minute soul'' for ``small.'' Now do you see those things to be not of a word but of a speech, which are connected from several transfers, as I have set out? But these — which I have said are either changed or to be otherwise understood than they are said — are transferred in a certain way.

\ciceroSection{170} So it happens that all the virtue and praise of single words arises from three things: if there be either an old word — which yet usage can bear; or a coined one, either by joining or by novelty, in which likewise the ears and habit must be spared; or a transferred one, which most, as it were by certain stars, marks and lights up speech.

\ciceroSection{171} The continuity of words follows, which especially desires two things: arrangement first, then a certain measure and form. Of arrangement, it is to compose and build words so that neither be their meeting harsh nor gaping, but in some way close-jointed and smooth. In which my father-in-law's character was wittily played upon by him who could most elegantly do it, Lucilius: ``\textit{How prettily are the words composed! As all the little dice fitted skillfully into the pavement and inlay-work, vermicular!}'' When he had said this, mocking Albucius, he did not refrain from me either: ``\textit{I have Crassus as son-in-law, that you not be more rhetorical.}'' What then? This Crassus, since he abuses my name — what does he do? That, surely (as he wishes and I should wish), somewhat better than Albucius. But on me he played, as he is wont.

\ciceroSection{172} Yet this arrangement of words must be kept, of which I am speaking; which makes the speech bound, coherent, smooth, evenly flowing. You will attain this if the last words shall so be joined with the first that follow that neither do they harshly meet nor too widely part.

\ciceroSection{173} This diligence is followed also by the measure and the form of words; which I now fear may seem to this Catulus to be childish. For those ancients thought verses ought to be applied even in this prose speech, that is, certain rhythms — for they wished there to be in speeches closes punctuated by the measure of words and thoughts (not by marks of our breathlessness or of weariness, nor by copyists' marks, but by the measure of words and thoughts); and Isocrates is said first to have established this — that he might bind the unkempt habit of the ancients in speaking, for the sake of pleasure and the ears, in rhythms (as his pupil Naucrates writes).

\ciceroSection{174} For these two things were contrived for pleasure by the musicians, who once were the same as the poets — verse and song — that by both the rhythm of words and the measure of voices they might overcome the satiety of the ears with delight. These two then — I mean the regulation of voice and the closing of words — they thought were to be transferred from poetry to eloquence, as far as the severity of speech could bear.

\ciceroSection{175} In which is most of all this: that if a verse is brought about in speech by the joining of words, it is a fault — and yet we wish that joining, like the verse, to fall in rhythm and to be squared and brought to fulness. Nor of many things is there one matter which more distinguishes the orator from the unskilled and ignorant of speaking, than that he, untaught, pours out unkemptly as much as he can and what he says by breath, not by art, he limits; while the orator so binds the thought with words that he embraces it by some rhythm both bound and free.

\ciceroSection{176} For when he has bound it by form and measures, he relaxes and frees it by changing the order, that the words be neither bound as it were by some fixed law of verse, nor so loose as to wander. In what way then shall we set ourselves to so great a duty, as to suppose we can attain this force of speaking rhythmically? The thing is not so difficult as necessary. For nothing is so tender, nor so flexible, nor that follows so easily wherever you lead, as speech.

\ciceroSection{177} From this verses, from this same different rhythms are produced; from this also this prose speech of various measures and many kinds. For there are not different words of conversation, of contention, nor are they taken from another kind for daily use, from another for the stage and pomp. But these we, when we have lifted them up lying in the middle, like the softest wax shape and fashion at our pleasure. Therefore as at one time we are weighty, at another subtle, at another we hold something middle: so the kind of speech follows our settled feeling, and is changed and turned to every pleasure of the ears and motion of minds.

\ciceroSection{178} But as in most matters nature itself has incredibly fashioned this, so in speech: that those things which contained in themselves the greatest utility, the same had also most either of dignity or even of grace. For the safety and the salvation of all, we see this state of this whole world and of nature: that the heaven is round, and the earth is in the middle and is held by its own force and will; that the sun is borne about it, that it approaches the wintry sign and thence gradually rises into the opposite quarter; that the moon by its approach and recess receives the sun's light; that the same five stars complete the same spaces by unequal motion and course.

\ciceroSection{179} These have such force that, a little changed, they cannot cohere; such beauty that no shape can even be thought more ornate. Refer your mind now to the form and shape of men, or even of the rest of living things. You will find no part of the body, without some necessity, added on, and the whole shape as it were perfected by art, not by chance. What in those trees? In which the trunk, the branches, the leaves are not, in short, save for keeping and preserving their nature; nowhere yet is any part except graceful.

\ciceroSection{180} Let us leave nature and look at the arts. What in a ship is so necessary as the sides, as the holds, as the prow, as the stern, as the spars, as the sails, as the masts? Which yet have in their look so much grace, that they seem to have been invented not only for safety, but also for pleasure. Columns hold up temples and porticoes; yet they have no more usefulness than dignity. The pediment of the Capitol and of other temples necessity itself, not grace, fashioned. For when account had been had how on each side of the roof water might run off, the dignity of the pediment followed the temple's utility — so that, even if the Capitol were set up in heaven, where there could be no rain, it seems it would have had no dignity without a pediment.

\ciceroSection{181} The same likewise happens in all parts of speech: that some sweetness and charm follows the utility and almost necessity. For the closes and the punctuations of words were brought by the closing of breath and the straits of the breath. That invention is so sweet that, even if anyone had infinite breath given him, we should not yet wish him to draw out his words. So pleasing has been found to our ears that which not only could be tolerable but also easy to men's lungs.

\ciceroSection{182} The longest combination of words, then, is what can be rolled out in one breath. But this is the measure of nature; another, of art. For since rhythms are several, the iamb and the trochee, the frequent ones, Aristotle, your Aristotle, Catulus, sets aside from the orator. They yet by nature run themselves into our speech and conversation. But marked are the strikings of those rhythms and minute feet. So first of all the heroic foot — the dactyl and anapest, the spondee — invites us. In which it is allowed to advance with impunity for two feet only, or a little more, lest we plainly fall into a verse or the likeness of a verse. ``\textit{High are the twin towers, by which}.'' These three heroic feet at the openings of continuing words fall sufficiently becomingly.

\ciceroSection{183} Most approved by that same Aristotle is the paean, which is double. For it either rises from a long syllable, which three short follow, as these words ``cease, begin, restrain''; or from three short syllables in succession, the last drawn out and long, as those are ``had subdued, hoof-beating.'' And that philosopher likes to begin from the former paean, to end with the latter; the latter paean is, not in the number of syllables but in the measure of the ears (which is keener and more certain judgment), about equal to the cretic, which is from a long, a short, and a long: as ``\textit{What guard shall I seek, or pursue? whither now —}'' From which rhythm Fannius began: ``\textit{If, Quirites, his threats}.'' This rhythm Aristotle thinks more apt for closes, which he wishes mostly to be ended by a long syllable.

\ciceroSection{184} These things, however, do not require so keen a care and diligence as does that of the poets — whom necessity compels, and rhythms and measures themselves, so to enclose words in a verse that nothing may be even by the smallest breath shorter or longer than is needed. Speech is more free, and plainly, as it is said, is so truly loose — not so as to flee or err, but so as to regulate itself without bonds. For I agree with Theophrastus, who thinks that speech, polished and made (in some way), ought to be rhythmical not strictly but more loosely.

\ciceroSection{185} For, as he conjectures, both from those measures by which this customary verse is brought about, the anapaest afterwards flowered out (a longer rhythm); thence flowed that more licentious and richer dithyramb, whose limbs and feet, as the same says, lie spread through every richer speech. And if it is rhythmical in all sounds and voices that has certain pressings and what we can measure by equal intervals, this kind of rhythms is rightly set in the praise of speech, provided they be not continuous. For if that breathless and flowing chatter without intervals is to be deemed rude and unpolished, what other reason is there why it should be rejected, save that nature itself modulates the voice for men's ears? Which cannot happen unless rhythm is in the voice.

\ciceroSection{186} But there is no rhythm in continuity. Distinction and the striking of equal — or often various — intervals brings about rhythm, which we can mark in falling drops (which are distinguished by intervals), but cannot in a rushing river. If then this loose continuity of words is much more apt and pleasant when it is distinguished by joints and limbs than when continued and prolonged, those limbs will have to be measured. If they are shorter at the end, that as it were ambit of words is broken; for so the Greeks call these turnings of speech. So either the latter must be equal to the former, the last to the first; or, what is even better and more pleasant, longer.

\ciceroSection{187} These things have indeed been said by those philosophers whom you most love, Catulus. So I more often call them as witnesses, that, by praising my authorities, I may flee the charge of frivolities.'' ``Of what frivolities?'' said Catulus. ``Or what can be brought into that disputation more elegant, or more subtly said at all?''

\ciceroSection{188} ``But I fear,'' said Crassus, ``that these may either seem too hard for these to follow, or, because they are not handed down in the common discipline, that we may seem to wish them to seem greater and more difficult.'' Then Catulus: ``You err, Crassus,'' he said, ``if you think either I or any of these expects from you these daily and common things. Those things you say, we wish to be said; nor so much to be said as to be said in that way. Nor do I answer this for myself only, but for all of these without any doubt.''

\ciceroSection{189} ``For my part,'' said Antonius, ``I have now found him whom I had said in the little book I wrote I had not found — the eloquent man. But I have not interrupted you even for the sake of praising, lest anything of this so brief time of your speech should be lessened by one of my words.''

\ciceroSection{190} ``By this rule, then,'' said Crassus, ``both by exercise and by the pen — which both other things and this most adorns and polishes — speech must be shaped by us. Yet this is not of so great labour as it seems, nor are these things to be directed by the keenest standard of rhythmicians or of musicians. Only this we must bring about: that speech does not flow on, does not wander, does not stop within, does not run too far out — that it be distinguished by limbs, that it have absolute turnings. Nor must we always use the unbroken and as it were turning of words, but often speech must be broken into smaller limbs, which yet themselves must be bound by rhythms.

\ciceroSection{191} Nor let the paean or the heroic disturb you. They will themselves run into speech; they themselves, I say, will offer themselves and answer when not called. Only let there be that habit of writing and speaking, that thoughts may be ended by words, and the joining of those words may be born from longer and freer rhythms — most of all the heroic, or the former paean, or the cretic, but variously and distinctly let it settle. For likeness is most marked in coming to rest. And if the first and last feet have been preserved by this rationale, the middle ones can lie hidden, provided the very circuit of words be not either shorter than the ears expect, or longer than strength and breath bears.

\ciceroSection{192} The closes I think must be even more diligently kept than the previous parts, because in them most of all is the perfection and absoluteness judged. For of a verse the first and middle and last part is equally attended to, which, in whatever part it stumbles, is weakened. In speech, however, few perceive the first parts, most the last. Which since they appear and are understood, must be varied, lest they be rejected either by the judgments of minds or by the satiety of ears.

\ciceroSection{193} Two or three feet at the end must be kept and noted, provided the previous parts will not be shorter and clipped. These ought to be either chorees or heroic, or alternating; or in that latter paean which Aristotle approves, or in the cretic equal to it. The changes of these will bring it about that neither shall those who hear be sated by disgust at likeness, nor we seem to do, by set design, what we shall do.

\ciceroSection{194} If that Antipater of Sidon — whom you well remember, Catulus — was wont to pour out hexameters and other verses on the spur of the moment in various measures and rhythms, and the practice of an ingenious and remembering man so prevailed that, when he had thrown himself in mind and will to the verse, the words followed: how much more easily shall we attain that in speech, with practice and habit applied!

\ciceroSection{195} But that no one may wonder how the multitude of unskilled in listening notices these things — both in every kind, and in this very thing, there is some great force and incredible nature. For all by some silent sense, without any art or reason, distinguish what is right and faulty in arts and reasonings. And while they do this in paintings and statues and other works, for the understanding of which they have less equipment from nature, much more do they show it in the judgment of words, rhythms, and voices, since these are fixed in common sense, and nature has wished none to be wholly without them.

\ciceroSection{196} So all are moved not only by skilfully placed words, but also by rhythms and voices. For how few there are who hold the art of rhythms and measures? But in these, if there is offence even a little — that anything is made shorter by contraction or longer by drawing out — whole theatres cry out. Does not the same happen in voices, that not only choirs and concerts, but also each one to himself, even singers in mismatch are cast out by the multitude and the people?

\ciceroSection{197} It is wonderful — when there is the greatest difference in doing between the learned and the rude — how little the difference is in judging. For art, since it has set out from nature, unless nature moves and delights, seems to have done nothing at all. There is nothing so akin to our minds as rhythms and voices, by which we are roused and inflamed and softened and made languid and led often to cheerfulness and to gloom. The highest force of which is more apt for songs and singing — not neglected, as it seems to me, by Numa the most learned king and by our ancestors, as the lyres and pipes and verses of the Salii at solemn banquets show; most of all celebrated by ancient Greece. With which and like things would that you had preferred to argue rather than these childish transfers of words!

\ciceroSection{198} But as in a verse the multitude sees if there is a fault, so in our speech, if anything halts, it feels — but it does not pardon the poet, while it grants it to us. Yet all silent see that what we have said is not apt and perfect. Therefore those ancients — as we see today some too — when they could not bring about a circuit and as it were an orb of words (for that we have lately even begun to be able or to dare), used to speak three or two, or some even single words; who in that infancy yet held that natural thing which men's ears demanded — that what they said should be equal, and they should use equal pauses for breathing.

\ciceroSection{199} I have set out, as far as I could, what I judged most pertained to the ornament of speech. For I have spoken of the praise of single words; I have spoken of their joining; I have spoken of rhythm and form. But if you also ask the habit of speech and as it were some colour, there is one full but yet smooth, and one slender, not without sinews and strength, and one which, partaking of each kind, is praised in a certain mediocrity. On these three figures should sit a colour of grace not painted with rouge, but spread by blood.

\ciceroSection{200} Then at last this orator must be so shaped by us, both in words and in thoughts, that, as those who use arms or the wrestling-school, he should think he has to have account not only of avoiding or striking, but also that he move with grace; as those who are versed in the handling of arms, so should he in words for the apt composition and becomingness, but in thoughts for the gravity of speech. Words and thoughts are shaped almost countlessly, which I know enough is known to you. But between the shaping of words and of thoughts there is this difference: that of words is removed if you change the words; that of thoughts remains, whatever words you wish to use.

\ciceroSection{201} Although you do this, I yet think you ought to be warned that you should not think there is anything else to be the orator's, that is outstanding and wonderful, except, in single words, to hold those three things — that we should use transferred words frequently, sometimes coined, rarely also very ancient ones. In the continuous speech, when we have held both the smoothness of joining and the rationale of rhythms which I have spoken of, then must all speech, as it were with lights, be distinguished and frequented with thoughts and words.

\ciceroSection{202} For dwelling on one subject moves very much; and the bright unfolding and the placing of matters under the eye, almost as if they were being done; which both in setting forth a matter avail most, and for illustrating what is set forth, and for amplifying — so that to those who hear, what we shall amplify shall seem to be as great as speech can effect. To this is contrary often a running over; and signification of more for understanding than you have said; and the distinctly clipped brevity, and lessening, and joined to it the irony not strange to Caesar's precepts;

\ciceroSection{203} and digression from the subject, in which when there has been pleasure, then the return to the subject ought to be apt and neat. And the laying down of what you are to say; and a setting-apart from what has been said; and a return to the proposition, and repetition, and the apt closing of the reasoning. Then for the sake of amplifying or diminishing, the over-statement and going-beyond truth; and questioning, and a kind of inquiry akin to it, and the setting out of one's own thought; then that which most as it were creeps into men's minds: dissimulation that says other things and means other; which is most pleasing when it is handled not with the contention of speech but with conversation. Then doubt; then distribution; then correction either before or after you have spoken or when you reject something from yourself. There is also pre-fortification for what you are to attack, and casting onto another;

\ciceroSection{204} consultation, which is as it were deliberating with those very ones before whom you speak; the imitation of morals and life either in persons or without them — a great ornament of speech, apt for conciliating minds especially, but often also for moving;

\ciceroSection{205} the bringing in of feigned persons, the weightiest light of amplifying; description; the leading into error; an impulse to mirth; pre-occupation; then those two which most move: likeness and example; division; interruption; contention; reticence; commendation; some free voice and even unreining for the sake of amplifying; anger; rebuke; promise; entreaty; supplication; a brief turning aside from the proposition (not as the earlier digression); purgation; conciliation; injury; wishing and execration. With these lights mostly thoughts illustrate speech.

\ciceroSection{206} Of speech itself, as of arms, there is either threat for use and as it were attack, or for grace itself. For doubling of words has sometimes force, sometimes charm; and a word a little changed and bent away; and the frequent repetition of the same word, now from the beginning, now turned to the end; and the rush upon the same words; and concourse and joining and progression and a kind of distinguishing of the same word more often placed; and the recalling of a word; and those that end alike or fall alike or are matched with equal returns or are like one another.

\ciceroSection{207} There is also a kind of climax and turning-back and neat transposition of words; and the contrary; and the loose; and turning aside; and reproof; and exclamation; and lessening; and what is set in many cases; and what, drawn from single matters proposed, is referred to single ones; and the reasoning subjoined to the proposition; and likewise in the distributed the reasoning placed under; and permission; and again another doubt; and something unforeseen; and enumeration; and another correction; and dispersion; and the continuous and the broken; and image; and answer to oneself; and changing; and disjunction; and order; and reference; and digression; and circumscription.

\ciceroSection{208} These mostly, and like to these, or even more, are those things which by thoughts and shapings of words illustrate speech.'' ``Which I see you, Crassus,'' said Cotta, ``— because you think they are known to us — have poured out without definitions and without examples.'' ``For my part,'' said Crassus, ``I did not think those things either, which I said before, were new to you, but I obeyed all your wish.

\ciceroSection{209} About these matters that sun has admonished me to be briefer, which itself now sliding has compelled me too to roll out these things almost slipping. But yet the demonstration of this kind, and the doctrine itself, is common; the use, however, is the weightiest, and in this whole study of speaking the most difficult.

\ciceroSection{210} Therefore, since about all the ornament of speech all topics have been, if not laid open, at least pointed out, now let us see what is apt — that is, what is most becoming in speech. Although that is plain — that not to every cause, nor every hearer, nor every person, nor every time does one kind of speech agree. For both capital cases require some kind of sound of words, another of private and small ones;

\ciceroSection{211} and another kind of speaking deliberations, another laudations, another judgments, another conversations, another consolation, another rebuke, another disputation, another history requires. It also matters who hear: senate or people or judges; many or few or single, and what kind. The orators themselves, of what age, honour, authority, must seem; the time, of peace or war, of haste or leisure.

\ciceroSection{212} So in this place there is nothing that can seem to be precepted, save that we choose the figure of fuller and slenderer speech, and likewise of that middle, suited to what we shall be doing. The same ornaments may be used, sometimes more contendedly, sometimes more loosely. To be able in every matter to do what becomes is of art and nature; to know what becomes when, is of prudence.

\ciceroSection{213} But all these things are as they are delivered. Delivery, I say, in speaking, alone is queen. Without this the highest orator can be in no number, and one of middling ability, equipped with this, often surpasses the highest. Demosthenes is said to have given to delivery the first place, when he was asked what was first in speaking; to it the second; to it the third. From which that of Aeschines also seems to me even better said: who, when on account of the disgrace of his trial he had withdrawn from Athens and had betaken himself to Rhodes, asked by the Rhodians, is reported to have read that outstanding speech which he had spoken in Ctesiphon against Demosthenes. When he had read it through, he was asked the next day to read also that which had been said on the other side by Demosthenes for Ctesiphon. When he had read it with most sweet and great voice, all wondering: ``How much,'' he said, ``would you have wondered, if you had heard the man himself!'' From which he sufficiently signified how much was in delivery, since he thought the same speech would be another with the speaker changed.

\ciceroSection{214} What was there in Gracchus, whom you better remember, Catulus, that, when I was a boy, was so greatly noted? ``\textit{Whither, wretched, shall I betake me? Whither shall I turn? To the Capitol? But it is wet with my brother's blood. Or home? That I see my mother lamenting, wretched, and cast down?}'' Which were so delivered by him, by his eyes, voice, gesture, that even his enemies could not hold back their tears. These things I say at greater length, because this whole kind orators (who are actors of truth itself) have left, while imitators of truth — actors — have seized it.

\ciceroSection{215} Without doubt in every matter truth conquers imitation; but if it could effect enough in delivery by itself, we should not have need of art at all. But because the movement of the soul, which most must be either declared or imitated by delivery, is often so disturbed that it is obscured and almost overwhelmed, those things must be shaken off which obscure, and those which are eminent and ready must be taken.

\ciceroSection{216} For every motion of the soul has from nature its own face and sound and gesture; and the whole body of a man, and all his face, and all his voices, like strings on a lyre, so sound, as they are struck by the motion of the soul. For voices, like chords, are stretched, which respond to each touch — sharp, low; quick, slow; great, small. Among all of which yet there is in each kind a moderate one; and from these have flowed several other kinds: smooth, harsh; contracted, diffuse; with continuous, with broken breath; broken, split; thinned, swelled with bent sound;

\ciceroSection{217} for there is none of these kinds which is not handled by art and moderation. These are set out for the actor, as for the painter, for varying his colours. For one kind of voice anger takes for itself: sharp, urged on, frequently breaking off: ``\textit{My brother himself urges me, wretched, to deliver my children to the jaws}'' — and what you, Antonius, brought forward a little while ago: ``\textit{Did you dare to set yourself apart}'' and ``\textit{Will any one notice this? Bind, beat}'' — and almost the whole of \textit{Atreus}. Another, pity and grief: flexible, full, broken, with a tearful voice: ``\textit{Whither now shall I turn me? What way shall I begin to enter? My father's house? Or to the daughters of Pelias?}'' And: ``\textit{O father, O country, O house of Priam!}'' — and what follows: ``\textit{All this I saw burning, and Priam's life torn from him by force.}''

\ciceroSection{218} Another, fear: lowered and hesitating and cast down: ``\textit{In many ways have I been beset — by sickness, exile, and want; then panic has crushed all my wisdom from my breath; my mother threatens dire torment of life and death — what no one is of so firm a mind and so great confidence that he will not feel his blood shrink with fear and grow pale.}''

\ciceroSection{219} Another, force: stretched, vehement, threatening with a kind of impulse of weight: ``\textit{Again Thyestes comes attended by Atreus; again now he comes upon me, rouses me from rest. A greater mass for me, a greater evil to be mixed, that I may dash and crush his bitter heart.}'' Another, pleasure: poured out, gentle, tender, cheerful, slack: ``\textit{But when she had brought to herself the wreath for binding the marriage, while she pretended to give it eagerly to herself, then to you, gaily, learnedly and daintily, she brought it.}'' Another, distress: without pity, somewhat heavy and brought through with one pressing and tone: ``\textit{In the season when Paris joined Helen in unwedded marriage, I was then heavy with child, having now fulfilled my months for bearing; at the same time Hecuba bore Polydorus, her last child.}''

\ciceroSection{220} All these motions gesture must follow — not that of the stage, expressing words, but declaring the whole matter and thought, not by demonstration but by signification, with the inflection of the sides, this strong and manly, not from the stage and from actors but from arms or even from the wrestling-school. The hand less talkative, with the fingers following the words, not expressing them; the arm thrown out longer, like a kind of weapon of speech; the stamp of the foot at the beginning or end of contention.

\ciceroSection{221} But all is in the face; in the face itself the eyes hold the lordship over all. Wherefore those old men of ours, who used not to praise even Roscius greatly when he wore a mask. For all action is of the soul, and the face is the image of the soul; the eyes its indicators. For this is the one part of the body which can produce as many significations and changes as there are motions of the soul. Nor is there anyone who could produce the same with eyes shut. Theophrastus indeed says one Tauriscus was wont to call an actor turned away who, in delivering, was looking at something.

\ciceroSection{222} So great moderation is of the eyes. For the look of the face must not be too much changed, lest we be carried either to ineptitudes or to some crookedness; the eyes are those by whose intentness, slackness, glance, mirth, we should signify the motions of soul aptly with the very kind of speech. For delivery is, as it were, the speech of the body, the more it must be agreeable to the mind. The eyes, however, nature gave us — as to a horse or lion bristles, tail, ears — for declaring the motions of the soul.

\ciceroSection{223} Therefore in our delivery, after the voice, the face avails most; and that is governed by the eyes. And in all those things which belong to delivery, there is some force given by nature; for which reason it is by this also that the unskilled, the multitude, indeed the foreigners are most moved. For words move no one save him who is joined by the fellowship of the same tongue. Often keen thoughts also fly past the senses of men not keen. Delivery, which carries on its face the motion of soul, moves all. For all are stirred by the same motions of soul, and recognise them by the same marks in others, and indicate them in themselves.

\ciceroSection{224} For the use and praise of delivery the voice without doubt holds the greatest part; which we must first wish for; then, whatever it shall be, we must guard it. About which it is no part of this kind of precepts how the voice should be served — although I much think it should be served. But this seems not foreign to the duty of this our conversation: that, as I said a little before, in most things what is most useful is also, somehow, most becoming. For nothing is more useful for keeping the voice than frequent change; nothing more harmful than uninterrupted, poured-out contention.

\ciceroSection{225} What is more apt for our ears and the sweetness of delivery than alternation, variety, change? So that same Gracchus — what you can hear, Catulus, from Licinius your client, a lettered man, whom Gracchus had as a slave at hand — was wont to have, with a little ivory pipe, a man standing secretly behind him when he spoke in assembly, an expert who would quickly blow that note by which he could either rouse him when slack or call him back from contention.'' ``By Hercules, I have heard,'' said Catulus, ``and I have often admired that man's diligence as also learning and science.''

\ciceroSection{226} ``For my part,'' said Crassus, ``I am also grieved that those men slipped into that ruin in the commonwealth. But that net is being woven, and that rationale of living in the state is shown to posterity, that we now desire to have like to those citizens whom our fathers did not bear.'' ``Let go, I beg, that conversation, Crassus,'' said Iulius, ``and bring yourself back to Gracchus's pipe; whose rationale I do not yet plainly understand.''

\ciceroSection{227} ``In every voice,'' said Crassus, ``there is something middle, but each voice has its own. From this it is useful and pleasant to ascend the voice gradually (for to shout from the start is something rustic), and the same is salutary for confirming the voice. Then there is some extreme of contention, which yet is more inward than the sharpest cry, beyond which the pipe will not let you advance, and now from the very contention will call you back. There is likewise on the contrary something heaviest in slackness, by which one descends as it were by steps of sounds. This variety, and this course through all sounds of voice, will both keep itself, and bring sweetness to delivery. But you will leave the piper at home; the sense of this habit you will bring with you to the Forum.

\ciceroSection{228} I have brought out what I could, not as I wished but as the straits of time have compelled me; for it is well noted to charge the matter on time, when, if you wished, you could not bring more.'' ``You yourself,'' said Catulus, ``have gathered everything, as far as I can judge, so divinely, that you seem not to have taken from the Greeks but to be able to teach those very ones these things. I am glad to have been made partaker of this conversation; and I should have wished my son-in-law, your friend, Hortensius, had been present, whom indeed I trust to be outstanding in all the praises which you have embraced in your speech.''

\ciceroSection{229} And Crassus: ``You say `to be'?'' he said. ``I now judge him so to be, and so judged when, in my consulship, he defended the cause of Africa in the senate, and lately, even more, when he spoke for the king of Bithynia. So you see rightly, Catulus. For I see that to that young man neither nature nor learning is wanting.

\ciceroSection{230} The more, Cotta, you must keep watch and labour, and you, Sulpicius. For he is no middling orator who, as it were, springs up into your generation, but with most keen talent and burning zeal and outstanding learning and singular memory. Whom, although I favour, I yet wish him to excel his own age; that he should outrun you who are so much younger is scarcely honourable.'' ``But now let us rise,'' he said, ``and let us care for ourselves and at last loosen our minds and care from this contention of disputation.''

